\documentclass[11pt,a4paper,openany]{book}
\usepackage[T1]{fontenc}
\usepackage[utf8]{inputenc}
\usepackage{lmodern}
\usepackage[a4paper,margin=27mm,headheight=15pt]{geometry}
\usepackage{amsmath,amssymb,amsthm,bm}
\usepackage{booktabs,longtable,tabularx,array}
\usepackage{graphicx,xcolor,microtype}
\usepackage{enumitem,listings}
\usepackage{xurl}
\usepackage[round,authoryear]{natbib}
\usepackage{tikz}
\usetikzlibrary{arrows.meta,positioning,fit}
\usepackage{fancyhdr}
\usepackage[colorlinks=true,linkcolor=blue!45!black,citecolor=teal!60!black,urlcolor=blue!55!black,breaklinks=true]{hyperref}
\usepackage{bookmark}
\definecolor{ink}{HTML}{17324D}
\definecolor{teal}{HTML}{23778A}
\definecolor{light}{HTML}{F1F5F8}
\pagestyle{fancy}\fancyhf{}
\fancyhead[LE,RO]{\small\thepage}
\fancyhead[RE]{\small Causal analysis for a universal bank}
\fancyhead[LO]{\small\nouppercase{\rightmark}}
\renewcommand{\headrulewidth}{0.3pt}
\setlength{\parskip}{4pt}
\setlength{\parindent}{12pt}
\setlength{\emergencystretch}{2em}
\setlength{\tabcolsep}{4pt}
\newcolumntype{P}[1]{>{\raggedright\arraybackslash}p{#1}}
\renewcommand{\tabularxcolumn}[1]{>{\raggedright\arraybackslash}p{#1}}
\setcounter{tocdepth}{1}
\let\cleardoublepage\clearpage
\renewcommand{\bibfont}{\small}
\setcounter{secnumdepth}{2}
\setlist{nosep,leftmargin=*}
\newcommand{\E}{\mathbb E}
\newcommand{\Prob}{\mathbb P}
\newcommand{\ind}{\mathbf 1}
\newcommand{\Var}{\operatorname{Var}}
\newcommand{\Cov}{\operatorname{Cov}}
\newcommand{\doop}{\operatorname{do}}
\newcommand{\ESS}{\operatorname{ESS}}
\newcommand{\ATE}{\operatorname{ATE}}
\newcommand{\ATT}{\operatorname{ATT}}
\newcommand{\given}{\,\vert\,}
\newtheorem{proposition}{Proposition}[chapter]
\newtheorem{assumption}{Assumption}[chapter]
\theoremstyle{definition}\newtheorem{example}{Example}[chapter]
\newenvironment{decision}{\begin{quote}\color{ink}\bfseries}{\end{quote}}
\lstset{basicstyle=\ttfamily\small,breaklines=true,columns=fullflexible,frame=single,
 backgroundcolor=\color{light},showstringspaces=false}
\hypersetup{pdftitle={Causal Analysis for a Universal Bank},pdfauthor={Research and product proposal}}

\newcommand{\NumStaticATE}{8.472}
\newcommand{\NumStaticSE}{0.254}
\newcommand{\NumStaticNaive}{55.586}
\newcommand{\NumLongNaive}{19.137}
\newcommand{\NumTWFE}{1.200}
\newcommand{\NumPanelATT}{3.500}
\newcommand{\NumFilterLikelihood}{0.023}
\newcommand{\NumHMMError}{0.050}
\newcommand{\NumBreakEven}{3.408}
\newcommand{\NumBaseNPV}{3.958}
\newcommand{\NumLowNPV}{-3.503}
\newcommand{\NumHighNPV}{13.905}

\newcommand{\SoftwareCount}{44}
\newcommand{\PublicRequestCount}{415}
\newcommand{\CandidateCount}{457}
\newcommand{\SourceCount}{187}

\title{\sffamily\bfseries\color{ink}\Huge Causal Analysis for a\\Universal Bank\\[16pt]
\Large A Shared Scientific Library and Independent Product\\[8pt]
\large Longitudinal inference, econometric panels, partial observation,\\causal discovery and the organization of banking data}
\author{Research and product proposal}
\date{7 October 2026\\Version 2.0}
\begin{document}
\frontmatter
\maketitle
\chapter*{Prologue}
\addcontentsline{toc}{chapter}{Prologue}
A deposit campaign appears successful when balances rise among customers who
receive an offer. Yet those customers may already have been about to move money
into the bank. A collections intervention may improve repayments this month
while changing who remains eligible for assistance next month. A new branch
policy may coincide with a local economic recovery. Each situation requires an
answer about an intervention, a population and a period of time.

The banking group has much of the raw information needed to investigate these
questions. It also has many useful statistical libraries. What it commonly lacks
is a reliable way to preserve the meaning of an analysis as it moves between
data preparation, estimation, policy comparison and another application. A
customer identifier changes; a feature uses a revised balance unavailable at
the decision; a treatment is recorded as acceptance rather than assignment; an
estimator returns an effect whose population differs from the business case.
These are scientific and product failures at the same time.

This proposal develops a common platform around that problem. It begins with
bank decisions, derives identifying conditions and estimators, evaluates
available software, and specifies the data and evidence that components must
exchange. Large-scale computation executes a valid study. Partial observation
remains a modeling and identification problem. Causal discovery proposes and
constrains explanations under explicit assumptions.

The product develops three workspaces together: policy and transmission,
sequential decisions with controlled latent-state estimation, and temporal
discovery with automatic data organization. They share an auditable account
of information, interventions and financial outcomes. Their scientific
qualification is explicit: an identified effect, a conditional model scenario
and an ambiguous graph are distinct useful outputs.

The argument proceeds from identification and estimator derivation through
complete synthetic banking studies to data construction, software selection
and an equal-output investment comparison. The studies fit unknown models,
retain unfavorable results and separate numerical verification from bank
validation. The product earns its common construction cost only through
measured reuse or additional decision value relative to expert composition.

The evidence cutoff is 7 October 2026. The accompanying archive preserves 87
scholarly PDFs, three official guidance documents, a dated census of \SoftwareCount{}
open-source projects, reading notes, discovery records and source hashes.
The survey is a reproducible purposive review with explicit access and coverage
limits. It is not an exhaustive systematic review. Numerical evidence includes exact checks, fitted synthetic banking studies,
repeated-sample experiments and bounded component measurements; it does not
measure bank performance. This is an
author-reviewed proposal and reusable specifications; it does not claim a
deployed product, independent validation or institutional endorsement.

\begingroup\setlength{\parskip}{0pt}
\tableofcontents
\endgroup
\mainmatter
\part{Bank decisions and scientific foundations}
\chapter{The bank decision and the product}\label{ch:product}
\section{An intervention has an economic meaning}
Consider a bank deciding whether to offer a temporary savings-rate enhancement
to households at risk of withdrawing funds. The action is the assigned offer,
including its rate, duration, communication channel and eligibility rules.
Acceptance is a subsequent event. The outcome is incremental discounted
relationship contribution over a declared horizon, including funding value,
interest expense, servicing cost and transfers from the bank's other products.
A positive effect on the promoted account's balance need not imply a positive
effect on the banking group's value.

The study must distinguish externally acquired funds, retained funds and
internal transfers. A customer who moves money from another group account
changes the product manager's balance without necessarily creating new funding.
Discounting, funds-transfer pricing and capital allocation belong to an explicit
valuation specification. They are not quantities a causal estimator should
invent. The treatment-effect output includes a currency, valuation
date, horizon, population and contrast direction.

This running example has empirical roots. Deposit pricing and monetary
transmission vary with local market structure \citep{drechsler2017};
bank relationships exhibit persistent cross-selling patterns whose
interpretation requires a design distinguishing stickiness from preferences
\citep{basten2026}. Those studies inform mechanisms and designs. Their
coefficients are not transferable estimates of a new bank's campaign.

\section{Four scientific routes, one shared record}
The product should make four routes available within a common study record.
An experiment or adjustment study asks how a specified action changes an
outcome when assignment is randomized or conditionally exchangeable.
A panel study exploits a policy change across entities and dates, with its
own counterfactual trend or instrument restrictions.
A sequential study compares strategies whose earlier actions alter
later information and eligibility.
A discovery study investigates graph structure under stated assumptions and
returns the appropriate equivalence class or structural hypothesis.

These routes share entity definitions, clocks, outcome components, lineage and
evidence. They do not share a universal identification theorem. A graph inferred
from transactions cannot automatically validate an instrument; an accurate
hidden-state filter cannot automatically validate a treatment policy.
The central object is a \emph{study specification}: a versioned
statement of the decision, information set, target population, causal contrast,
design assumptions, estimator, uncertainty and acceptable use.

\begin{table}[htbp]\centering\small
\begin{tabularx}{\textwidth}{P{25mm}XXX}
\toprule
Application & Decision and outcome & Plausible identifying variation & Principal difficulty\\
\midrule
Deposit retention & Rate offer; group contribution and external funds & Household holdout or justified adjustment & Internal transfers, expectations, selective offers\\
Collections & Assistance and restructuring; recoveries and customer outcomes & Sequential randomization or measured-history design & Feedback, private agent information\\
Credit limits & Limit change; loss-adjusted contribution & Randomized eligible bands or threshold design & Selection, binding limits, delayed default\\
Mortgage pricing & Price or refinancing assistance; surplus and margin & Eligibility, instruments, structural demand & Take-up, competition, equilibrium\\
SME relationship & Service/price bundle; wallet contribution & Assigned offers with verified receipt & Competitor balances, selected quotes\\
Branch policy & Hours or service model; regional contribution & Staggered rollout or synthetic comparison & Anticipation, migration, common shocks\\
Treasury transmission & Funding strategy; balance-sheet outcomes & Identified shocks and bank exposure & Few shocks, changing regimes, aggregation\\
\bottomrule
\end{tabularx}
\caption{Proposed applications. Valid identifying variation must be established within the bank.}
\label{tab:applications}
\end{table}

Banking research demonstrates why different routes are needed.
\citet{agarwal2020} exploit refinancing-program eligibility rather than simply
comparing borrowers who refinance with those who do not.
\citet{mian2006} compare lending by differently affected banks to the same firm
to separate borrower demand from bank liquidity supply.
\citet{benetton2021} combine mortgage demand and lender behavior to analyze
market-wide regulation. \citet{beraja2019} explain why regional refinancing
responses alone cannot determine aggregate policy effects. These are
economically meaningful designs, not interchangeable regression templates.

\section{What the independent product sells}
The standalone offering is a workbench for bank economists, decision scientists
and model owners. A user defines a study, builds a population as of a decision
date, examines support and design diagnostics, runs declared estimators and
challengers, and exports an evidence package for a decision meeting. The
package connects a financial recommendation to data, assumptions, results and
limitations. An API offers the same operations to other applications.

The product has three substantive workspaces from the outset of development.
The \emph{policy and transmission workspace} analyzes interventions across
customer panels, branches and macroeconomic dates. The \emph{sequential
decision workspace} estimates longitudinal effects and controlled latent-state
models, then compares policies using the information actually available to
their decision maker. The \emph{discovery and data workspace} reconstructs
temporal datasets and investigates candidate mechanisms under explicit
structural restrictions. These are the central requested capabilities.
Campaign effects are one application of the scientific core.

The workspaces have different qualification levels within each release.
For example, a first release can provide useful finite-state POMDP estimation
and model-conditional scenarios while withholding an unsupported claim to
identify a historical intervention. A discovery result can be a useful PAG
or a ranked experiment proposal without becoming a unique directed causal
graph. The build sequence develops these capabilities together with the
data foundation; it does not defer the central scientific problems to an
unspecified later product.

The library owns study compilation, temporal transformations,
adapter contracts, uncertainty aggregation, graph semantics, evidence export
and validation. It should reuse established estimators rather than reproduce
their entire implementations. The independent product adds workspaces,
collaboration within bank systems, study search, controlled access and
operational support. Both forms invoke the same scientific core.

\section{A testable differentiation claim}
The strongest alternative is an expert team composing its warehouse,
catalog, notebooks and existing causal packages. That team can already perform
excellent analyses. The proposed product earns its cost only if it reduces
repeated data and review work, catches consequential errors earlier, and
increases the number of usable studies without weakening scientific standards.

Evaluation should compare matched studies with that alternative: analyst time,
time to reconstruct another analyst's result, frequency of decision-time
leakage, material review findings, compute cost and the value of decisions
supported. A polished graph display or faster fitting is insufficient.
The incremental value must be measured against what the same experts could
achieve with existing tools.

The defensible assets are banking temporal semantics, reusable
financial outcome definitions, accepted study patterns, and
evidence that travels between applications. A proprietary estimator collection
is less compelling given the available open-source ecosystem.

\section{Institutional setting}
A universal bank is not one homogeneous dataset. Legal entities, booking
locations, currencies, product systems and access rights constrain analysis.
The default deployment brings computation to approved bank data,
exports authorized results, and records the population each result
covers. Cross-border pooling is an explicit design decision. A group-wide
report can combine compatible local estimates without unrestricted
movement of transaction records.

The initial sponsor should be a central analytics or economics function with
one retail product owner and named data owners. Model risk, security and data
governance participate in acceptance. They define permissible
use, validate material models, and determine which artifacts another business
unit may rely on. No named banking group is asserted to have requested or
accepted this design.

\section{The decision system the bank needs}
A bank-wide platform must answer questions at several levels. A collections
manager compares assistance strategies for an eligible cohort. A funding team
asks how a rate change alters deposit persistence and the cost of funds. A
chief economist estimates transmission across regions and banking entities.
A model-development team asks whether new measurements can distinguish two
causal explanations. The same transaction may support all four questions,
but the interventions, populations and independent units differ.

The common scientific record therefore has six linked objects. The
\emph{population} records who was eligible before a decision and preserves
later departures. The \emph{information set} reconstructs what the actor
could know. The \emph{intervention} identifies assignment, execution,
duration and permissible subsequent actions. The \emph{outcome account}
specifies cash, impairment, funding and customer outcomes on a dated horizon.
The \emph{identification argument} describes which features of the observed
law determine the requested contrast. The \emph{estimate} carries uncertainty,
diagnostics and the exact evidence used. Shared storage is useful because
these objects travel together, not because every dataset becomes one table.

For example, a credit-limit application may consume a treatment-effect
estimate from this library and a cash-flow engine from another repository.
The effect might concern utilization after an assigned limit increase among
eligible current cardholders. The valuation engine must not interpret it as
an effect on new-account acceptance or on default among all applicants.
The product checks those declared types and retains the empirical study that
supports them. Human and model review must still examine whether the declared
types correctly describe the bank's source data. Mechanical compatibility
is necessary for composition and insufficient for truth.

\section{A usable end-to-end research journey}
Suppose a funding economist receives a proposal for a 25-basis-point retention
offer. The first screen asks for the decision, eligible population, comparison
policy and outcome horizon. It displays an observed accounting baseline:
balances, external inflows, internal transfers and interest paid. The economist
can inspect the definitions and reconciliation totals before requesting an
effect. A large balance movement with no external funding movement is then
visible before model fitting begins.

The next step selects a design. A randomized assignment supports an intention-
to-treat analysis. A staggered branch rollout supports a cohort-time comparison
only with a defensible untreated trend. Discretionary account-level offers
need measured-history exchangeability or another identification route.
The interface explains the resulting population and comparison in the
language of the funding decision. It does not present these choices as
interchangeable algorithms ranked by predictive accuracy.

The economist then sees overlap or donor diagnostics, the main estimate,
predeclared alternatives and a financial decomposition. In the executed
rollout study in Chapter~\ref{ch:deposits}, the result is unfavorable even
though additional deposits are identified under the stipulated design.
The correct product behavior is to retain the evidence and explain the cost
of paying the higher rate on balances that would have remained. It is not
to search for a more favorable estimator. A new targeted offer is a new
policy requiring its own evaluation.

For a sequential question, the same journey opens a trajectory view. The
reader sees how assistance affects next-month stress, eligibility and later
costs. A latent-state model has an additional observation and logging-law
view: which signals measure the state, which variables the historical actor
used, and how that differs from the proposed actor. The fitted state is
presented as model-dependent uncertainty. This makes a POMDP estimation
product useful to a bank researcher before any automated action is authorized.

\section{Why shared construction may be worth buying}
The strongest integration benefit is the removal of repeated scientific
translation. A team today often reconstructs availability dates, rewrites
entity-level folds, interprets a package's treatment coding and reassembles
an evidence report for every study. The proposed library would preserve those
decisions in executable study objects, allowing a new estimator to be a
challenger for the same target. It would also make a source correction
propagate to the affected studies instead of silently changing a dashboard.

That benefit is falsifiable. If every business unit needs a different outcome
ontology, if adapters consume as much work as a custom analysis, or if the
evidence package cannot shorten a real review, a shared product may have
little advantage over a capable internal team using open source. The economic
comparison in Chapter~\ref{ch:economics} holds scientific quality and source
reconciliation obligations fixed. It counts shared costs once, measures the
marginal application cost and shows the portfolios for which reuse fails.

The standalone buyer is consequently an institution with a portfolio of
repeated causal questions and a team able to own the assumptions. It is not
a buyer seeking an automatic answer to every causal question from raw data.
The initial procurement can be a scientific platform and supported analytical
applications, with the bank retaining authoritative data and model decisions.
The same core can be installed as a library in the user's repositories.
Both forms should produce the same estimand, score contributions and report
on a shared conformance fixture.

\section{What this monograph constructs}
The scientific chapters work from intervention definitions to estimators,
dependence-aware uncertainty, partially observed dynamics and graph search.
The banking studies then use those constructions in complete synthetic
examples. They distinguish estimating balances from valuing funding,
evaluating a fixed policy from choosing one, and fitting a hidden-state law
from identifying a causal action. The data and software chapters explain
what existing tools supply and what a reusable platform must add.

The numerical studies are deliberately transparent. Their assumptions,
financial units, training information and exact or model-based comparison
values are visible. They establish that the stated calculations can be
constructed and expose several failure modes. They do not estimate any
named bank's business opportunity. The proposal's next investment is to
test the same scientific and engineering contracts on owned bank data,
with equivalent-task comparisons against existing practice.

\chapter{Identification before estimation}\label{ch:identification}
Identification determines whether a requested answer is a function of the
available distribution and stated assumptions. Estimation approximates that
function from a finite sample. Computation evaluates an estimator. The three
questions fail in different ways: an observational equivalence cannot be
repaired by a larger cluster, a noisy identified effect may benefit from more
independent observations, and a slow valid estimator may benefit from better
software. A product must diagnose which problem it is facing.
\section{The estimand is a product contract}
For an eligible household let $A\in\{0,1\}$ denote assigned treatment, $X$
information available before assignment, and $Y(a)$ the horizon outcome under
assignment $a$. The average treatment effect is
\begin{equation}
 \theta=\E[Y(1)-Y(0)],\qquad
 \tau(x)=\E[Y(1)-Y(0)\mid X=x].
 \label{eq:ate}
\end{equation}
The expectation refers to a declared population. An effect among existing
high-balance customers is not an effect among all households. Changing
eligibility, outcome horizon or economic valuation creates a new
estimand. For offers, the intention-to-treat effect concerns assignment;
the effect of receipt requires a separate compliance design.

Potential outcomes make the contrast explicit. Structural causal models
describe how interventions replace mechanisms. A single-world intervention
graph links the two by splitting a treatment node into its fixed intervention
value and its random component \citep{hernan2020,peters2017,swig2013}.
The product should store the estimand and, where used, a graph with
substantive variable definitions. Neither notation substitutes for a plausible
assignment mechanism.

Under consistency, conditional exchangeability
$\{Y(0),Y(1)\}\perp A\mid X$, and positivity
$0<e(X)=\Prob(A=1\mid X)<1$ on the target population,
\begin{equation}
 \theta=\E[m_1(X)-m_0(X)],\qquad m_a(x)=\E[Y\mid A=a,X=x].
 \label{eq:adjust}
\end{equation}
Consistency includes the absence of consequential undisclosed treatment
versions. An offer is inadequately defined if its price, delivery and duration
are unspecified. Exchangeability concerns actual decision information:
a relationship manager's unrecorded knowledge can invalidate it even if the
warehouse contains thousands of predictive fields. Positivity concerns actions
within histories; large row counts do not create support for actions never taken.

\section{Experiments and observational studies}
Randomize at household level when accounts share money, or at branch/region
level when capacity or communications spill across customers. Preserve actual
allocation probabilities and distinguish assignment, delivery and take-up.
Experiments also need a prespecified horizon, contamination diagnostics and
attrition analysis.

Adaptive experiments require their time-varying assignment probabilities and
inference valid for the adaptation and stopping rule. Repeatedly examining
ordinary fixed-sample confidence intervals does not make them valid under
optional stopping. In a cluster experiment, the number of randomized
clusters constrains inference even with millions of transactions per cluster.
Randomization inference uses the actual allocation mechanism, not an
analyst's preferred regression model.

Observational adjustment is appropriate when assignment information and
overlap are defensible. Diagnostics include balance, propensity support,
outcome sensitivity and negative controls with justified causal roles.
They can discover defects but do not prove that all confounders were observed.
The uncertainty object should distinguish sampling uncertainty from
assumptions about omitted confounding \citep{ovb2022}.

\section{Why orthogonal estimation is useful}
Flexible prediction helps estimate nuisance functions $m_a$ and $e$, but
direct substitution can introduce regularization bias. An augmented
inverse-probability score is
\begin{equation}
 \phi(O;m,e)=m_1(X)-m_0(X)
 +\frac{A\{Y-m_1(X)\}}{e(X)}
 -\frac{(1-A)\{Y-m_0(X)\}}{1-e(X)}.
 \label{eq:aipw}
\end{equation}
Estimate nuisances outside the evaluation fold and average held-out scores.
Cross-fitting limits direct dependence between nuisance overfitting and
evaluation residuals. Orthogonality reduces first-order sensitivity to
nuisance error \citep{dml2018,wager2025}.

\begin{proposition}[Conditional bias of the AIPW score]
With true nuisances $e,m_0,m_1$, fixed candidate nuisances
$\widetilde e,\widetilde m_0,\widetilde m_1$ and denominators bounded away
from zero,
\begin{align}
&\E[\phi(O;\widetilde m,\widetilde e)\mid X]-(m_1-m_0)(X)\nonumber\\
&\quad=(\widetilde e-e)
 \left\{\frac{\widetilde m_1-m_1}{\widetilde e}
       +\frac{\widetilde m_0-m_0}{1-\widetilde e}\right\}(X).
 \label{eq:drbias}
\end{align}
\end{proposition}
\begin{proof}
Conditioning on $X$ replaces the treated residual term by
$e(m_1-\widetilde m_1)/\widetilde e$ and the untreated residual term by
$-(1-e)(m_0-\widetilde m_0)/(1-\widetilde e)$. Collecting coefficients of
each outcome error gives the expression.
\end{proof}

The bias vanishes if either nuisance component is correct; otherwise it
depends on products of errors. With overlap bounds, Cauchy--Schwarz bounds
these terms by products of nuisance-error norms. Under the remaining
independent-sampling and moment conditions, products
$o_p(n^{-1/2})$ support asymptotically linear inference. This is not a
result for arbitrarily dependent transaction rows.

For independent observations, the empirical held-out score variance divided
by $n$ estimates the variance of the mean under suitable conditions.
Household trajectories and common shocks require appropriate score
construction and covariance aggregation. The interface therefore requires
both an estimator and an inference specification.

\section{A campaign example with a known answer}
The synthetic construction sets $X\sim\mathrm{Bernoulli}(1/2)$,
$e(X)=0.2+0.6X$, and
\[
 Y=100+80X+(10-4X)A+\varepsilon,\qquad
 \varepsilon\sim N(0,20^2).
\]
The true effect is $8$. In the population,
$\Prob(X=1\mid A=1)=0.8$ and $\Prob(X=1\mid A=0)=0.2$.
The observed difference is $80(0.8-0.2)+10-4(0.8)=54.8$:
assignment targets customers with larger untreated outcomes.

In 40,000 independent synthetic households, five-fold saturated nuisance
fits give AIPW $\NumStaticATE$ with standard error $\NumStaticSE$,
compared with the raw difference $\NumStaticNaive$.
The 95\% interval is approximately $[7.974,8.969]$.
This checks a computation against a known process. One seed and an interval
containing the truth do not demonstrate coverage.

Heterogeneity matters economically. Effects are $10$ and $6$ in the two
groups. At an additional cost of $7$ per offer, treating everyone has true
incremental value $1$ per eligible customer; treating only $X=0$ has
value $1.5$. The code evaluates these fixed policies with paired AIPW scores
and simultaneous intervals over three declared comparisons.
This demonstrates policy comparison without optimizing and evaluating a
learned policy on the same outcomes.
The targeted-versus-all estimate is $0.336$ with a simultaneous interval
approximately $[-0.096,0.769]$. Although the generating model makes the
targeted policy better by $0.5$, this sample does not establish superiority
at the declared simultaneous confidence level. Both policies have positive
lower bounds relative to treating nobody. The example thus distinguishes
positive value from evidence that a more selective policy is better.

\section{Heterogeneity, instruments and thresholds}
Meta-learners and orthogonal learners estimate conditional average effects
\citep{kunzel2019,dml2018}. Compare subgroup calibration, policy value and
uncertainty, not only factual prediction. Suppress unsupported segments.
A customer's CATE is an average among comparable customers, not an
observable personal causal effect.

When assignment has unmeasured determinants, a valid instrument can offer
a different route \citep{applied2024}. With binary randomized encouragement $Z$, exclusion,
relevance, independence and monotone take-up identify the complier effect
through the Wald ratio
\begin{equation}
 \theta_{\mathrm{LATE}}=
 \frac{\E[Y\mid Z=1]-\E[Y\mid Z=0]}
      {\E[A\mid Z=1]-\E[A\mid Z=0]}.
 \label{eq:wald}
\end{equation}
Monotonicity and exclusion are not consequences of a strong first stage.
Weak first stages need suitable inference. Competitor rates may directly
affect switching and therefore fail exclusion. A price instrument valid in
one market can fail in another because information and policy differ.

Regression discontinuity uses a treatment rule at a threshold $c$ of a
running variable $R$. With continuity of potential-outcome means, the
outcome jump identifies a local effect in a sharp design. In a fuzzy design
divide the outcome jump by the treatment-probability jump, with appropriate
compliance restrictions. Retain the score and rule version used at the time;
inspect manipulation, competing thresholds and bandwidth sensitivity.
The local effect does not by itself value a policy far from the threshold.

\section{What a result is allowed to say}
Every result identifies its design, target, comparison, assumptions and
uncertainty. Its scientific state may be \emph{identified under declared
assumptions}, \emph{partially identified}, \emph{unsupported by observed
action overlap}, or \emph{not identified under the supplied model}.
Numerical convergence is a separate field. A consuming application can then
abstain or request a better experiment without mistaking prediction for
identification.

\section{Why the augmented score works}
The AIPW score has a useful interpretation for a banker. It first predicts
each customer's outcome under each action, then corrects the prediction
using the observed residual among customers who actually received that
action. A rare action receives a larger correction because fewer observed
customers represent its potential outcome. The estimator does not create
counterfactual observations; it uses exchangeability to give those residuals
their intervention interpretation.

Equation~\eqref{eq:drbias} makes that forecast-and-correction account
precise. Conditioning on training information treats fitted functions as
fixed for an independent evaluation observation; averaging the conditional
remainder gives the population bias. The relevant protection concerns
the product of nuisance errors, with the overlap and inference conditions
already stated. It does not assert that two inaccurate learners
necessarily cancel.

For bank data, overlap can vary sharply by product, jurisdiction and
eligibility rule. A risk department may never approve a limit increase for
a certain population. No flexible outcome learner makes the intervention
identified there from those records. Trimming changes the population;
capping weights changes the estimator's bias; replacing a deterministic
policy by a feasible stochastic policy changes the estimand. Each can be
reasonable, but the report must describe the actual choice and recompute
the economic target accordingly.

\section{Residualization can change the target}
Partially linear DML is attractive because it combines flexible nuisance
prediction with a low-dimensional coefficient. Its canonical model is
\[
 Y=\theta_0 A+g_0(X)+U,\quad \E[U\mid A,X]=0,
 \qquad A=e_0(X)+V,\quad \E[V\mid X]=0.
\]
Let $\ell_0(X)=\E[Y\mid X]$. Regress the residualized outcome
$Y-\ell_0(X)$ on $A-e_0(X)$, or solve the orthogonal moment
\begin{equation}
 \E[(A-e_0(X))\{Y-\ell_0(X)-\theta(A-e_0(X))\}]=0.
 \label{eq:plr-score}
\end{equation}
Cross-fitting allows flexible estimates of $e_0$ and $\ell_0$ while
retaining a tractable target under the model and rate assumptions.

Now suppose the true effect is $\tau(X)$ rather than a constant. For a
binary action with conditional mean $m_0(X)+A\tau(X)$, the population
solution of the same moment is
\begin{equation}
 \theta_{\mathrm{PLR}}=
 \frac{\E[e_0(X)\{1-e_0(X)\}\tau(X)]}
      {\E[e_0(X)\{1-e_0(X)\}]}.
 \label{eq:plr-weighted}
\end{equation}
To obtain it, condition the numerator
$\E[(A-e_0)(Y-\ell_0)]$ on $X$; it equals the conditional
treatment variance times $\tau(X)$. The denominator averages that same
variance. The resulting coefficient emphasizes customers with treatment
probabilities near one half. It need not equal the eligible population's
ATE or an effect among treated customers.

This distinction is operational. A bank could compare an AIPW ATE to a
partially linear DML coefficient, see different values and incorrectly
attribute the difference to software quality. The estimates may target
different weighted populations. A challenger comparison must either align
the target or state the different question. The adapter registry therefore
needs an estimand declaration, including aggregation weights, rather than
only a label such as ``DML supported.''

\section{Selecting covariates requires temporal and causal reasoning}
A variable measured before an action can still be inappropriate to adjust
for if it is a collider or an instrument that aggravates variance without
addressing confounding. A variable measured after an action can be a mediator,
a common effect or a confounder of a later decision. The temporal clock is
necessary, but it does not decide every adjustment question.

For example, an offer $A_0$ changes engagement $L_1$, which affects both
future assistance $A_1$ and final recoveries $Y$. Adjusting for $L_1$
in a single regression intended to estimate the total $A_0$ effect can
block part of that effect. Failing to adjust for it when evaluating $A_1$
can leave confounding. Chapter~\ref{ch:longitudinal} resolves this through
the sequential law. It cannot be resolved by a global ``pre-treatment
features'' list shared by all decision dates.

The same issue appears in customer service data. Complaints can result from
both an offer and latent dissatisfaction. Conditioning on complaint receipt
may select a subgroup whose latent dissatisfaction differs across offers.
An analysis of customers with a recorded complaint then targets neither
the original cohort nor an automatically valid mediator effect. A mediation
study needs its own intervention definition and identifying assumptions;
the total assignment effect remains useful even when the mechanism is
unresolved \citep{hernan2020,swig2013}.

The proposed data organizer can suggest which columns describe complaints,
assignments and account outcomes. It can identify timestamps and draw
candidate graphs. It cannot infer that a complaint is an acceptable
adjustment variable from its business name. The study must retain the
reason it is included, excluded or treated as a time-varying variable.

\section{Measurement, missingness and transport are part of identification}
Suppose true treatment is an assigned offer but the source records only an
opened message. The measurement is then affected by attention, device use
and customer preference. Relabeling message opening as treatment changes
the intervention. Recovering a hidden treatment from surrogates requires
additional measurement and proxy restrictions; recent work does not justify
universal recovery of missing action logs \citep{hidden2026}.

Likewise, a missing competitor balance is not a zero balance. If observed
bank balance equals liquid wealth times bank share, $B=W\omega$, the
same $B$ is compatible with many $(W,\omega)$ pairs. Neither precise
measurement of $B$ nor a flexible causal learner identifies the outside
wallet without additional information. Data-recovery draws can propagate
uncertainty under a supplied missingness model, but they do not establish
that model. The corrupted-data literature makes clear that different
measurement and observation operators require different identification
conditions \citep{corrupt2022}.

For a source population $S=1$ and a target population $S=0$, transporting
an effect requires more than matching feature names. If conditional causal
effects are stable and every target $X$ has source support, the target
ATE is $\E[\tau(X)\mid S=0]$. This is a useful standardization formula.
It does not follow merely because a classifier has difficulty distinguishing
the populations. Product terms, laws, delivery systems, hidden effect
modifiers and competitive responses may change the response mechanism.
The target population and stability assumptions must remain attached to
the transported estimate.

\section{From an effect to a financial comparison}
Let $C_k(a)$ be dated cash or economic contribution components with signs
defined from the group's perspective, and let $d_k$ be fixed discount and
currency-conversion factors. Define $Y(a)=\sum_k d_k C_k(a)$.
Linearity gives
\[
 \E[Y(1)-Y(0)]=\sum_k d_k\E[C_k(1)-C_k(0)].
\]
However, uncertainty is joint. If $\widehat{\boldsymbol\theta}$ estimates
the component effects with covariance $\widehat\Sigma$, the variance
of the financial contrast is $\boldsymbol d'\widehat\Sigma\boldsymbol d$.
Adding component confidence intervals or treating correlated components as
independent does not reproduce this variance. Fitting all components on
the same folds and retaining customer or cluster score vectors allows
the covariance to be reconstructed.

This formula conditions on fixed financial factors. An uncertain transfer
price, loss severity or capital charge adds parameter and scenario uncertainty.
If the policy changes aggregate market rates, the factors can themselves
be intervention-dependent. Then a common fixed conversion is an approximation
whose adequacy needs an economic argument. Deposit principal is a funding
quantity; its full value is not an immediate profit. Similarly, a loan
advance is an outflow paired with a claim, not an expense that should be
counted again when expected principal losses are recorded.

Scientific uncertainty and economic sensitivity should be displayed separately
and, where a defensible joint model exists, combined. A narrow sampling interval
under one accounting convention does not settle a decision that reverses
under plausible funding-cost or customer-response assumptions.

\chapter{Panels, policy changes and banking economics}\label{ch:panels}
\section{Repeated measurement is not a design}
Bank panels combine households, firms, accounts, branches and calendar dates.
They separate persistent differences from changes over time but do not
automatically remove time-varying confounding. In
\begin{equation}
 Y_{it}=\alpha_i+\lambda_t+\beta D_{it}+u_{it},
 \label{eq:twfe}
\end{equation}
unit effects absorb time-constant additive heterogeneity and time effects
absorb common additive shocks. Neither absorbs a local recession affecting
both rollout and outcomes. The coefficient must match the economic target.
Modern designs include group-time effects, imputation, synthetic controls
and interactive factors \citep{panelreview2024,roth2023}.

\subsection{What the within transformation actually removes}
For a balanced panel define
$\widetilde Z_{it}=Z_{it}-\bar Z_{i\cdot}-\bar Z_{\cdot t}
+\bar Z_{\cdot\cdot}$. Applying this linear operation to
Equation~\eqref{eq:twfe} removes both additive effects. A scalar coefficient
is estimated by
\begin{equation}
 \widehat\beta=
 \frac{\sum_{it}\widetilde D_{it}\widetilde Y_{it}}
      {\sum_{it}\widetilde D_{it}^{2}}.
 \label{eq:within}
\end{equation}
The denominator shows the source of information: treatment variation left
after unit and date effects. A national policy adopted by every bank at the
same time has no such variation in this specification. Millions of accounts
cannot rescue that rank failure. For an unbalanced panel, successive simple
demeaning does not generally equal the exact joint projection; a suitable
least-squares or converged absorption routine is needed.

A sufficient model-based condition for the usual within moment is strict
exogeneity, $\E[u_{it}\mid D_{i1},\ldots,D_{iT},\alpha_i]=0$,
together with the appropriate common-effect specification. This excludes
feedback from an unobserved outcome shock into later treatment. Bank
managers often change offers after unexpected withdrawals, so the condition
can fail even if today's treatment predates today's recorded outcome.
Predeterminedness is weaker and does not in general justify the same within
estimator. The policy-study route should therefore start with its design,
not infer exogeneity from a familiar regression formula.

The transformation also cannot remove an arbitrary function of latent
heterogeneity. If untreated deposits are
$Y_{it}(0)=\exp(\alpha_i+f_t)+\epsilon_{it}$, subtracting unit and
date means leaves interactions between the customer and the macro state.
A machine learner fitted to demeaned data inherits that problem. The panel
DML discussion in \citet{paneldml2024} is relevant precisely because
nonlinear latent heterogeneity and dependence complicate otherwise useful
orthogonal-learning constructions.

\section{Staggered adoption and explicit aggregation}
Let $G_i$ be the first treatment date, with $G_i=\infty$ for never-treated
entities. An absorbing rollout defines
\[
 \ATT(g,t)=\E[Y_{it}(g)-Y_{it}(\infty)\mid G_i=g],\qquad t\geq g.
\]
With no anticipation and conditional parallel trends between cohort $g$ and
valid controls $C$,
\begin{align}
 \ATT(g,t)=\E\big[&
 \E[Y_t-Y_{g-1}\mid G=g,X]\nonumber\\
 &-\E[Y_t-Y_{g-1}\mid C=1,X]\mid G=g\big].
 \label{eq:gtatt}
\end{align}
Controls remain untreated over the comparison interval and have support for
the treated cohort. The outer expectation targets the treated covariate
distribution. Never-treated and not-yet-treated comparison groups imply
different assumptions and available dates \citep{callaway2021}.

Aggregation is part of the estimand:
$\theta=\sum_{g,t}w_{gt}\ATT(g,t)$ for declared nonnegative weights
summing to one. Equal customer, branch and treated-period weights answer
different questions. Event-time effects also change cohort composition when
long follow-up is available only for early adopters.

\citet{bacon2021} decompose TWFE into two-by-two comparisons.
\citet{sun2021} show that lead/lag coefficients can mix effects across
event times. \citet{borusyak2024} make target weights and the untreated
outcome model explicit through imputation. The product should expose these
alternatives rather than label every fixed-effects regression valid DiD.

\begin{example}[Perfect untreated trends, misleading TWFE]
Three equal cohorts are observed over four dates. The early cohort adopts
at date 2 with effects $(1,4,8)$; the late cohort adopts at date 4 with
effect $1$; the third never adopts. Untreated outcomes equal a cohort
intercept plus a common linear trend. The four-treated-cell average is
$(1+4+8+1)/4=3.5$.

Residualizing treatment on cohort and time effects yields underlying effect
weights $(0.3,0.3,-0.1,0.5)$. Thus TWFE is
$0.3+1.2-0.8+0.5=1.2$. Comparisons with the never-treated cohort recover
every group-time effect and their average $3.5$.
\end{example}

\begin{figure}[htbp]\centering
\includegraphics[width=\textwidth]{figures/panel_bias.pdf}
\caption{Exact synthetic panel calculation. Contamination occurs despite correct untreated parallel trends.}
\label{fig:panel}
\end{figure}

This is a target and comparison problem; clustering standard errors does
not repair it. A rollout with gradually growing adoption can exhibit this
pattern.

\subsection{A doubly robust cohort comparison}
For one cohort and comparison interval, retain only cohort $g$ and controls
that are untreated over the interval. Let $D=1$ indicate cohort $g$,
$\Delta Y=Y_t-Y_{g-1}$, $p(X)=\Prob(D=1\mid X)$, and
$m_0(X)=\E[\Delta Y\mid D=0,X]$ in this restricted population.
Write $q=\Prob(D=1)$. A population score for the cohort ATT is
\begin{equation}
 \frac{1}{q}\E\left[
 D\{\Delta Y-m_0(X)\}
 -(1-D)\frac{p(X)}{1-p(X)}\{\Delta Y-m_0(X)\}
 \right].
 \label{eq:did-dr}
\end{equation}
Under no anticipation, overlap and conditional parallel trends, the first
term compares treated changes with their counterfactual untreated change.
The second corrects a fitted untreated-change regression using reweighted
control residuals. Its weights map control covariates to the treated
distribution. With either the relevant propensity or untreated-change
regression correct, the mean has the desired robustness; a sample version
also estimates the cohort share and carries its uncertainty. This is a
simple form of the identification and estimation logic developed by
\citet{callaway2021}.

Notice what differs from a static treatment-effect model. The identifying
assumption concerns untreated \emph{changes}, not the levels of untreated
outcomes. A treated branch can have a permanently different customer mix
and deposit level. It becomes a problem when its untreated change would
differ after the rollout, conditional on the stated covariates. Including
post-rollout customer composition as a control can itself remove a policy
mechanism, since the rollout may attract a different customer population.

The common score representation is useful for software. Once each
$\ATT(g,t)$ has cluster contributions, an aggregate effect uses the
same declared weights on those contributions. A covariance matrix over
event times then supports simultaneous bands and financial aggregation.
Repeatedly calling a package for separate dates and discarding cross-date
covariance loses information required to value a multi-month policy.

\subsection{The comparison set is an economic choice}
Never-treated branches may be structurally unlike rollout branches.
Not-yet-treated branches can be more comparable but can anticipate a
scheduled change. A one-month anticipation allowance changes the last valid
baseline and excludes controls close to their own adoption. When every
branch eventually adopts, late effects may have no untreated controls.
The honest output is then a shorter identified horizon or a more model-
dependent extrapolation, rather than a graph with silently declining support.

Event-time aggregation needs a second decision. At event month 12 only early
cohorts may remain observable. Comparing that average with event month 1
can conflate effect dynamics with a different cohort mix. A balanced event
window estimates the same cohorts at every horizon but discards information.
The report should show both the target weights and the set of contributing
cohorts. A bank sponsor can then decide whether the retained target
population matches the planned rollout.

\section{Trend sensitivity and reversals}
An insignificant pretrend test does not establish parallel trends. Low power
and post-selection affect interpretation. \citet{rambachan2023} allow
specified smoothness or relative-magnitude deviations. Show conclusions
over a justified sensitivity range and the point at which the decision
changes.

\subsection{Sensitivity gives the assumption an economic scale}
Suppose a simple two-period estimate is $\widehat\tau$. Let
$\delta$ be the treated-minus-control change in untreated potential
outcomes, so that $\E[\widehat\tau]=\tau+\delta$ under the rest of
the design. If substantive reasoning limits $\delta$ to $[-M,M]$,
the compatible treatment effects shift by that range. An approximate
sampling interval $[L,U]$ becomes a conservative union
$[L-M,U+M]$ when that bound is treated as fixed. This elementary
construction is deliberately separate from the more structured event-study
confidence sets in \citet{rambachan2023}.

The parameter should be expressed in business units. If a positive policy
decision reverses when untreated deposits could have grown only USD\,20
per eligible customer more rapidly in treated branches, that is useful
information for a funding economist. A sensitivity parameter whose
relation to deposit growth is hidden behind a default slider is less
useful. Historical untreated variation, known local shocks and institutional
rollout rules can motivate a range, but do not prove its bound.

Pre-treatment coefficients have two roles: they can reveal contradictions
with a design, and they can help calibrate a sensitivity analysis. Passing
a pretrend test is not a license to assume $M=0$. Selection of the final
design after searching for quiet pretrends can distort the subsequent
interval. The product records the original study plan and exploratory
alternatives so that the final inference can be interpreted honestly.

Bank treatments also change intensity or reverse: rates move repeatedly,
contact policies stop and restart, and branch services expand in stages.
Absorbing binary-adoption methods may be inappropriate.
\citet{dechaisemartin2024}, in the archived May 2026 revision, study
intertemporal effects with status-quo comparisons for groups sharing initial
treatment. No anticipation, comparison support and treatment-path meaning
remain essential. Preserve intensity and version rather than collapsing
every history to first treatment.

\section{Dynamic panels and latent factors}
In
\[
 Y_{it}=\rho Y_{i,t-1}+\beta A_{it}+\alpha_i+\varepsilon_{it},
\]
demeaning correlates the transformed lagged outcome and disturbance in a
short panel. Differencing removes $\alpha_i$ but not endogeneity. Under
suitable serial-correlation and initial-condition restrictions, lagged
levels instrument differenced equations. Moments such as
$\E[Y_{i,t-s}\Delta\varepsilon_{it}]=0$, $s\geq2$, require justification.
\citet{arellano1991} develop relevant specification tests. First-order
correlation in differenced errors may be mechanical; higher-order tests,
instrument validity and sensitivity to instrument count matter.

\subsection{Deriving an internal-instrument moment}
First differencing the dynamic model gives
\[
 \Delta Y_{it}=\rho\Delta Y_{i,t-1}+\beta\Delta A_{it}
              +\Delta\varepsilon_{it}.
\]
The regressor $\Delta Y_{i,t-1}$ contains $\varepsilon_{i,t-1}$,
whereas the disturbance contains $-\varepsilon_{i,t-1}$.
Even independent level errors therefore create a correlation. Under serially
uncorrelated level innovations and suitable initial conditions, $Y_{i,t-2}$
is orthogonal to both $\varepsilon_{it}$ and $\varepsilon_{i,t-1}$.
It can instrument $\Delta Y_{i,t-1}$. Older levels provide additional
moments, but the allowable lags for $A$ depend on whether actions are
strictly exogenous, predetermined or contemporaneously endogenous.

Stack admissible instruments in $Z_i$ and differenced residuals in
$v_i(\vartheta)$. The sample moment is
$\bar g(\vartheta)=N^{-1}\sum_i Z_i'v_i(\vartheta)$, and GMM
minimizes $\bar g(\vartheta)'W\bar g(\vartheta)$ for a suitable
weighting matrix. Robust covariance estimation uses unit-level moment
contributions; the two-step weighting matrix and its finite-sample
behavior need explicit treatment. This construction explains why dynamic
panel software requires the instrument set, transformation and inference
convention to be part of the saved model \citep{arellano1991}.

Persistence creates another problem. If levels are highly persistent,
lagged levels may be weak instruments for changes. Enlarging the instrument
set can overfit endogenous changes and make overidentification diagnostics
unhelpful. A sensible comparison limits lag depth, considers collapsed
instrument sets and reports sensitivity; it does not select the specification
with the most reassuring diagnostic $p$-value. A system of additional
level-equation moments needs additional stationarity or initial-condition
restrictions. It is a different identifying model, not an automatic
upgrade whenever difference GMM is imprecise.

Interactive-factor models
$Y_{it}(0)=\alpha_i+\lambda_t+\ell_i^\top f_t+\varepsilon_{it}$
permit different exposures to common shocks. Factor rank, untreated-data
availability, treatment selection and extrapolation govern their use.
An estimated factor is not automatically a measured confounder.

Panel DML combines flexible nuisances with panel structure, but nonlinear
latent heterogeneity need not disappear under a within transformation.
\citet{paneldml2024} examine point estimation in simulation.
\citet{cluster2021} establish multiway-cluster DML results under specified
exchangeability and dissociation conditions. These are not guarantees
for arbitrary serial dependence.

\section{Synthetic controls and structural time series}
For a small number of policy changes, construct
$\widehat Y_{1t}(0)=\sum_{j\in\mathcal D}w_jY_{jt}$ from suitable
untreated donors, often with nonnegative weights summing to one.
Pre-fit assesses one aspect of plausibility; interpolation bias and
spillovers remain concerns \citep{abadie2010,abadie2021}.
Outcome-model augmentation can correct imperfect fit while implying
negative weights \citep{benmichael2021}.

Synthetic DiD combines unit and time weighting \citep{synthdid2021}.
Its inspected package documents block treatment: paper extensions do not
prove implementation support. Bayesian structural time-series methods offer
another counterfactual forecast \citep{brodersen2015}. Controls must be
unaffected and their counterfactual relationship must remain credible.

Nearby branches receiving displaced customers after a closure may be
invalid donors. National campaign controls may share exposure. Placebo
dates, donor sensitivity and pre-period holdouts can challenge a design;
they cannot create an untreated comparison where none exists.

\subsection{What a donor fit identifies}
One useful model writes untreated outcomes as
$Y_{it}(0)=\alpha_i+\lambda_t+\ell_i'f_t+\epsilon_{it}$.
If donor weights reproduce the treated unit's intercept and factor loading,
their weighted untreated mean reproduces its systematic counterfactual.
A long and informative pre-period can constrain those loadings through
outcome fit. It does not ensure that a novel post-period factor lies in
the same span or that donor outcomes remain unaffected.

The conventional nonnegative simplex constraint controls extrapolation:
the synthetic unit lies in the donors' convex hull. If the treated unit is
outside that hull, poor fit can persist even with many donors. Augmentation
uses an outcome model to correct residual imbalance and can permit negative
effective weights \citep{benmichael2021}. It exchanges a stricter interpolation
restriction for a model-dependent extrapolation. Synthetic DiD balances
both units and dates in its main block-adoption setting
\citep{synthdid2021}. These are distinct responses to imperfect comparison
data; neither is uniformly preferable.

Inference must match the design. A placebo distribution across donors can
be informative when units are suitably exchangeable under the null and
fit differences are handled. It is not automatically a randomized test
when a regulator deliberately selected one troubled institution.
Time-series posterior predictive intervals condition on the counterfactual
model and its unaffected controls \citep{brodersen2015}. Their coverage is
not a test of the causal exclusion. A bank report should distinguish
sampling, model and design uncertainty instead of merging all three into
one shaded band labeled ``confidence.''

\section{Financial interpretation and uncertainty}
Independent policy changes constrain information more than raw account
counts. \citet{bertrand2004} demonstrate the consequences of ignored
serial correlation. Preserve assignment clusters and disclose few-cluster
or short-time limitations.

The lending design of \citet{mian2006} compares bank supply among
multi-bank borrowers; it does not directly estimate all borrowers' losses.
The mortgage equilibrium analysis of \citet{benetton2021} adds demand
and supply assumptions. Regional effects in \citet{beraja2019} cannot
simply be summed into an aggregate monetary-policy effect. Preserve the
distinction between an identified local effect, a transported effect and
a structural scenario.

\part{Repeated decisions, hidden states and discovery}
\chapter{Longitudinal interventions and repeated decisions}\label{ch:longitudinal}
\section{Why a sequence changes the problem}
In collections, yesterday's assistance affects today's arrears, today's arrears
influence the next assistance decision, and both predict eventual repayment.
Today's arrears are therefore both a consequence of earlier treatment and
a confounder for later treatment. Regressing repayment on every action and
every intermediate covariate can remove part of the total effect of earlier
assistance. Omitting the covariate can confound the later effect.

Let $L_t$ be information available immediately before decision $A_t$,
$H_t=(L_0,A_0,\ldots,A_{t-1},L_t)$, and let decisions occur at
$t=0,\ldots,T-1$. A policy $g_t(a\mid h)$ specifies a distribution over
feasible actions at each history. A deterministic rule is the special case
concentrated at $a=d_t(h)$. The target is
\begin{equation}
 \psi(g)=\E[Y^g],\qquad \Delta(g,g_0)=\psi(g)-\psi(g_0),
 \label{eq:policytarget}
\end{equation}
with population, horizon, outcome and eligibility stated. A comparison of
one action at baseline differs from a comparison of sustained strategies.

\section{Identification through the longitudinal g-formula}
Assume consistency, sequential exchangeability of actions given $H_t$,
and policy support: $g_t(a\mid h)>0$ only where the behavior probability
$e_t(a\mid h)>0$ at histories reached under $g$.
With no censoring for now, intervention replaces the assignment factors
in the observed law:
\begin{equation}
 p^g(l_0,a_0,\ldots,l_T)
 =p(l_0)\prod_{t=0}^{T-1}
   g_t(a_t\mid h_t)p(l_{t+1}\mid h_t,a_t).
 \label{eq:gformula}
\end{equation}
Integrating $Y$, included in terminal information $L_T$, under this law
gives $\psi(g)$ \citep{hernan2020}. The replacement is justified by causal
assumptions, not by a fitted transition model alone.

Two computational forms follow. G-computation fits outcome/covariate
transitions and simulates or integrates the intervention. Inverse weighting
uses the change of measure
\begin{equation}
 \psi(g)=\E\!\left[Y\prod_{t=0}^{T-1}
       \frac{g_t(A_t\mid H_t)}{e_t(A_t\mid H_t)}\right].
 \label{eq:longipw}
\end{equation}
The product cancels observed assignment factors, leaving the intervention
law. This derivation exposes both why weighting works and why it becomes
unstable when supported trajectories are rare.

Stabilization, clipping and normalization are different operations.
Clipping importance weights generally introduces bias relative to the
unclipped estimand; report the threshold and its effect. Effective sample
size $(\sum_iw_i)^2/\sum_iw_i^2$ describes weight concentration but is not
a replacement for a dependence-aware variance estimator or a proof of
identification.

\section{Recursive regression and orthogonal scores}
For a fixed policy $g$, define $V_T(H_T)=Y$ and work backward:
\begin{align}
 Q_t(h,a)&=\E[V_{t+1}(H_{t+1})\mid H_t=h,A_t=a],\\
 V_t(h)&=\sum_a g_t(a\mid h)Q_t(h,a),\qquad \psi(g)=\E[V_0(H_0)].
 \label{eq:recursiveq}
\end{align}
For continuous actions the sum becomes an integral. Define
$W_t^g=\prod_{s=0}^t g_s(A_s\mid H_s)/e_s(A_s\mid H_s)$.
A population orthogonal representation has score
\begin{equation}
 \Phi_g=V_0(H_0)+
 \sum_{t=0}^{T-1}W_t^g
       \{V_{t+1}(H_{t+1})-Q_t(H_t,A_t)\}.
 \label{eq:longdr}
\end{equation}
At true nuisance functions, each weighted residual has mean zero by
iterated expectation. Cross-fitted recursive estimators use this structure
to limit nuisance bias. Nested-function and longitudinal-rate analyses
explain the needed products of nuisance errors
\citep{nested2022,singh2022}. Sequential double robustness can hold when
an appropriate outcome or treatment component is correct at each stage
\citep{sequentialdr2018}; this property requires the corresponding recursive
construction and cannot be assigned to every algorithm labeled DR.

Longitudinal targeted estimation updates an initial outcome regression
along a fluctuation informed by treatment/censoring weights, aiming to
solve the influence-function estimating equation while respecting outcome
bounds. It is a useful alternative to a one-step estimator, especially
for bounded risks. The target, nuisance fits, targeting convergence and
support diagnostics must all be retained. A point estimate from a neural
sequence predictor does not supply these inferential properties.

Dynamic DML via g-estimation uses orthogonal residualization and recursive
removal of future effects in a structured dynamic model
\citep{dynamicdml2021}. Its independent-trajectory formulation suits some
many-customer panels. It does not directly supply a theorem for one
unrestricted aggregate bank series. The product must expose the chosen
effect parameterization, not relabel all dynamic models a single method.

\section{An exact two-decision example}
Set $L_0\sim\mathrm{Bernoulli}(1/2)$ and
\begin{align*}
 \Prob(A_0=1\mid L_0)&=0.2+0.6L_0,\\
 \Prob(L_1=1\mid L_0,A_0)&=0.1+0.5L_0+0.2A_0,\\
 \Prob(A_1=1\mid L_0,A_0,L_1)&=0.15+0.5L_1+0.1A_0,\\
 Y&=20+3A_0+4A_1+10L_1+6L_0.
\end{align*}
Every relevant action has positive probability. Under always versus never
treat, direct contributions add $3+4=7$ and the mediated change in $L_1$
adds $10(0.2)=2$. The total is $9$.
The intervention means are $35.5$ and $26.5$.

Enumerating all 16 observed cells verifies that the g-formula and exact
IPW reproduce both means. The observed always-adherer versus never-adherer
contrast is $\NumLongNaive$. A correctly specified conditional outcome
regression includes $L_1$ but its two action coefficients sum only to $7$;
that sum excludes the mediated component. The lesson concerns the causal
contrast, not an inability of regression to fit the observed outcome.

\begin{figure}[htbp]\centering
\includegraphics[width=\textwidth]{figures/effect_bias.pdf}
\caption{Synthetic static and longitudinal comparisons. Large apparent effects can reflect assignment and intermediate-variable conditioning.}
\label{fig:effects}
\end{figure}

\section{Feasible interventions and weak support}
An always-assist policy may be unsupported for low-risk customers or
ineligible accounts. Incremental interventions instead multiply treatment
odds by $\delta>0$:
\begin{equation}
 g_{\delta,t}(1\mid h)=
 \frac{\delta e_t(1\mid h)}{\delta e_t(1\mid h)+1-e_t(1\mid h)}.
 \label{eq:incremental}
\end{equation}
These preserve deterministic zero/one assignment regions
\citep{kennedy2019}. They answer how outcomes change under an odds shift,
not under universal treatment. Since the target itself depends on the
observed propensity law, its efficient score is not obtained by blindly
treating $g_\delta$ as externally fixed in Equation~\eqref{eq:longdr}.

Modified treatment policies can shift a rate or limit while respecting
conditional support \citep{lmtp2023}. A natural-value modification and
its stochastic analogue may require different sequential randomization
conditions. An economically feasible shift also needs a clear operational
implementation: a mathematical intervention on utilization is not necessarily
an action the bank can directly perform.

\section{Censoring, eligibility and competing events}
Account closure, debt sale, default, prepayment and loss of observation have
different meanings. If an outcome remains well-defined but becomes
unobserved, identification can require conditional independent censoring
and positive observation probability. Artificial censoring at deviation
from a treatment regime introduces an additional weighting or modeling task.
If treatment changes eligibility for a later action, simply restricting to
observed survivors changes the target population.

\citet{selective2025} explicitly model treatment-dependent eligibility;
their monotone eligibility setup must be checked against banking re-entry
and repeat episodes. \citet{competing2026} derive dynamic-regime cumulative
incidence methods with competing risks in a recent preprint. These motivate
separate outcome contracts for default and prepayment.

For disjoint terminal events with regime-specific marginal hazards
$\lambda^g_{D,k}$ and $\lambda^g_{P,k}$ among event-free accounts,
\begin{equation}
 F_D^g(K)=\sum_{k=1}^K
  \left[\prod_{j<k}(1-\lambda^g_{D,j}-\lambda^g_{P,j})\right]
  \lambda^g_{D,k}.
 \label{eq:cif}
\end{equation}
Estimating these regime hazards still requires longitudinal identification.
Censoring prepayment and calling the result the actual probability of default
would answer a different, potentially hypothetical question. A customer-value
outcome can incorporate both terminal events directly, provided cash-flow
ownership and observation are defined.

\section{Deep sequence models and the product boundary}
Counterfactual recurrent networks, causal transformers and neural
g-computation supply flexible representations or transition estimators
\citep{bica2020,melnychuk2022,gnet2020}. The inspected texts retain
consistency, sequential ignorability and overlap. Attention or adversarial
balance does not verify those assumptions.

They are appropriate challengers when long histories and nonlinear dynamics
defeat simpler models. Evaluate simulated interventions, misspecification,
support and policy value alongside factual prediction. Use simpler
g-computation, IPW and recursive DR comparators with the same target.
The first product should support a modest set of interpretable policies and
explicit horizons before claiming arbitrary long-horizon optimization.

\section{The estimating recursion and its error propagation}
For a cash-flow outcome, it is often more natural to keep rewards on their
actual dates than to aggregate them first into $Y$. Let $R_t$ be contribution
between decisions $t$ and $t+1$, let $h(H_T)$ settle the terminal account,
and fix discount factor $\gamma$. The action value satisfies
\[
 Q_t^g(h,a)=\E[R_t+\gamma V_{t+1}^g(H_{t+1})\mid h,a],
 \quad V_t^g(h)=\sum_a g_t(a\mid h)Q_t^g(h,a),
 \quad V_T^g=h.
\]
The terminal settlement must exclude cash already included in earlier
rewards. In lending, a remaining claim or expected recovery can belong
there; subtracting the same principal loss both monthly and at termination
would value a different account.

For fitted action values, define $\widehat V_t=\sum_a g_t\widehat Q_t$
and recursively set
\begin{equation}
 D_T=h,\qquad
 D_t=\widehat V_t(H_t)+\frac{g_t(A_t\mid H_t)}{\widehat e_t(A_t\mid H_t)}
       \{R_t+\gamma D_{t+1}-\widehat Q_t(H_t,A_t)\}.
 \label{eq:backward-dr}
\end{equation}
Expanding this recursion gives the full cumulative-ratio residual score.
It is the sequential forecast-and-correction construction of
\citet[Section 4.1]{jiang2016}, with the terminal account made explicit.
Every early ratio enters the correction for a later reward because early
actions determine which histories arrive at that later decision.

\begin{proposition}[Two routes to the mean of a fixed policy]
Under the stated sequential identification assumptions, with independent
evaluation trajectories, correct terminal settlement,
policy support and integrable terms, the score has mean $V^g$ if all
logging probabilities are correct. It also has mean $V^g$ if all action
values $\widehat Q_t=Q_t^g$ are correct and the fitted positive logging
probabilities are fixed history-measurable functions.
\end{proposition}
\begin{proof}
With correct ratios, each weighted dated quantity has its expectation under
the intervention law. Averaging $\widehat Q_t$ over the policy gives
$\widehat V_t$; the successive value terms telescope. With correct action
values, start at $T$ and use backward induction. The conditional expected
Bellman residual at each date is zero given $H_t,A_t$, so multiplying by
the history-measurable ratio leaves its mean zero.
\end{proof}

This is a mean identity. It does not state that every finite-sample normal
interval has nominal coverage, that a learned policy can be evaluated on
its own training outcomes, or that any estimated latent representation is
a valid state. The collections study will show visible undercoverage for
some unbiased or approximately unbiased constructions when path weights
are concentrated.

The propagation of nuisance error can be made explicit. Define
$b_t(h)=\E[D_t\mid H_t=h]-V_t^g(h)$ and
$\epsilon_t(h,a)=\widehat Q_t(h,a)-Q_t^g(h,a)$. Conditional on the
training fits,
\begin{align}
 b_t(h)&=\sum_a g_t(a\mid h)
       \left(1-\frac{e_t(a\mid h)}{\widehat e_t(a\mid h)}\right)
       \epsilon_t(h,a)\nonumber\\
 &\quad+\gamma\sum_a
   \frac{g_t(a\mid h)e_t(a\mid h)}{\widehat e_t(a\mid h)}
   \E[b_{t+1}(H_{t+1})\mid h,a],\qquad b_T=0.
 \label{eq:dynamic-remainder}
\end{align}
Substitution of $\E[D_{t+1}\mid H_{t+1}]=V_{t+1}^g+b_{t+1}$
into Equation~\eqref{eq:backward-dr} and collection of the regression
errors proves the identity. The first line is a local product of nuisance
errors; the second transports all later error backward. Very large early
ratios can amplify an otherwise modest continuation error. A small
one-step prediction loss is therefore an inadequate diagnostic for a
long-horizon causal value estimator.

\section{Learning the continuation rather than the whole trajectory}
Forward simulation of the g-formula needs a conditional distribution for
every covariate whose
future realization affects decisions or rewards. For a large mixed record
of balances, transactions, arrears and categorical eligibility, building a
coherent joint transition can be difficult. Recursive regression instead
learns only the conditional continuation needed for a fixed policy and
outcome. It can be much cheaper when the target is narrow. The price is
that a different policy, horizon or outcome can require a new sequence of
regressions.

A simulation model is more reusable for new policy scenarios but exposes
more mechanisms to misspecification. A reward model, balance evolution,
observation process and action eligibility must evolve together. Simulating
arrears independently of credit limits can generate impossible account
states. Neural g-computation offers flexible conditional laws
\citep{gnet2020}; flexibility does not remove the obligation to preserve
account identities and interventions. The product should compare simulated
policy values with recursive estimates when both are identified, not
judge the simulator only by factual one-step prediction.

For either approach, partition complete customer or household histories
before training. Within each training fold, generate future pseudo-outcomes
without using the held-out customer's outcomes. A naive implementation can
cross-fit the last regression while leaking the held-out outcome through
an earlier fitted continuation. The complete dependency graph of nuisance
fits must respect the fold boundary. Nested and longitudinal orthogonal
methods formalize rate and error-control requirements for such constructions
\citep{nested2022,singh2022}.

This distinction also affects hyperparameter tuning. Tuning on final policy
value can adapt to noise in the evaluation trajectories. Prediction-based
tuning within training folds preserves the declared evaluation, though it
may not minimize policy-value error. An outer evaluation of the entire
training and selection procedure is needed when the scientific claim is
about that procedure rather than one frozen policy.

\section{Targeted estimation and bounded outcomes}
For a binary terminal event, a fitted conditional probability should remain
between zero and one. A simple logistic fluctuation uses an initial
$\widehat Q_t^0$ and a one-dimensional parameter $\epsilon$:
\[
 \operatorname{logit}\widehat Q_t^\epsilon
   =\operatorname{logit}\widehat Q_t^0+\epsilon H_t^*,
\]
where the chosen covariate $H_t^*$ or equivalent observation weight is
constructed from the treatment and censoring mechanism for the target.
Fitting the fluctuation so the relevant weighted residual is zero aligns
the plug-in estimate with an influence-function equation. A longitudinal
procedure updates backward through the continuation regressions.

The displayed expression describes the mechanism, not a complete universal
TMLE algorithm. Different targets require different clever covariates,
weight conventions and recursion. The implementation must state whether
outcomes have been bounded or rescaled, what happens at fitted probabilities
zero and one, and how targeting convergence is checked. The relevant
sequential robustness analysis is estimator-specific
\citep{sequentialdr2018}. This is why a longitudinal adapter should retain
its native algorithm and diagnostics instead of wrapping every estimator
behind a function that returns only a scalar ``effect.''

Targeting can respect the parameter space while a one-step correction falls
outside it in a small sample. That is a useful finite-sample feature, but
does not repair practical positivity or unmeasured confounding. Conversely,
an out-of-range one-step estimate is a diagnostic of instability; silently
clipping it to a probability conceals the issue and changes inference.

\section{Observation loss has its own clock and model}
Let $C_{t+1}=1$ mean the next required record remains observed, conditional
on being observed through $t$. Write its conditional probability as
$c_t(H_t,A_t)$. If observation is conditionally independent of the relevant
future potential outcomes given that information, a change-of-measure
weight contains
\[
 W_t^{g,C}=\prod_{s=0}^{t}
 \frac{g_s(A_s\mid H_s)}{e_s(A_s\mid H_s)}
 \frac{C_{s+1}}{c_s(H_s,A_s)}.
\]
This formula separates two missing pieces: the action law and the reporting
law. Exact randomized treatment probabilities do not make a selectively
missing outcome ignorable. A customer whose debt is sold may have an
observed sale price, an unobserved later recovery and a changed economic
owner. The appropriate bank contribution can be observed at sale even
though a different lifetime recovery outcome is censored.

Administrative censoring at a fixed cutoff also changes maturity. A recent
cohort may have complete six-month outcomes but no twelve-month outcomes.
A twelve-month effect then needs older cohorts, a shorter target or an
extrapolation model. More recent transactions cannot fill calendar time
that has not elapsed. The data contract stores both the outcome end date
and the date by which the outcome could actually be known.

Competing events need the same discipline. Prepayment is not an arbitrary
missing default label. It changes the bank's remaining claim and may itself
be affected by assistance or pricing. Estimating default cumulative incidence,
event-free survival, a hypothetical no-prepayment risk and discounted bank
contribution are four different questions. Equation~\eqref{eq:cif} answers
the first under a regime; the cash-flow recursion can answer the fourth
without pretending that prepayment never occurs.

\section{Structural blips and dynamic effect parameters}
Sometimes the bank wants an interpretable effect parameter rather than a
complete regime value. Consider a two-action additive model in which the
conditional causal contribution of the final action is $\psi_1 A_1$.
Removing this contribution creates a blipped outcome
$Y^{\mathrm{blip}}=Y-\psi_1 A_1$. A g-estimation moment compares that
outcome across earlier actions after conditioning on their histories.
With an appropriate structural nested model, one estimates the final blip
first and moves backward, avoiding ordinary conditioning that blocks the
earlier mediated effect.

Dynamic DML combines this recursive removal with orthogonal nuisance
residualization \citep{dynamicdml2021}. The gains are a compact effect
parameterization and flexible adjustment; the costs are the restrictions
of the structural model, the need for support and dependence-compatible
sampling. Constant blips, heterogeneous blips and arbitrary dynamic regimes
are different model classes. The bank should choose among them according
to the decision and interpretability requirement rather than the apparent
modernity of the implementation.

The recommended comparison therefore has three tiers. An interpretable
structured model estimates a small set of economically meaningful dynamic
effects. Recursive DR or targeted estimation evaluates a prespecified set
of feasible regimes. A flexible transition or representation model explores
where those simpler constructions miss important history. Agreement is
useful evidence under shared assumptions; disagreement is a reason to
inspect targets, support and model restrictions, not to average incompatible
answers.

\chapter{Time series, dependence and computational scale}\label{ch:scale}
\section{Three different large-data regimes}
A bank may observe millions of short customer histories, a smaller panel of
branches exposed to common shocks, or one long aggregate balance-sheet
series. These regimes have different asymptotics and different independent
information. Pooling them into one large table does not make them independent.

For a stationary scalar score series with summable autocovariances, the
variance of its average is approximately
\begin{equation}
 \Var(\bar\phi_T)\simeq
 \frac{1}{T}\left\{\gamma_0+2\sum_{k\geq1}\gamma_k\right\},
 \qquad \gamma_k=\Cov(\phi_t,\phi_{t-k}).
 \label{eq:lrv}
\end{equation}
For an AR(1)-like correlation pattern, the variance-based illustrative
effective sample size is $T(1-\rho)/(1+\rho)$.
At $\rho=0.9$, it is about $T/19$. This calculation concerns a stationary
mean; it is not a general causal-information measure for nonstationary
bank data.

\section{A fold is an information boundary}
Customer-level cross-fitting keeps the complete customer's trajectory
outside its nuisance-training fold. Household or connected-business
interference can require a larger unit. A prospective evaluation also uses
training data available before its evaluation origin, with embargoes covering
feature lookback and overlapping outcome windows where needed.

Retrospective estimation under a justified stationary sampling design may
use different folds from a prospective deployment simulation. The product
must distinguish the two. Randomly splitting account-month rows leaks
individual histories. Pure time splitting can still share a common economic
shock between training and test entities. There is no universal split that
repairs every dependence structure.

Multiway-cluster DML splits along cluster dimensions under a specific
sampling framework \citep{cluster2021}. For crossed clustering, a familiar
score covariance construction adds cluster contributions and subtracts their
intersection, schematically
\[
 \widehat V=\widehat V_{\text{entity}}
            +\widehat V_{\text{date}}
            -\widehat V_{\text{entity,date}}.
\]
Scaling and finite-sample corrections depend on the estimator. Few clusters,
unbalanced groups and indefinite finite-sample estimates require disclosure
and suitable alternatives. Increasing the number of accounts cannot replace
missing independent policy dates.

HAC estimators, block bootstraps and cluster bootstraps must match the score
process and dependence assumptions. The report records bandwidth/block
length, cluster counts, sensitivity and any finite-sample correction.
Generic package defaults are not a scientific inference specification.

\section{Changing regimes and transport}
Rates, underwriting policies and customer behavior change. A pooled
coefficient across pre-pandemic, crisis and high-rate periods may have little
decision relevance. Stability diagnostics should examine treatment policy,
covariate support, residual distributions and effect heterogeneity by regime.
Change-point detection is evidence of distributional change, not a unique
explanation for it.

Transporting a conditional effect from source $S=1$ to target $S=0$
may use $\E[\tau(X)\mid S=0]$ if the relevant conditional response is
invariant and target covariates have source support. The target distribution
must be recorded. Unobserved effect modifiers, product versions and changes
in interference can invalidate transport even when feature distributions
overlap. Banking macro and equilibrium responses can change the mechanism
itself \citep{beraja2019,benetton2021}.

\section{Interference and shared constraints}
Customers can share households, firms can share owners, and offers can
compete for branch capacity. In such settings write
$Y_i(a_i,z_i)$, where $z_i=f_i(\bm a_{-i},\mathcal G)$ summarizes
others' assignments through a declared network $\mathcal G$.
The exposure mapping is a substantive restriction: it asserts which
features of others' treatment matter. Individual potential outcomes
indexed only by $a_i$ are inadequate under interference
\citep{hernan2020}.

Direct effects compare $Y_i(1,z)-Y_i(0,z)$ at fixed exposure;
spillover effects compare exposures at fixed own action.
A group policy effect changes both. Household or branch randomization,
or randomized treatment saturation across groups, can create the relevant
variation when operationally appropriate. Independent-unit inference must
then use the randomized or plausibly independent groups.

For observational networks, homophily, endogenous links and unobserved
shared shocks complicate identification. Neither a graph database nor
clustered standard errors eliminates these problems. The initial product
should support declared partial-interference groups and report when
cross-group spillovers remain plausible. General network identification
is a research extension, not a promised default feature.

\section{Compute around scientific boundaries}
The execution engine should materialize a versioned eligible population,
partition by the independent unit, compute fold-specific nuisances, and
aggregate scores with the declared cluster keys. Map tasks may fit independent
folds or bootstrap replicates; reductions retain sufficient statistics and
covariance blocks. Every retry uses the same input snapshot and random seed.
Hyperparameter search occurs inside nuisance-training partitions.

Dense data footprints are already material. At $10^8$ rows and 256
float32 features, the feature array alone is $102.4$ GB in decimal units,
before identifiers, outcomes, fold indicators, copies and model state.
Use columnar projection, predicate pushdown, partition pruning, streamed
score evaluation and sparse representations when justified. Avoid materializing
all lagged columns and all horizons for every exploratory study.

Finite-state filtering costs $O(NTK^2)$ for dense $K$-state transitions,
or less with sparse structure. The forward pass needs $O(K)$ working
state per active sequence; storing all smoothing messages costs $O(TK)$
per sequence unless checkpointing/recomputation is used. State dimension,
not just row count, can dominate. Particle methods trade computation for
approximation variance and degeneracy diagnostics.

Temporal graph search can be more difficult still. A coarse
$p^2L$ candidate count already gives three million directed lag candidates
for $p=500,L=12$, before conditioning-set search. Apply justified temporal
and domain constraints, variable grouping and staged screening. Screening
can discard true parents and must be evaluated as part of the discovery
procedure. Distributed execution does not eliminate the statistical
multiple-testing or combinatorial problem.

\section{A measurable scale qualification}
The benchmark program should use three tiers: $10^5$ rows for correctness,
$10^7$ for multi-team workflows, and $10^8$ for large group studies. These
are future test sizes, not achieved performance claims. Benchmark fixed
estimands and equivalent estimators on single-node and distributed execution.
Measure peak memory, wall time, compute cost, serialization, retry behavior
and numerical agreement. Report nuisance-fitting and data-preparation costs
separately.

Use an explicitly documented reference machine before stating latency
targets. A candidate target is interactive cohort diagnostics within 30
seconds on a precomputed study snapshot and scheduled large-study completion
within one working day. These are design goals contingent on workload and
hardware. They must not be advertised as measured service levels.
Scientific stopping rules take precedence: a fast job with unsupported
actions or invalid uncertainty has not completed an acceptable study.

\section{Dynamic responses: local projections and structural VARs}
A treasury question often concerns a path: how does a funding shock today
change deposits, lending and contribution over the next twelve months?
Write a stable reduced-form VAR as
\begin{equation}
 X_t=\sum_{\ell=1}^{L}A_\ell X_{t-\ell}+u_t,
 \qquad u_t=C\varepsilon_t,\quad
 \E[\varepsilon_t\varepsilon_t']=I.
 \label{eq:var}
\end{equation}
The lag coefficients describe conditional dynamics. The structural impact
matrix $C$ maps economically interpretable shocks to reduced-form
innovations. Estimating $A_\ell$ and $\Sigma_u=CC'$ does not in general
identify $C$. For any orthogonal $Q$, $(CQ)(CQ)'=\Sigma_u$.
An identifying restriction must choose among observationally compatible
shock decompositions or return a set of responses.

For a VAR(1), the response at horizon $h$ to shock $j$ is
$A^h C e_j$. A recursive timing restriction selects a triangular factor
of $\Sigma_u$. Its interpretation is a statement about which variables
can respond contemporaneously. Ordering funding before lending may be
plausible at one frequency and implausible at another. The selected order
must be attached to the impulse-response report as an assumption, not
hidden in a default Cholesky routine.

Local projections estimate each horizon directly, for example
\begin{equation}
 Y_{t+h}=\alpha_h+\beta_h D_t+\Gamma_h'W_t+v_{t+h}^{(h)},
 \label{eq:lp}
\end{equation}
where $D_t$ is the shock or endogenous intervention variable and $W_t$
contains the declared controls. Without an identification argument,
$\beta_h$ is a projection coefficient. When $D_t$ represents a suitably
identified shock, it can estimate a dynamic causal response. Direct horizon
estimation permits flexible responses and avoids imposing all finite-order
VAR restrictions, often with a variance cost. The review by
\citet{jorda2024} explains both the connection and the finite-sample
tradeoffs; it does not justify a universal preference for one method.

\subsection{External instruments identify a particular shock}
An external instrument $Z_t$ must correlate with the shock of interest and
be orthogonal to the other shocks that enter the dynamic outcome.
Under the moving-average representation
$X_t=\sum_{h\geq0}\Theta_h\varepsilon_{t-h}$, assume
$\E[Z_t\varepsilon_{1t}]=\alpha\ne0$ and
$\E[Z_t\varepsilon_{j,t+k}]=0$ for other contemporaneous shocks
and for the relevant nonzero leads and lags. Normalize the impact of shock
one on $X_{1t}$ to unity. Then
\begin{equation}
 \Theta_{h,i1}=\frac{\E[Z_t X_{i,t+h}]}{\E[Z_t X_{1t}]}.
 \label{eq:lpiv}
\end{equation}
The numerator equals $\alpha\Theta_{h,i1}$ because all other shock
covariances vanish; the denominator equals $\alpha$ under the normalization.
This is the dynamic covariance-ratio argument in
\citet[Section 2.3]{stock2018}. Controls require the corresponding
partialled-out restrictions.

The lead-lag exclusion is substantive. A monetary-policy announcement
surprise might also reveal information about future economic conditions.
A bank funding instrument may reflect a local demand shock. A strong first
stage does not exclude either channel. Direct LP-IV and a proxy SVAR have
different information requirements: the latter normally needs a suitable
invertibility or recoverability argument, while the direct approach in the
stated source can avoid invertibility under stronger lead-lag exclusions.
The product must retain that distinction when offering both estimators.

An intervention that raises a bank's posted rate for three months is also
different from a one-time structural shock with endogenous future rate
responses. An impulse response describes the latter under a maintained
system. A sustained policy path requires a different intervention or
structural policy simulation. Calling both a ``25-basis-point scenario''
would conceal an economically material difference.

\section{Inference concerns the score process}
For a scalar moment estimator satisfying
$T^{-1}\sum_t\psi_t(\widehat\theta)=0$, a first-order expansion gives
\[
 \sqrt T(\widehat\theta-\theta_0)
 \simeq -J^{-1}T^{-1/2}\sum_t\psi_t(\theta_0),
 \qquad J=\E[\partial_\theta\psi_t(\theta_0)].
\]
Under an appropriate dependent-data central limit theorem, its covariance
uses the long-run variance of $\psi_t$, not simply that of outcomes or
regression residuals. With matrix moments the same expression uses
$J^{-1}\Omega J^{-\prime}$. A kernel HAC estimate can be written
\[
 \widehat\Omega=\widehat\Gamma_0+
 \sum_{k=1}^{b}K(k/b)(\widehat\Gamma_k+\widehat\Gamma_k'),
 \quad \widehat\Gamma_k=T^{-1}\sum_{t=k+1}^{T}
     \widehat\psi_t\widehat\psi_{t-k}'.
\]
Bandwidth $b$, kernel, centering, finite-sample adjustment and the range of
available $T$ determine the approximation. Short series and persistent
scores can yield poor finite-sample coverage even with a familiar HAC label.

Overlapping horizons make the distinction useful. In an AR(1), repeated
substitution gives
$X_{t+h}=\rho^h X_t+\sum_{j=1}^{h}\rho^{h-j}u_{t+j}$.
The projection residuals overlap and are serially correlated. Some
lag-augmented local-projection constructions make the relevant score
serially uncorrelated under stronger innovation conditions, allowing
heteroskedasticity-robust inference. That conclusion comes from the
regressor--innovation structure, not from denying residual overlap.
\citet[Section 6]{jorda2024} distinguishes these cases. Its reviewed
uniform-inference result should not be transferred without proof to arbitrary
multiway bank panels or learned nonlinear nuisance scores.

Pointwise intervals for twelve horizons also do not give a simultaneous
95\% band for the response path. A maximum-statistic bootstrap, a justified
joint Gaussian approximation or a conservative multiplicity adjustment
can answer the pathwise question. The resampling must preserve the time
and cluster structure. A funding decision depending on the discounted sum
of the response needs the covariance across horizons, which can be much
more informative than reading individual significance stars.

\section{A repeated-sample demonstration of common shocks}
The supplied experiment draws a score panel
\[
 \phi_{it}=a_i+f_t+\epsilon_{it},\qquad
 a_i,f_t,\epsilon_{it}\ \hbox{independent }N(0,1),
\]
with 96 entities and 36 dates. The target mean is zero. Direct variance
calculation yields
\begin{equation}
 \Var(\bar\phi)=\frac1{96}+\frac1{36}+\frac1{96\cdot36}.
 \label{eq:crossed-variance}
\end{equation}
The resulting true standard error is approximately $0.1962$.
There are 3,456 rows, but shared entity and date shocks dominate the
uncertainty. Treating rows as independent counts those shocks repeatedly.

\begin{table}[htbp]\centering\small
\begin{tabular}{@{}lrr@{}}
\toprule
Variance method & Mean SE & 95\% coverage\\
\midrule
Rows independent & 0.029 & 0.227\\
Entity clusters & 0.103 & 0.693\\
Entity and date clusters & 0.193 & 0.944\\
\bottomrule
\end{tabular}

\caption{Coverage over 1,200 independently generated panels. Critical values
are normal; the crossed construction uses entity and date corrections. This
is a stated Gaussian score design, not a theorem for arbitrary bank panels.}
\label{tab:dependence-coverage}
\end{table}

Entity clustering protects against repeated observations of a customer but
misses the date component in this design. The crossed construction is much
closer to the intended coverage, while only 36 date clusters still make
finite-sample approximation relevant. The code retains the Monte Carlo
standard error of every reported coverage estimate. Adding a million
customers would reduce the first and third terms in
Equation~\eqref{eq:crossed-variance}; it would leave the date term.
This is why the execution planner must preserve the scientific unit rather
than simply distribute all rows uniformly over workers.

\begin{figure}[htbp]\centering
\includegraphics[width=\textwidth]{derivations/dependence_and_irf.pdf}
\caption{Left: dependence changes uncertainty in the repeated score experiment.
Right: two structural timing assumptions produce different responses while
preserving the same reduced-form Gaussian VAR law.}
\label{fig:dependence-irf}
\end{figure}

\section{Large scale has several different bottlenecks}
The scientific workload can be decomposed as
\[
 T_{\rm job}=T_{\rm snapshot}+T_{\rm design}+T_{\rm nuisance}
            +T_{\rm score}+T_{\rm inference}+T_{\rm export},
\]
with separate memory and storage accounts. A fast score reduction can be
irrelevant when source reconstruction or repeated nuisance fitting dominates.
Likewise, a GPU can accelerate dense model fitting while increasing the
cost of copying sparse event data into device memory. The product should
measure each stage and retain the workload configuration with the result.

Consider $N$ independent trajectories, horizon $T$, $K$ latent states,
$p$ observed variables, $L$ candidate lags, $F$ folds and $B$ bootstrap
replicates. These dimensions describe different jobs. Dense HMM smoothing
requires order $NTK^2$ arithmetic per likelihood pass. Refitting all
parameters in every bootstrap can multiply that work by $B$, not merely
add a cheap variance reduction. A temporal graph learner may have
$p^2L$ candidate edges and far more conditioning tests. A flexible
longitudinal nuisance pipeline may fit a distinct learner for each date,
fold and target policy. Reporting only total rows misses each of these
sources of difficulty.

Several computational reductions have scientific conditions. Homogeneous
transition counts can be summed across customers only when the same
parameterization applies. State compression must preserve the decision-
relevant conditional law. Pooling across jurisdictions can improve precision
only under a defensible sharing or hierarchical model. Sparse graph
screening reduces search only by making a restriction on candidate parents.
Each reduction belongs in the model specification and sensitivity study.

The recommended architecture caches reusable source snapshots and
fold-consistent nuisance predictions, streams score evaluation, and reduces
cluster contributions without materializing unnecessary dense copies.
It parallelizes independent fits or sequences before splitting a dependent
sequence. For long-sequence smoothing, checkpointing trades recomputation
for storage; truncating a sequence instead changes the conditional
information available to the estimator. Those are different operations.

\section{A scale result must preserve the statistical result}
The construction chapter supplies a measured streaming reduction workload
and a temporal-revision implementation. These bounded measurements establish
the operation and resource use of those components. They do not establish
end-to-end performance for a causal forest, temporal graph search or
controlled-state likelihood on a banking group's records.

A future large-scale qualification uses identical source snapshots, folds,
target weights and random seeds for a single-node reference and a distributed
implementation. It compares estimates, influence contributions, covariance,
unsupported-cell flags and financial totals, not just runtime. Differences
caused by floating-point summation are bounded; differences caused by changed
folds or dropped records are scientific defects. Retry tests must show that
a failed worker does not duplicate a customer's contribution.

Resource budgets should be expressed per workload: bytes scanned per study,
peak worker memory, total CPU/GPU time, shuffled bytes, model checkpoints,
and the number of complete refits. A claim such as ``supports one hundred
million rows'' becomes useful only after specifying these quantities, the
estimator and the attainable uncertainty. Computation can make a valid
analysis feasible; it cannot create additional independent policy shocks.

\chapter{Partial observation, estimation and policy value}\label{ch:pomdp}
\section{A precise event clock}
Repeated bank decisions may depend on unobserved financial resilience,
liquidity needs or customer intent. A POMDP represents a latent state $S_t$,
an observation $O_t$ available before action $A_t$, a reward $R_t$, and
the next state and observation. Let
$H_t=(O_0,A_0,R_0,\ldots,O_t)$ be the recorded decision history.
The policy available to the bank must be measurable with respect to this
history, or to an explicitly enlarged logged information set.

For clarity, suppose a finite-state model has initial law $\mu_\theta$,
initial emission $B_\theta(o_0\mid s_0)$, transition
$P_\theta(s'\mid s,a)$, reward density $F_\theta(r\mid s,a,s')$
and next emission $B_\theta(o'\mid s',a)$. These factorization choices
are modeling assumptions. A joint emission for rewards and signals may be
needed when conditional independence is inappropriate.

Let the behavior policy be $q_\eta(a\mid s,h)$, potentially using private
information summarized by $s$. The observed-data likelihood for one
trajectory is
\begin{align}
 L(\theta,\eta)=\sum_{s_{0:T}}&
 \mu_\theta(s_0)B_\theta(o_0\mid s_0)\nonumber\\
 {}\times\prod_{t=0}^{T-1}&
 q_\eta(a_t\mid s_t,h_t)
 P_\theta(s_{t+1}\mid s_t,a_t)
 F_\theta(r_t\mid s_t,a_t,s_{t+1})
 B_\theta(o_{t+1}\mid s_{t+1},a_t).
 \label{eq:pomdplikelihood}
\end{align}
Omitting the behavior factor is justified for estimating environment
parameters when action assignment depends only on the recorded history
and is parameter-distinct: it can then be factored out of the latent-state
sum in the joint likelihood. If the logger uses unobserved information,
the action itself is evidence about latent state and the factor cannot
generally be removed. Fitting a model conditional on the action sequence
without this distinction can misinterpret selection as a transition effect.

\section{Filtering, smoothing and estimation}
The belief-state construction connects partial observation to planning
\citep{kaelbling1998}; estimation additionally requires an explicit likelihood.
For an observed-history policy, after observing $O_t$ and previous rewards,
let $b_t(s)=\Prob_\theta(S_t=s\mid H_t)$.
The action is not an additional latent-state signal conditional on $H_t$.
After observing $r_t,o_{t+1}$,
\begin{equation}
 b_{t+1}(s')=
 \frac{B_\theta(o_{t+1}\mid s',a_t)
 \sum_s b_t(s)P_\theta(s'\mid s,a_t)F_\theta(r_t\mid s,a_t,s')}
 {\sum_{u}B_\theta(o_{t+1}\mid u,a_t)
 \sum_s b_t(s)P_\theta(u\mid s,a_t)F_\theta(r_t\mid s,a_t,u)}.
 \label{eq:belief}
\end{equation}
If the logger depends on latent state, first update
$b_t^+(s)\propto q_\eta(a_t\mid s,H_t)b_t(s)$ and use $b_t^+$ in
the observation update. Under a new intervention policy based only on
recorded history, this observational action update is removed.

Normalizers from the forward recursion accumulate the log likelihood,
avoiding numerical underflow. Backward messages produce smoothed state
probabilities and expected transition counts. In an EM algorithm,
\[
 \widehat P_a(s,s')=
 \frac{\sum_{i,t:A_{it}=a}\xi_{it}(s,s')}
      {\sum_{i,t:A_{it}=a}\sum_u\xi_{it}(s,u)}
\]
updates tabular transitions, with explicit prior/regularization when
counts are small. Emission and reward parameters are updated consistently
with their likelihood. Unobserved-logger parameters, if included, require
their own model and identification restrictions.

Alternatives include direct differentiation of the marginalized likelihood,
Bayesian inference, variational approximations and particle methods for
nonlinear continuous states. Compare multiple initializations, held-out
predictive likelihood, residual dynamics, posterior predictive checks,
parameter recovery and prior sensitivity. Label permutations require a
normalization or permutation-invariant comparison. Numerical convergence
and good predictive fit do not establish causal identification.

\section{Statistical and causal identification differ}
\citet{hmmident2016} establish HMM identification under full-rank ergodic
transitions, distinct emissions and other stated conditions, up to state
labels. This is a result about an observed distribution within an HMM
class. A controlled bank model additionally needs sufficiently informative
actions, identifiable reward/observation mechanisms and an intervention
interpretation. The stationary uncontrolled theorem does not supply all
of these conditions.

\begin{proposition}[Observed fit need not identify an action effect]
There exist two binary-action binary-outcome models with the same strictly
positive observed law and different average treatment effects.
\end{proposition}
\begin{proof}
In model I, $U\sim\mathrm{Bernoulli}(1/2)$,
$\Prob(A=1\mid U)=0.2+0.6U$ and
$\Prob(Y=1\mid U,A)=0.2+0.6U$. Action has no causal effect.
In model II, $A\sim\mathrm{Bernoulli}(1/2)$ independently and
$\Prob(Y=1\mid A)=0.32+0.36A$. Its effect is $0.36$.
Both imply
$\Prob(A,Y)=(0.34,0.16,0.16,0.34)$ in cell order
$(0,0),(0,1),(1,0),(1,1)$. Their observational likelihoods agree
while their intervention distributions differ.
\end{proof}

This one-step counterexample embeds in a sequential problem. Observed
action positivity and perfect predictive fit therefore cannot guarantee
causal policy identification. The practical remedy is an explicit route:
randomized or measured-history assignment, defensible structural restrictions,
valid instruments/proxies, or partial identification.

\section{Proxies and instruments are measurement designs}
In a static proximal model let $U$ be hidden confounding, $Z$ a
treatment-inducing proxy, $W$ an outcome-inducing proxy and $X$ observed
covariates. One bridge equation is
\begin{equation}
 \E[Y\mid Z,A,X]=\E[h(W,A,X)\mid Z,A,X].
 \label{eq:bridge}
\end{equation}
Under latent exchangeability, proxy exclusions such as
$Y\perp Z\mid A,U,X$ and $W\perp(A,Z)\mid U,X$, completeness,
bridge existence and support, a causal mean can be expressed as
$\E[h(W,a,X)]$ \citep{miao2018,proximal2020}.
The exclusions describe how variables are generated. Two highly predictive
features are not necessarily two valid proxies.

The bridge is an inverse problem. Weak proxies can make it ill-conditioned;
regularization trades variance against bias. Minimax methods and regression
bridges provide estimation strategies under corresponding restrictions
\citep{ghassami2022,proxreg2024}. The product should show proxy-role
justification, bridge residuals, conditioning diagnostics and regularization
sensitivity; these do not empirically prove completeness.

POMDP/proximal OPE theories impose particular information and proxy
structures \citep{tennenholtz2020,bennett2021,bennettprox2023}.
IV-based sequential evaluation also requires relevance, exclusion,
interaction restrictions and dynamic conditional independences
\citep{xu2023}. These are research-grade options for a bank with the
right measurement design, not generic repairs for absent action logs.
Recent hidden-treatment work similarly requires informative surrogates
and additional proxy restrictions \citep{hidden2026}.

\section{Policy evaluation before optimization}
The actor must also be clear. A POMDP policy can describe a bank choosing
actions given customer-state uncertainty. A customer's refinancing or
account-closing choice is a different behavior model, potentially with
private preferences, switching costs and beliefs. Estimating that choice
model requires restrictions on utility, expectations and state transitions.
Optimizing the bank's supplied reward function does not identify customer
preferences. Structural demand or dynamic choice components can supply
this behavior only with their own identifying evidence.

With sufficient observed histories and unconfounded logging, a target
policy $\pi$ has finite-horizon value
\[
 V^\pi=\E^\pi\!\left[\sum_{t=0}^{T-1}\gamma^tR_t\right].
\]
Importance ratios $\rho_t=\pi(A_t\mid H_t)/e(A_t\mid H_t)$
and their products correct behavior-policy selection when the relevant
conditions hold. A sequential DR value score can combine action-value
regression with weighted Bellman residuals. For a Markov state $X_t$,
one familiar finite-horizon form is
\begin{equation}
 \widehat V^\pi=
 \frac1n\sum_i\left[
 \widehat V_0(X_{i0})+
 \sum_{t=0}^{T-1}\gamma^tW_{it}
 \{R_{it}+\gamma\widehat V_{t+1}(X_{i,t+1})
                   -\widehat Q_t(X_{it},A_{it})\}\right],
 \label{eq:opedr}
\end{equation}
where $\widehat V_t(x)=\sum_a\pi_t(a\mid x)\widehat Q_t(x,a)$,
$\widehat V_T=0$, and $W_{it}=\prod_{s\leq t}\rho_{is}$.
Markov restrictions permit more efficient marginalized-ratio approaches
\citep{kallus2020}; replacing $X_t$ with an estimated belief does not
automatically inherit their identification or variance theory.

For two fixed policies, estimate the difference with paired scores and
preserve their covariance. A candidate deployment rule is a positive
lower confidence bound on incremental net value, with customer-outcome
constraints and correction for policy selection. For learned policies,
use held-out trajectories or valid uniform inference; ordinary pointwise
intervals after a large policy search are optimistic
\citep{policy2021}.

An ambiguity set $\mathcal M$ gives a conservative criterion
$\inf_{M\in\mathcal M}\{V_M^\pi-V_M^{\pi_0}\}$.
Its meaning depends on why $\mathcal M$ contains the relevant model.
Sharp dynamic bounds in \citet{sharpdynamic2026} assume stage-local
confounding; persistent customer fragility can violate the rectangular
structure required for local backward optimization. A recent negative result
constructs exponentially difficult OPE despite favorable behavior-marginal
coverage under history-dependent logging \citep{opehard2026}. It cautions
against diagnosing difficulty with only one-step action counts.

\section{A reproducible estimation and filtering illustration}
The supplied two-state model has initial law $(0.7,0.3)$, emissions
\[
 B=\begin{pmatrix}.8&.2\\.2&.8\end{pmatrix},\quad
 P_0=\begin{pmatrix}.9&.1\\.25&.75\end{pmatrix},\quad
 P_1=\begin{pmatrix}.95&.05\\.55&.45\end{pmatrix}.
\]
For actions $(1,0,1)$ and observations $(1,0,1,1)$, forward
likelihood $0.0234985$ equals explicit enumeration of all 16 latent
paths. Rewards are omitted from filtering in this illustration.
Expected immediate reward vectors are $(12,-8)$ for no assistance and
$(10,1)$ for assistance. With fragile-state belief $b$, assistance
is myopically preferable when $-2+11b>0$, or $b>2/11$.
This threshold is not the long-horizon optimal policy.

\begin{figure}[htbp]\centering
\includegraphics[width=\textwidth]{figures/belief_value.pdf}
\caption{Exact filtering and a myopic value comparison in a supplied causal model.}
\label{fig:belief}
\end{figure}

A separate simulation generates 700 trajectories of 14 randomized actions.
With emissions and initial law fixed and known, likelihood optimization
recovers the four transition-to-fragile probabilities as approximately
$(0.0935,0.7002,0.0774,0.4192)$ against truth
$(0.1,0.75,0.05,0.45)$. Three starts converge to closely agreeing
objectives; maximum absolute error is $\NumHMMError$.
This verifies a small controlled estimation construction. It does not
establish unknown-emission identification, uncertainty coverage, real-bank
state recovery or bank policy value.

Counterfactual trajectories introduce another distinction.
A Gumbel-max structural model specifies an exogenous-noise coupling across
actions \citep{oberst2019}. Different couplings can share transition
probabilities while producing different individual counterfactual paths.
The product should label such paths as model-conditional scenarios.

\chapter{Constructing a controlled latent-state estimator}\label{ch:latentestimation}
\section{The estimation problem the product must own}
A planning library generally expects transition, observation and reward
models. A bank has transactions, incomplete measurements and a historical
decision process. Bridging those objects is a substantive product capability.
The estimation module must specify a latent-state model, reconstruct its
event clock, fit the observed-data likelihood, quantify uncertainty and
state the conditions under which changing the action mechanism has a
causal interpretation. Passing a fitted transition matrix to a
POMDP solver leaves most of that work unresolved.

For $K$ states, $M$ signal categories and a finite action set, the model
has $K-1+|\mathcal A|K(K-1)+K(M-1)$ free probabilities before
reward, covariate and logging parameters. Small $K$ makes exact
marginalization possible and is appropriate for a first implementation.
Increasing $K$ creates flexibility, weakly separated states and more local
optima; state count is a model-selection question, not a quality score.
A state name such as ``fragile liquidity'' is an economic hypothesis.
Ordering states by expected delinquency fixes a label convention; it does
not establish a unique financial mechanism.

\section{Forward and backward messages}
Suppress a separate reward emission and define the observed transition kernel
\[
 M_t(s,s')=q_\eta(a_t\mid s,h_t)
          P_\theta(s'\mid s,a_t)B_\theta(o_{t+1}\mid s',a_t).
\]
When logging uses only observed history and its parameters are distinct,
its factor is constant across latent paths and can be omitted from the
objective for environment parameters. This is a joint-likelihood
factorization argument. It does not assert that adaptively chosen actions
are independent of observations, or that the remaining factor is the
ordinary conditional density of observations given the entire future
action sequence.

Let $\alpha_t$ be the normalized filter after the current signal and
preceding observations and actions. Start with
$c_0=\sum_s\mu(s)B(o_0\mid s)$ and
$\alpha_0(s)=\mu(s)B(o_0\mid s)/c_0$. Then
\begin{align}
 \widetilde\alpha_{t+1}(s')&=\sum_s\alpha_t(s)M_t(s,s'),&
 c_{t+1}&=\sum_{s'}\widetilde\alpha_{t+1}(s'),\\
 \alpha_{t+1}(s')&=\widetilde\alpha_{t+1}(s')/c_{t+1},&
 \log L&=\sum_{t=0}^{T}\log c_t.
 \label{eq:scaled-forward}
\end{align}
Consistently scaled backward messages satisfy $\beta_T(s)=1$ and
$\beta_t(s)=\sum_{s'}M_t(s,s')\beta_{t+1}(s')/c_{t+1}$.
The smoothed transition and state probabilities are
\begin{equation}
 \xi_t(s,s')=
 \frac{\alpha_t(s)M_t(s,s')\beta_{t+1}(s')}{c_{t+1}},
 \qquad \gamma_t(s)=\sum_{s'}\xi_t(s,s'),
 \label{eq:controlled-smoothing}
\end{equation}
with $\gamma_T=\alpha_T$. The identities can be checked by direct
enumeration on a short sequence. They extend the expected-count
construction described by \citet[pp.264--266]{rabiner1989} to the
declared controlled kernel. A separate reward emission multiplies $M_t$
when rewards convey additional state information.

The distinction between filtering and smoothing matters operationally.
Filtering uses information available at the decision. Smoothing uses later
observations and actions to revise the historical state estimate. A plot
of smoothed fragility may help diagnose a fitted model, but using that
series as a feature in a retrospective ``prospective'' policy evaluation
would leak future information.

\section{The M step is a constrained count calculation}
The expected complete log likelihood contains
\[
 \sum_s N_s^0\log\mu_s+
 \sum_{a,s,s'}N^P_{ass'}\log P_a(s,s')+
 \sum_{s,o}N^B_{so}\log B_s(o)+
 \sum_{s,a}N^q_{sa}\log q_s(a).
\]
The counts sum $\gamma$ and $\xi$ over trajectories with the
appropriate observed action or signal. A Lagrange multiplier shows that
maximizing $\sum_jN_j\log p_j$ subject to $\sum_jp_j=1$ gives
$p_j=N_j/\sum_kN_k$. Row-wise application updates the initial law,
transitions, emissions and, when modeled, logging probabilities.
Structural zeros must follow from the model; small positive pseudocounts
are disclosed regularization.

The E step allocates each trajectory probabilistically among states.
The M step updates the specified likelihood from those expected counts.
A Viterbi path instead chooses a single most probable state sequence.
Treating that path as observed data discards state uncertainty and generally
does not maximize the same likelihood. A convenient visualization must
not silently replace the estimator.

EM need not reach a global maximum. Multiple starts, likelihood increments,
iteration caps, boundary probabilities and state separation belong in
the fit receipt. Direct differentiation of the scaled marginalized
likelihood is a useful challenger. Bayesian inference adds priors and
posterior computation; variational estimation introduces approximation
error. None supplies causal identification through numerical convergence.

\section{Unknown emissions and hidden logging: a fitted study}
The study generates 800 independent trajectories with twelve actions,
two latent states and three signal categories:
\[
 \mu=(.72,.28),\qquad
 P_0=\begin{pmatrix}.93&.07\\.22&.78\end{pmatrix},\quad
 P_1=\begin{pmatrix}.94&.06\\.52&.48\end{pmatrix},\quad
 B=\begin{pmatrix}.78&.18&.04\\.16&.34&.50\end{pmatrix}.
\]
The learner receives only signals and actions. Four starts estimate the
initial law, transitions and emissions; labels are aligned by expected
signal severity before parameter errors are evaluated. The generator's
latent states never enter estimation.

Randomized logging chooses action one with probability $0.4$. In the
observational design its probabilities are $0.15$ and $0.80$ in the two
hidden states. The same observational trajectories are fitted with and
without a state-dependent logging factor. The joint fit estimates that
factor as well. The experiment stipulates the finite model and invariance
of transitions and observations under an action intervention. Those
restrictions make the generated environment causal; the estimation code
does not discover their truth.

\begin{table}[htbp]\centering\small
\begin{tabular}{@{}lrrr@{}}
\toprule
Estimation design & Max $P$ error & Max $B$ error & Policy gain (USD)\\
\midrule
Randomized & 0.026 & 0.054 & 24.90\\
Hidden logging, omitted action law & 0.384 & 0.066 & -22.40\\
Hidden logging, joint action law & 0.045 & 0.007 & 13.61\\
\bottomrule
\end{tabular}

\caption{Unknown-emission estimation. The policy contrast is assistance
after the most severe current signal versus no assistance, over eight
decisions. The true gain is USD\,22.34.}
\label{tab:hmm-estimation}
\end{table}

The omitted-action model reverses the sign of the policy contrast in this
sample. Historical assistance reveals fragility; ignoring that selection
can interpret a selected post-assistance distribution as a harmful action
effect. The joint model corrects this particular mechanism, but its
single-sample value remains materially different from truth. Raw objective
values from the two fitting routes are not directly comparable because
the joint objective includes the action density.

\begin{table}[htbp]\centering\small
\begin{tabular}{@{}lrrr@{}}
\toprule
Action model & Gain bias (USD) & Gain RMSE & Transition RMSE\\
\midrule
Omitted & -38.95 & 39.08 & 0.203\\
Joint & 3.41 & 21.90 & 0.051\\
\bottomrule
\end{tabular}

\caption{Thirty independently generated observational datasets of 500
trajectories, with both fitting routes applied to each dataset. Correct
logging specification improves parameter recovery but leaves substantial
policy-value variability.}
\label{tab:hmm-repeated}
\end{table}

The omitted model's mean policy-gain error is approximately
USD\,$-38.95$, with Monte Carlo standard error USD\,0.59. The joint
model's mean error is USD\,3.41 with Monte Carlo standard error USD\,4.02,
and RMSE USD\,21.90. These results expose a failure and a conditional
remedy. They do not show that every controlled hidden-state model is
identified or that the remaining uncertainty permits a bank decision.
The runs retain finite iteration budgets. All 30 selected omitted-action
fits and seven of the 30 joint fits hit the 120-iteration repetition cap.
The repeated table therefore compares the implemented finite-budget
procedures, not fully optimized population estimators. Its error contains
optimization as well as sampling and specification effects.

A separate same-data convergence check extends the main observational
fits to a maximum of 2,000 iterations without changing data or starting
seeds. The omitted-action fit meets the stopping criterion at 163 iterations
for its best start and still gives a policy gain near negative USD\,22.40;
the joint fit again gives positive USD\,13.61. Thus the main sign reversal
persists after the stated numerical stopping test. This does not establish
a global optimum or isolate all specification bias in the repeated study.

\section{Uncertainty includes the observation model}
Seventy-nine trajectory bootstrap resamples of the randomized dataset
refit initial probabilities, transitions and emissions, holding the two
policies fixed. The percentile interval for their value difference is
approximately USD\,$[14.61,38.91]$. The bootstrap count gives coarse
tail quantiles; one selected bootstrap fit hits its 120-iteration cap.
This is a model-conditional interval for independent
trajectories, not a coverage study or an interval for unrestricted hidden
confounding in bank records.

Additional uncertainty arises from the state specification, intervention
invariance and policy selection. A bootstrap holding those choices fixed
does not address them. A fitted three-state model may describe historical
observations as well as a two-state model and value an untested policy
differently. An assistance action may alter reporting intensity as well
as liquidity. A policy selected from hundreds of candidates may exploit
noise. Each requires a separate comparison, sensitivity analysis or outer
evaluation design.

The product should show state separation, transition support, likelihood
profiles or posterior dependence, predictive residuals by regime, and
financial decision sensitivity. A large parameter interval need not imply
a large value interval, and a small parameter error can reverse a close
policy comparison. The decision determines which uncertainty matters.

\section{Planning with a distribution over states}
For a correctly specified environment and an actor restricted to recorded
history, let $b$ be its state belief, $p(o'\mid b,a)$ the next-signal
law and $\tau(b,a,o')$ its Bayes update. With feasible actions
$\mathcal A(b)$ and expected reward $r(b,a)$,
\begin{equation}
 V_t(b)=\max_{a\in\mathcal A(b)}\left[
 r(b,a)+\gamma\sum_{o'}p(o'\mid b,a)
       V_{t+1}\{\tau(b,a,o')\}\right].
 \label{eq:belief-bellman}
\end{equation}
Observed eligibility must be part of the information state if it changes
the action set. The continuation averages future state uncertainty and
future information. A myopic threshold neglects both continuation rewards
and any value of information \citep{kaelbling1998}.

The study values next-signal contribution at $(18,0,-90)$ dollars and
charges eight dollars for assistance. Reward is a deterministic function
of the next signal and action, so no separate reward information is
omitted from the filter. A 501-point belief grid solves eight-decision
control. Doubling resolution changes initial value by less than
USD\,0.0001. Under the supplied true model, values are approximately
$-21.92$ for no assistance, $0.42$ for assistance after a severe signal,
and $5.30$ for the grid policy. These numerical comparisons do not show
that the extra optimized value is statistically resolved by the fitted data.

\begin{figure}[htbp]\centering
\includegraphics[width=.9\textwidth]{derivations/pomdp_value.pdf}
\caption{Eight-decision value and initial action on a two-state belief grid,
conditional on the supplied causal environment.}
\label{fig:pomdp-plan}
\end{figure}

For bounded rewards and $\gamma<1$, the infinite-horizon Bellman
operator is a contraction of modulus $\gamma$: expectation changes by
at most the supremum change in continuation; discounting multiplies the
bound by $\gamma$; taking a maximum preserves it. This establishes a
unique fixed point for the supplied belief-MDP. It is a control result,
not an identification theorem for observational banking data.

\section{Partial observation is not always hidden confounding}
Under randomized or recorded-history assignment, partial observation can
make modeling difficult without confounding the action. History-based
importance evaluation remains identified under the sequential assumptions,
although its variance may be prohibitive. A latent model can impose
structure that improves precision, with a corresponding specification risk.

If historical staff used unrecorded information predictive of future
outcomes, an accurate $\Prob(A_t\mid H_t)$ is generally insufficient.
It describes selected historical actions without reconstructing hidden
information under a new policy. \citet{tennenholtz2020,bennettprox2023}
study this problem using specified information structures and informative
proxy or bridge conditions. Their results do not turn every fitted HMM
into a causal environment.

The empirical questions are about measurement: which private information
did staff use; which independently recorded proxies have defensible roles;
could limited randomization identify a missing transition; could a narrower
policy be evaluated with existing support? These questions can be more
valuable than increasing the representation dimension.

\subsection{A finite bridge exposes the inverse problem}
At fixed $(a,x)$ in Equation~\eqref{eq:bridge}, let
$M_{zw}=\Prob(W=w\mid Z=z,A=a,X=x)$ and
$m_z=\E[Y\mid Z=z,A=a,X=x]$. A discrete bridge solves
$M\boldsymbol h=\boldsymbol m$. Full column rank gives at most one
bridge in this class. A small minimum singular value lets modest noise
in $\boldsymbol m$ produce large changes in $\boldsymbol h$.
Regularization stabilizes the calculation while introducing bias.

This explains why two predictive but almost redundant proxies may be weak
for identification. It also limits a numerical diagnostic: empirical full
rank does not establish the substantive proxy exclusions or population
completeness. Sequential bridge methods additionally preserve temporal
roles and the target actor's information. The time-independent-sampling
identification construction and proximal estimating strategy discussed in
\citet{bennettprox2023} are distinct approaches, with different computational
and regularity requirements.

\section{Customer behavior and the bank's objective}
A bank POMDP is not automatically a model of customer utility. Distinguish
bank offer $b$ from household choice $a$. A declared dynamic logit
construction has
$v(a,s;b)=u(a,s;b)+\beta\E[\overline V(s')\mid s,a,b]$ and,
under centered independent Type-I extreme-value shocks of fixed scale
$\tau$, $\overline V(s)=\tau\log\sum_a\exp\{v(a,s)/\tau\}$.
Choice probabilities are proportional to $\exp\{v(a,s)/\tau\}$.
Utility normalization, discounting, information and transition restrictions
must be identified; a softmax-looking policy does not supply them.

Given a supported customer response law, the bank's induced objects are
\[
 K_b(s,s')=\sum_aP_\theta(a\mid s,b)F_\theta(s'\mid s,a,b),
 \qquad r_b(s)=\sum_aP_\theta(a\mid s,b)m(s,b,a).
\]
Household utility $u$ and bank contribution $m$ have different owners
and need not be proportional. Mortgage demand and lender behavior in
\citet{benetton2021}, and refinancing mechanisms in
\citet{agarwal2020,beraja2019}, motivate this separation. The synthetic
POMDP's supplied reward schedule does not estimate household utility.

Composition therefore needs three contracts: the response module declares
supported offer interventions; finance declares ownership, discounting
and terminal claims; control declares the bank's information and feasible
actions. A fitted incumbent-choice probability alone does not support an
unmodeled marketing or pricing intervention.

\section{A concrete initial module}
The initial product should support finite states, categorical and simple
parametric emissions, observed covariates, variable trajectory lengths,
explicit censoring and a declared logging-law family. Results contain
likelihood diagnostics, parameter and state uncertainty, transition support,
the observation clock and a separate causal interpretation. Prospective
trajectory views use filtering; retrospective views label smoothing.

Small models compare EM, direct optimization and exact path enumeration.
Planning consumes an explicitly accepted environment and preserves the
actor's information. Evaluation fixes a policy or runs an outer training
and evaluation protocol. Continuous states, particle likelihoods and
amortized inference have separate approximation and identification work.
BayesFilter can supply useful probabilistic infrastructure, but the causal
action model, intervention evidence and temporal semantics remain owned by
this product. No current bank or sibling integration is claimed.

\chapter{Causal discovery as constrained scientific search}\label{ch:discovery}
\section{What a graph represents}
A bank may want to understand whether deposit-rate changes precede withdrawals,
whether a service complaint mediates attrition, or whether an economic shock
affects several products. A discovery procedure can help organize hypotheses,
but its output must state the variables, time scale, sampling mechanism,
latent-variable assumptions and edge semantics.

A DAG specifies acyclic causal relations among variables. A CPDAG represents
a Markov equivalence class under an appropriate causally sufficient model.
A PAG represents invariant ancestral information across a class that can
include hidden confounding and selection. Circle endpoints encode uncertainty;
they are not weak versions of directed arrows. Causal graphs also differ
from data-lineage graphs, whose edges represent transformations or provenance.

Conditional independences alone often leave multiple DAGs compatible with
the same observed law \citep{peters2017}. Temporal precedence rules out
some directions, but contemporaneous dependence, subsampling, delayed
recording and hidden common causes still matter. A graph of monthly
aggregates need not preserve daily intervention semantics.

\section{Comparison of method families}
\begin{longtable}{P{27mm}P{42mm}P{38mm}P{37mm}}
\toprule Method family & Identifying structure & Output/use & Failure to examine\\\midrule
\endhead
PC/FCI & Markov and faithfulness assumptions; suitable independence tests; different latent-variable assumptions & CPDAG or PAG; background constraints & Selection, measurement error, weak signals, finite-sample CI errors\\
LiNGAM / structural VAR & Linear equations and independent non-Gaussian innovations; acyclic contemporaneous structure & Oriented coefficients under the model & Gaussianity, hidden confounding, aggregation\\
NOTEARS / DAGMA & Smooth acyclicity constraint plus a chosen score and structural model & Numerical DAG candidate & Nonconvex solutions, sufficiency, thresholding, score equivalence\\
PCMCI & Stationary lagged system with suitable parent conditioning and tests & Lagged temporal relations under assumptions & Unmodeled contemporaneous dependence and omitted lags\\
PCMCI+ & Sufficiency, stationarity, adequate lag, Markov/faithfulness, valid tests & Lagged graph and contemporaneous equivalence information & Selection and measurement error excluded by the theorem\\
LPCMCI & Latent-confounded temporal ancestral-graph assumptions & Temporal PAG & Lag-window marginalization and unresolved endpoints\\
DYNOTEARS & Structural VAR with independent errors and functional/noise restrictions & Contemporaneous and lag coefficients & Fixed graph, stationarity/model fit, latent confounding\\
\bottomrule
\caption{Method assumptions determine what the output may assert. Runtime and finite-sample accuracy need separate evaluation.}
\label{tab:discovery}
\end{longtable}

The PC literature separates pointwise consistency from stronger
high-dimensional uniform consistency requirements \citep{kalisch2007}.
LiNGAM and its structural-VAR extension obtain additional directionality
through non-Gaussian independent noise \citep{shimizu2006,hyvarinen2010}.
Granger-style prediction instead asks whether past values improve
conditional prediction. It becomes interventionally causal only with
additional structural assumptions; a common macro driver can create
predictive lead-lag relations without a direct causal effect.

\section{Continuous optimization is a computational result}
NOTEARS replaces a combinatorial DAG constraint with
\begin{equation}
 h(W)=\operatorname{tr}\{\exp(W\circ W)\}-p=0.
 \label{eq:notears}
\end{equation}
For $S=W\circ W\geq0$, each $\operatorname{tr}(S^k)$ sums weights
of closed directed walks. Expanding the matrix exponential shows that
$h(W)=0$ precisely when no directed cycle exists
\citep{notears2018}. This elegant characterization makes continuous
optimization possible; it does not establish which direction is causal.
DAGMA gives another smooth characterization using a log determinant
\citep{dagma2022}. Both still need a suitable structural model, regularization
and numerical checks.

DYNOTEARS adds lagged coefficients to a structural VAR and constrains the
contemporaneous graph \citep{pamfil2020}. Time order already constrains
lagged directions, while identification of within-period effects depends
on functional and innovation assumptions. A group-wide monthly model that
mixes accounting identities, price-setting and endogenous demand is not
automatically within those assumptions.

\section{Time-series discovery and latent causes}
PCMCI reduces irrelevant conditioning by estimating useful parent sets
before momentary conditional-independence tests \citep{runge2019}.
PCMCI+ adds contemporaneous links. Its inspected assumption statement
requires causal sufficiency, stationarity, sufficient maximum lag,
appropriate faithfulness and consistent tests, and excludes selection and
measurement error \citep{runge2020}. Banking measurement defects therefore
need explicit treatment before applying its theorem.

LPCMCI permits hidden common causes through temporal ancestral-graph
representations \citep{gerhardus2020}. The result is a PAG with uncertainty,
not a fully observed latent-state reconstruction. Changing the lag window
can change the marginal graph being represented. Store the window,
conditioning tests, significance thresholds, multiplicity treatment,
background constraints and bootstrap stability with every result.

Missing data adds another graph. Conditioning on observed records can create
selection dependencies; naive imputation can create apparently strong edges
\citep{tu2019}. A null result can reflect weak measurements or low power,
not absence of an economic mechanism. The discovery workspace must keep
measurement and missingness assumptions visible.

\section{Continuous time, feedback and recent directions}
Dynamic SCMs model variables as paths and define interventions at that level
\citep{boeken2024}. Connections between dynamical systems, interventions and
equilibria require solvability and stability assumptions
\citep{bongers2022}. Recent causal SDE results extend Markov properties and
do-calculus under specified solvability conditions \citep{sde2026}.
These are relevant to irregular transactions and continuous-rate dynamics,
but do not yet provide an automatically qualified bank estimator.

The recent search found several useful research comparators.
ITSY addresses irregular series but explicitly assumes causal sufficiency,
stationarity and MCAR \citep{itsy2026}. LUCID routes corrections according
to pervasive or local confounding signatures \citep{lucid2026}.
The Causal-TS paper describes nonstationary temporal workflows and accelerated
independence tests \citep{causalts2026}. Their author-reported performance
has not been reproduced here; the most recent versions are preprints.
They motivate benchmarks rather than immediate default dependencies.

Group-level causal background knowledge can reduce a large search problem
\citep{vargas2026}. For example, an agreed temporal ordering of eligibility,
assignment, receipt and outcomes excludes implausible edges. An enterprise
ontology alone supplies names and relationships, not that causal ordering.
Domain experts must justify it.

\section{A discovery-to-study workflow}
Start with a narrowly defined question and a time-respecting dataset.
Record domain constraints, plausible omitted variables and observation
mechanisms. Run transparent baseline methods and selected challengers;
compare equivalence classes, graph stability and sensitivity to aggregation,
lag windows and missingness.

Next, choose a small set of economically meaningful candidate mechanisms.
For each one, state a proposed intervention and ask whether its effect is
identified in every graph still under consideration. If not, report a range
or identify the experiment or measurement that would distinguish models.
Only then create an effect-study specification with a documented graph
choice. The change from discovered hypothesis to accepted identification
assumption must be an explicit event in the evidence record.

This workflow can reduce wasted study design and suggest better measurement.
It should not promise to discover the bank's unique true causal graph from
all data, especially when policies adapt and hidden information drives actions.

\section{A financial graph needs a unit and an intervention}
Consider daily balances, the posted deposit rate, customer withdrawals,
funding transfers and a market rate. ``Rate causes balances'' is incomplete.
The rate may be an account-specific offer, a regional posted rate or a
group-average accounting yield. Balances may be end-of-day stocks, intraday
liquidity or monthly averages. A causal edge between two aggregate series
does not necessarily describe the effect of changing one customer's offer.

The graph specification should therefore name the unit, aggregation window,
decision actor and hypothetical intervention for every actionable node.
Some nodes are accounting identities. End balance equals beginning balance
plus net flows and valuation adjustments. Treating all terms as independent
stochastic variables can create deterministic conditional relations that
violate the generic faithfulness assumptions of a discovery algorithm.
The data layer should record the identity and choose a scientifically
appropriate variable representation.

The distinction is useful even when the final graph remains uncertain.
An economist may discover that all plausible graphs leave a rate effect
unidentified without a funding-cost instrument or an experiment. That is
an actionable output: obtain the missing variation or narrow the claim.
A unique-looking graph created by an arbitrary tie-breaker would be less
useful because it conceals the unresolved measurement decision.

\section{Why observational equivalence survives unlimited data}
For three variables, the DAGs $X\to M\to Y$ and
$X\leftarrow M\to Y$ have the same skeleton and no unshielded collider.
Both can imply $X\perp Y\mid M$ while allowing marginal dependence.
Under ordinary observational conditional-independence assumptions, the
data cannot distinguish their causal direction. A collider
$X\to M\leftarrow Y$ instead predicts marginal independence and
conditional dependence when its other assumptions hold. This is why
constraint-based discovery records equivalence information
\citep{peters2017,kalisch2007}.

The economic consequences differ. In the chain, intervening on $X$ can
change $Y$ through $M$. In the common-cause graph, an intervention on
$X$ may leave $Y$ unchanged. More observations can estimate the common
observed law more precisely without resolving that causal ambiguity.
Temporal order or an intervention can add information; a more confident
conditional-independence test alone cannot.

A similar ambiguity occurs in the executed two-variable VAR. Its lag
matrix and innovation covariance are
\[
 A=\begin{pmatrix}.7&0\\.15&.6\end{pmatrix},\qquad
 \Sigma_u=\begin{pmatrix}1&.8\\.8&1\end{pmatrix}.
\]
Ordering funding first gives impact column $(1,.8)'$ for a unit funding
shock. Ordering it second gives normalized column $(1,0)'$.
Both decompositions reproduce exactly the same Gaussian reduced-form law.
Their loan responses are the second components of $A^h(1,.8)'$ and
$A^h(1,0)'$, plotted in Figure~\ref{fig:dependence-irf}. The impact
response is $0.8$ in one model and zero in the other. No forecasting
comparison can choose between them under that common law.

Independent non-Gaussian structural innovations can break some linear
equivalences under the assumptions of LiNGAM and its time-series extension
\citep{shimizu2006,hyvarinen2010}. This is additional model information,
not a consequence of fitting a VAR. Hidden common causes, Gaussian
innovations, misspecified lag dynamics and aggregated signals can invalidate
the argument. The discovery interface should show which restriction
contributed each orientation.

\section{Conditional-independence testing is a statistical component}
For jointly Gaussian variables, zero partial correlation corresponds to
conditional independence. One can residualize $X$ and $Y$ on a
conditioning set $S$, calculate their correlation, and use an appropriate
sampling law. For nonlinear relationships, kernel, nearest-neighbor or
learned tests require different regularity and calibration. A strong
nonlinear dependence can coexist with zero ordinary correlation.

Temporal dependence complicates the sampling distribution. A test designed
for iid rows cannot be applied blindly to overlapping lag vectors. Parent
conditioning in PCMCI seeks to remove irrelevant persistence and reduce
the conditioning dimension before its final momentary tests
\citep{runge2019}. PCMCI+ incorporates contemporaneous structure under
its stated sufficiency and stationarity restrictions
\citep{runge2020}. These methods improve a specified search problem;
they do not exempt the observations from their own measurement model.

High-dimensional search also makes the testing family adaptive. With
$p$ variables and $L$ lags, even one test per candidate has order
$p^2L$ comparisons. Conditioning-set exploration can add many more.
An edgewise nominal level does not by itself imply graphwise error control.
Any false-discovery correction must match the tests and dependence it
actually receives. Bootstrap selection frequency is useful stability
information, but it is not automatically a posterior edge probability
or a calibrated confidence level.

The proposed benchmark therefore measures at least three objects.
First, test calibration on known independent or dependent processes under
the relevant serial dependence. Second, graph recovery under declared
causal models, reporting adjacency and orientation errors separately.
Third, the downstream effect or decision loss induced by using the graph.
A small structural Hamming distance can still conceal the one orientation
that reverses an important funding decision.

\section{Acyclic optimization and time-series structure}
In a linear contemporaneous model $X=XW+U$, a squared-error score plus
sparsity penalty and Equation~\eqref{eq:notears} defines an optimization
problem. The acyclicity condition is exact for the weighted adjacency
matrix, but identification also depends on what distribution of $U$ and
what score distinguish observationally equivalent parameterizations.
Thresholding small coefficients after optimization can alter the fitted
distribution and the uncertainty of the selected graph.

For lagged systems, write
\[
 X_t=B_0X_t+\sum_{\ell=1}^{L}B_\ell X_{t-\ell}+\varepsilon_t.
\]
Only $B_0$ needs a within-period acyclicity restriction in a finite
unrolled time graph. Solving gives reduced-form coefficients
$A_\ell=(I-B_0)^{-1}B_\ell$ and innovations
$u_t=(I-B_0)^{-1}\varepsilon_t$. Estimating $A_\ell$ alone
therefore does not identify the structural $B_\ell$. DYNOTEARS fits
contemporaneous and lagged structure together under its structural model
\citep{pamfil2020}; the product should retain both matrices rather than
present reduced-form predictive coefficients as direct mechanisms.

Feedback can be economically real within an observation window. Rates
affect flows and flows induce repricing during the same month. Forcing
an acyclic monthly contemporaneous graph may be a convenient approximation,
not a correct mechanism. Finer timing, a cyclic structural model with
well-defined solutions, or a path-level intervention can be more suitable.
Dynamic SCM and dynamical-systems work provides a language for these
questions \citep{boeken2024,bongers2022}, while practical estimation
still requires an observation scheme and identifying restrictions.

\section{Aggregation, missingness and changing regimes}
Sampling every second period of a VAR(1) produces
\[
 X_t=A^2X_{t-2}+Au_{t-1}+u_t.
\]
Even if $A$ is sparse, $A^2$ can contain a nonzero $(i,j)$ entry
because of an intermediate path $j\to k\to i$. The coarser lag graph
can therefore contain an edge that represents a multi-step mechanism.
Its innovation covariance also changes. Learning the monthly graph and
interpreting each edge as a daily direct intervention would be a category
error, not merely a loss of temporal precision.

Irregular observations add selection. An account can be measured more
often after an adverse event; staff may record a customer conversation
only when assistance is considered. Treating the resulting absence as
MCAR can create or remove conditional dependencies. Missingness-aware
discovery requires assumptions about the observation process
\citep{tu2019}. ITSY's explicit MCAR and stationarity restrictions
\citep{itsy2026} are useful precisely because they reveal when the method
does not fit a bank's informative observation schedule.

A pooled graph across low-rate, crisis and high-rate regimes can also mix
mechanisms. One approach fits separate graphs in sufficiently supported
regimes; another models changing parameters or environments under declared
invariance restrictions. Separate fits reduce sample size; a shared graph
can force false stability. The choice should be evaluated with temporal
holdouts and mechanism-specific diagnostics. A change-point detector
locates distributional change without establishing which causal mechanism
changed or whether the cause was policy.

Recent LUCID and Causal-TS proposals address important confounding or
nonstationarity challenges \citep{lucid2026,causalts2026}. They should be
included as research comparators on banking-specific synthetic mechanisms
and carefully bounded real studies. Their newness and author-reported
performance do not establish that they solve aggregation, hidden logging
or measurement selection simultaneously.

\section{From graph uncertainty to a useful answer}
Suppose a discovery method returns a set of compatible graphs
$\mathcal G$. For each $G$, let $\Theta_G$ be the compatible effect
set under its assumptions. The union $\bigcup_{G\in\mathcal G}\Theta_G$
describes what remains compatible when graph uncertainty is retained.
If a candidate policy is economically preferable across that set, the
decision can be robust even without unique orientation. If the set spans
both a gain and a loss, reporting one graph's point effect would hide a
material uncertainty.

The representation must distinguish an equivalence class justified by
theory from a convenient collection of bootstrap graphs. The latter is a
stability sample and need not contain every compatible graph. Likewise,
a PAG endpoint carries ancestral information under a latent-variable
model; replacing a circle by an arrow because its score is slightly larger
changes the identification claim.

A research workbench can then rank additional observations or interventions
by their potential to resolve the consequential ambiguity. For example,
randomized rate offers in a supported population may distinguish a
rate-response edge, while an independent liquidity measurement may help
separate demand from funding pressure. The first release can present these
as analyst-authored study proposals tied to graph alternatives. Automatic
optimal experiment selection would require a model of information value
and operating constraints and is a separate research feature.

\section{The discovery module's construction contract}
Inputs include variable definitions, population, observation clock,
aggregation, lag window, background constraints, hidden-variable and
selection assumptions, conditional-independence test and multiplicity
policy. Outputs include graph type, endpoint marks, edge evidence,
unresolved alternatives, stability information and sensitivity to timing
and missingness. Every display of a causal effect identifies the graph
and additional assumptions used to compute it.

The product should initially wrap established temporal and latent-aware
algorithms, retain their native diagnostics and compare simple structural
baselines. It should use domain groupings to reduce search only when those
groupings supply justified background restrictions
\citep{vargas2026}. Data-lineage links, semantic similarity and causal
arrows remain different edge types. Their integration in one interface
is useful only if their meanings remain explicit.

\part{Banking studies from decision to value}
\chapter{A deposit-policy study from rollout to funding value}\label{ch:deposits}
\section{The owner's decision and the financial account}
Suppose the group funding committee is considering a temporary savings-rate
enhancement. The retail business expects retained deposits; treasury values
the funding; finance pays interest on qualifying balances; and the customer
receives a higher rate. The relevant bank decision is whether the incremental
funding value exceeds the incremental interest and operating costs over a
specified horizon. An effect on account balances is one input to that
decision, not its answer.

The committee needs three quantities kept separate: balances in the promoted
product, total external funding retained by the group, and incremental
economic contribution. An internal transfer from another group account can
increase the first with little change in the second. Paying a higher rate
on incumbent balances can make the third negative even when both balance
measures rise. This decomposition should be visible in the first report.

The empirical literature suggests mechanisms and designs. The deposits
channel in \citet{drechsler2017} connects monetary policy, local deposit
market power, funding flows and lending. Its within-bank branch comparison
uses geographic concentration while controlling bank-specific opportunities;
its interpretation relies on banks being able to move funding between
branches. The paper's subsequent lending comparisons operate at different
aggregation levels. It does not supply the effect of a new bank's account-
level retention offer. The product should make that distinction explicit
when an analyst imports a mechanism from a published study.

\section{The data contract before estimation}
The proposed bank implementation begins with a fixed eligible cohort before
announcement. It preserves branch and household identity as of that date,
the original assignment, the announced terms, delivery, take-up and later
closure. It then reconstructs daily balances, interest, external flows and
internal transfer links under a stated accounting cutoff. A corrected
transaction changes a later readout while leaving the original information
set and assignment intact.

\begin{table}[htbp]\centering\small
\begin{tabularx}{\textwidth}{P{31mm}XX}
\toprule Object & Required meaning & Consequence if absent\\\midrule
Eligibility & Population defined before announcement & Survivor-only analysis changes the target\\
Assignment & Offer terms, duration, branch/cohort and date & Acceptance cannot substitute for assignment\\
Balance & Stock, currency, product and economic date & Daily and monthly-average balances are not interchangeable\\
Flows & External source or linked internal transfer & Group funding cannot be inferred from one product's growth\\
Interest & Actual accrual and promotional increment & Repricing of incumbent funds is omitted\\
Funding value & Dated marginal funding convention & Principal may be mistaken for contribution\\
Observation & First available date and correction lineage & Retrospective leakage changes the study information set\\
\bottomrule
\end{tabularx}
\caption{The proposed deposit study's minimum economic objects. A bank may
support a narrower balance study before all funding objects are available.}
\end{table}

Households moving between branches require a declared attribution rule.
Freezing baseline assignment preserves an intention-to-treat interpretation;
following current branch membership can select on the policy response.
Cross-branch displacement also challenges untreated donors. A group-wide
study can use household-level outcomes while clustering by randomized or
rollout branch, provided the mapping is stable and spillovers are addressed.

\section{A complete synthetic rollout}
The executed study observes 96 branches for 30 months. Twenty-four branches
adopt in each of months 12, 18 and 24, with 24 never adopting. Months are
indexed from zero. Cohort labels are fixed and branch disturbances are
independent of them. Untreated balances, in millions of dollars, are
\begin{align}
 B_{it}(0)&=30+\alpha_i+0.07t+0.4\sin(t/3)+\epsilon_{it},\\
 \alpha_i&\sim N(0,3^2),\qquad
 \epsilon_{it}=0.65\epsilon_{i,t-1}+\eta_{it},\quad
 \eta_{it}\sim N(0,0.28^2),
 \label{eq:deposit-dgp}
\end{align}
with the first disturbance drawn at the innovation scale. For cohort $g$
and event month $k=t-g\geq0$, the balance effect is
\[
 \tau_{gk}=\{0.85+0.025(g-12)\}(1+0.14k).
\]
This law imposes untreated parallel trends while allowing heterogeneous,
growing effects. It is designed to compare targets and estimators rather
than simulate every feature of deposit demand.

For every supported $(g,t)$, estimate the change from $g-1$ among that
cohort minus the same change among never-treated branches. Average all
treated cohort-month cells equally. Because all cohorts have 24 branches,
this equals an equal-weight average over treated branch-months. It gives
early cohorts more total weight because they contribute more treated
months. A bank wanting an equal-cohort or balanced-event-window target
would use different declared weights.

The challenger fits a single treatment indicator with branch and month
effects. It estimates a different weighted combination when effects evolve
with adoption time, as explained in Chapter~\ref{ch:panels}. The study
also calculates a known reference by comparing the generated potential
outcomes for the same branch-months. Those potential outcomes are available
only to the synthetic evaluator.

\section{Converting balances into contribution}
The teaching account assumes an annual net funding margin of $2.5\%$
on balances before the promotional increment, a 25-basis-point promotional
rate increase on all treated balances, and USD\,1,800 monthly program
cost per treated branch. Under the stipulated monthly convention,
\begin{align}
 C_{it}(1)&=10^6 B_{it}(1)\frac{0.025}{12}
           -10^6 B_{it}(1)\frac{0.0025}{12}-1800,\\
 C_{it}(0)&=10^6 B_{it}(0)\frac{0.025}{12}.
 \label{eq:deposit-cash}
\end{align}
Subtracting gives
\begin{equation}
 \Delta C_{it}=10^6\frac{0.025-0.0025}{12}\Delta B_{it}
       -10^6\frac{0.0025}{12}B_{it}(0)-1800.
 \label{eq:deposit-hurdle}
\end{equation}
The second term is the cost of repricing balances that would have remained.
It is easy to omit when a business case values only additional deposits.

At an untreated balance of USD\,30 million, break-even requires
\[
 \Delta B>
 \frac{0.0025(30)+12(1800)/10^6}{0.025-0.0025}
 \simeq 4.293\ \hbox{million dollars}.
\]
The hurdle follows directly from the financial account. It changes with
rate duration, funding value, churn, servicing and the share of incumbent
balances receiving the offer. The numerical study holds those terms fixed
so that the effect of repricing is transparent.

\section{Fitted results and uncertainty}
The same cohort-time estimator is applied separately to balances and
observed contribution. Whole branches are resampled within each adoption
cohort in 499 bootstrap samples, retaining each trajectory and the fixed
cohort counts. The percentile intervals condition on this rollout design
and the supplied common macro path. They are not a repeated-sample
coverage qualification for a bank rollout.

\begin{table}[htbp]\centering\small
\begin{tabular}{@{}lrrrr@{}}
\toprule
Outcome & True effect & Cohort ATT & 95\% interval & TWFE\\
\midrule
Balances (USD m) & 1.780 & 1.689 & [1.564, 1.816] & 1.442\\
Contribution (USD) & -5030.635 & -5218.875 & [-5498.095, -4935.700] & -5607.561\\
\bottomrule
\end{tabular}

\caption{One fixed synthetic rollout. The reference is the realized
finite-population potential-outcome contrast; intervals resample complete
branches. Positive balance effects coexist with negative contribution.}
\label{tab:deposit-results}
\end{table}

\begin{figure}[htbp]\centering
\includegraphics[width=\textwidth]{derivations/deposit_study.pdf}
\caption{Cohort balance responses and the funding-contribution comparison.
The differing event profiles matter for aggregation; their positive signs
do not establish a profitable promotion.}
\label{fig:deposit-study}
\end{figure}

The contribution estimate is approximately USD\,$-5,219$ per treated
branch-month, with interval $[-5,498,-4,936]$. Under the stipulated design
and account, the rollout loses money despite retaining balances. The
appropriate decision is to reject this broad offer or investigate a
different, explicitly evaluated policy. Selecting a more favorable
balance estimator would not remove the cost of repricing incumbent funds.

The balance and contribution estimates share outcomes and branches.
For a portfolio of more detailed cash components, the implementation
retains their joint bootstrap or influence contributions. It should not
form a purported interval for total contribution by adding separately
calculated marginal endpoints. The simple account here allows contribution
to be observed directly, avoiding an unnecessary two-stage uncertainty
approximation.

\section{A targeted alternative is a new study}
A natural response is to offer the enhancement only to households likely
to withdraw external funds. That may reduce repricing cost, but it changes
the intervention and target population. A withdrawal-risk model predicts
which customers leave under historical practice. It does not necessarily
predict whose withdrawal is preventable by the offer. High-risk customers
may be insensitive to price; low-risk customers may accept a subsidy
without changing their behavior.

The appropriate targeting comparison uses conditional effects on the
group funding account and evaluates the whole targeting rule on held-out
or randomized data. It includes contact costs, take-up, rate duration and
customer constraints. The existing static example demonstrates why a
known positive average effect does not establish that one targeted policy
beats another at the declared simultaneous confidence level.

For repeat offers, the problem becomes longitudinal. A first offer changes
balances and may teach customers to wait for future offers. Selection for
the next offer then depends on a variable affected by the first.
The sequential methods in Chapter~\ref{ch:longitudinal} are needed to
compare sustained strategies. A one-period uplift score should not be
repeated indefinitely under an assumption of unchanged effects.

\section{What can transfer from the literature}
\citet{drechsler2017} supports a mechanism linking market power,
deposit pricing and transmission, with specific within-bank and local
comparisons. \citet{basten2026} investigates relationship stickiness and
cross-selling using mover and prior-location variation. Both suggest
heterogeneity worth measuring. Neither establishes that a bank's current
customers share the same elasticities, product alternatives or information.

Refinancing evidence adds a complementary lesson. In
\citet{agarwal2020}, eligibility for HARP provides a policy comparison;
the paper distinguishes a program effect from the selected experience
of borrowers who refinance and studies competitive frictions in pass-through.
This is analogous to keeping an assigned retention offer distinct from
acceptance. \citet{beraja2019} further explains why local responses do
not settle aggregate policy effects. A bank-wide price change can alter
competitor behavior and internal funding substitution beyond the local
partial-equilibrium study.

The platform should consequently store an evidence-transfer statement:
which mechanism or design comes from the paper, which coefficient is
estimated on owned data, and which equilibrium or stability assumptions
remain scenario inputs. Literature informs construction without becoming
an unexamined parameter database.

\section{The reusable application and acceptance test}
The deposit application exports a study specification, reconstructed cohort,
balance/flow/contribution definitions, cohort-time effects, covariance,
target weights, trend sensitivity and a funding report. A FinMarketing
consumer can use the supported causal response in its own financial
workflow if the population, intervention and account match. A treasury
consumer can use the funding contrast while retaining its transfer-price
and liquidity assumptions.

Bank acceptance would require reconciliation to authoritative balances
and interest, verified external-versus-internal flows, a defensible rollout
or experimental design, explicit spillover analysis, adequate independent
branches and outcome maturity. The equivalent-task comparator is a bank
economist building the same study with existing panel software and the
warehouse. The product earns its role if it reconstructs and updates that
study more reliably and cheaply while preserving the same scientific
answer, including this kind of unfavorable result.

\chapter{Repeated collections decisions and policy value}\label{ch:collections}
\section{Define assistance before optimizing it}
Collections is a natural longitudinal application because an action changes
the state that determines later action. A payment reminder, temporary
forbearance and a contractual restructuring are different interventions.
They affect cash timing, costs, eligibility and customer outcomes differently.
The initial product should evaluate a small set of operationally defined
assistance strategies, with ordinary servicing as the comparator, before
attempting unrestricted policy optimization.

The economic owner matters. A servicer, a loan owner and a purchaser of
distressed debt receive different cash flows. The study preserves the
original eligible cohort, includes assistance delivery costs and defines
the remaining claim at the horizon. It reports customer outcomes separately,
such as persistent hardship or repeated delinquency, rather than assuming
that bank contribution fully represents them.

The following complete synthetic study isolates the observed-history case.
The controlled hidden-state study in Chapter~\ref{ch:latentestimation}
then shows what changes when relevant information is hidden from the analyst.
Neither study represents a named bank's customers or treatment effectiveness.

\section{A three-state decision process}
The observed pre-action state is $S_t\in\{0,1,2\}$, representing
current, moderate stress and severe stress. Initial proportions are
$(0.6,0.3,0.1)$. Extra assistance is feasible only in states one and two.
The historical action-one probabilities are $(0,0.35,0.60)$.
These are known randomized probabilities conditional on the observed state;
there is no hidden action confounding in this experiment.

The transition matrices are
\[
 P_0=\begin{pmatrix}.88&.10&.02\\.30&.50&.20\\.08&.22&.70\end{pmatrix},
 \qquad
 P_1=\begin{pmatrix}.88&.10&.02\\.50&.40&.10\\.18&.42&.40\end{pmatrix}.
\]
The first row for action one is not used because assistance is infeasible
there. The process has eight decisions and monthly discount $\gamma=0.99$.
Contribution associated with the next state is $(20,-10,-160)$ dollars,
less USD\,12 for each assistance action. These are signed monthly
contribution increments, not a repeated write-off of the same principal.
A terminal adjustment of USD\,$-100$ in severe stress represents the
remaining account beyond the eight-month cash window.

This law lets assistance improve future states while costing money today.
It also creates practical support difficulties: the probability of following
a deterministic strategy over all eight dates can be small even though
each feasible action has positive probability at every state. The known
finite Markov law gives an exact policy-value reference without requiring
Monte Carlo approximation of truth.

\section{Policies, training information and exact evaluation}
Three fixed policies are declared before fitting: no extra assistance,
assistance only in severe stress, and assistance in all eligible states.
They use observable information and respect the same feasible action set.
The study draws 3,000 training trajectories and 8,000 independent evaluation
trajectories. Training data estimate transition rows by observed counts
with a one-half pseudocount in each next-state category.

For each target policy $\pi$, backward recursion under the fitted
transition law gives
\[
 \widehat Q_t(s,a)=\sum_{s'}\widehat P_a(s,s')
 \{r(s')+\gamma\widehat V_{t+1}(s')\}-12a,
 \quad \widehat V_t(s)=\sum_a\pi(a\mid s)\widehat Q_t(s,a).
\]
The terminal condition is the declared remaining account. Repeating the
recursion with the generating transition law gives the exact reference
value, available only to the evaluator. The fitted policy scores never
consume generated latent variables or future information unavailable at
the relevant decision.

The direct method averages $\widehat V_0(S_0)$ over evaluation initial
states. Stepwise importance sampling reweights each dated reward by the
full product of target-to-behavior action ratios through that date.
Sequential DR adds cumulative-ratio Bellman residuals to the direct
prediction. All three target the same discounted contribution and
terminal account. This alignment is essential for a fair comparison.

\begin{table}[htbp]\centering\small
\begin{tabular}{@{}lrrrrr@{}}
\toprule
Policy & True USD & Direct & IPS & DR & Final ESS\\
\midrule
No extra assistance & -139.58 & -123.44 & -135.08 & -137.56 & 146\\
Severe state only & -58.26 & -49.57 & -59.95 & -59.55 & 1267\\
All eligible & -2.34 & 0.48 & 2.93 & 1.70 & 489\\
\bottomrule
\end{tabular}

\caption{Eight-decision fixed-policy values in dollars per original eligible
customer. Final ESS describes concentration of final trajectory weights;
it is not the independent sample size of every dated score.}
\label{tab:collections-values}
\end{table}

The direct estimates depend entirely on the fitted transition law.
Importance sampling uses known logging and remains highly variable on
rare complete paths. DR corrects fitted continuation errors but cannot
remove all realized transition and reward uncertainty. Its empirical
value need not lie between the other two estimates.

\section{A plausible shortcut estimates the wrong law}
Suppose the analyst uses only the current-action ratio at each date.
This changes the distribution of the action conditional on the historical
incoming state, but leaves that incoming state distributed under historical
practice. Earlier assistance has already changed who reaches severe
stress. The shortcut therefore does not evaluate a sustained target policy.

For the all-eligible strategy in this run, that current-ratio calculation
gives approximately USD\,$-36.79$, compared with exact value USD\,$-2.34$.
For no extra assistance it gives about USD\,$-82.67$, compared with exact
value USD\,$-139.58$. The error is not fixed by a larger evaluation sample:
it is a mismatch between the reweighted distribution and the requested
intervention. A marginalized state-ratio estimator would need the target-
to-behavior state occupancy ratio as well as the action ratio.

The software consequence is a required weight-scope declaration.
``Propensity-weighted'' is too vague. A result should distinguish
current-action, full-history, marginal occupancy and censoring weights,
with their target and normalization. A consumer can then reject a score
that lacks the history correction required by its policy contrast.

\section{Evaluate contrasts with their covariance}
The practical question is the incremental contribution relative to the
incumbent. Compute both policy scores on each evaluation trajectory and
form their difference. The standard error of the mean uses the variance
of this paired difference, preserving covariance from shared customers.
There are two prespecified comparisons against no extra assistance;
Bonferroni-adjusted normal intervals give a conservative simultaneous
approximation for that family.

\begin{table}[htbp]\centering\small
\begin{tabular}{@{}lrrr@{}}
\toprule
Policy versus no extra & True USD & DR difference & Simultaneous interval\\
\midrule
Severe state only & 81.32 & 78.01 & [44.58, 111.44]\\
All eligible & 137.24 & 139.26 & [103.83, 174.69]\\
\bottomrule
\end{tabular}

\caption{Prespecified contrasts using paired DR scores. These intervals
condition on the training fits and on the synthetic complete-trajectory law.}
\label{tab:collections-contrasts}
\end{table}

The account and intervention are deliberately simple enough to reconstruct.
For a bank implementation, inference must also handle estimated observation
probabilities, household or regional clustering, delayed losses and policy
selection. The positive synthetic comparisons cannot be treated as measured
collections benefit for the investment case.

\section{Double robustness and interval reliability differ}
The repeated experiment uses 160 independent datasets, each with 600
training and 2,000 evaluation trajectories. It evaluates the all-eligible
policy under four nuisance specifications: fitted transition values and
true logging; zero interior action values and true logging; exact action
values and a wrong logging model assigning probability $0.5$ in eligible
states; and both the zero action values and wrong logging.

The terminal account is correct in all four cases. The first two illustrate
the protection of known assignment; the third isolates the correct-value
route; the fourth has no double-robustness protection. The generating law
and declared policies remain fixed across repetitions, while the independent
datasets vary. No favorable subset of repetitions is selected.

\begin{table}[htbp]\centering\small
\begin{tabular}{@{}lrrr@{}}
\toprule
Specification & Bias (USD) & RMSE & 95\% coverage\\
\midrule
DR & 0.15 & 17.33 & 0.938\\
Wrong Q; true logging & 0.71 & 18.28 & 0.869\\
True Q; wrong logging & 0.55 & 10.72 & 0.912\\
Both wrong & 1.41 & 12.67 & 0.856\\
\bottomrule
\end{tabular}

\caption{Repeated-sample error and nominal normal-interval coverage for
all-eligible policy value. Monte Carlo standard errors for coverage range
from about 0.019 to 0.028 and are retained in the result file.}
\label{tab:collections-robustness}
\end{table}

The zero-value/true-logging case has mean protection but coverage well below
95\% in this finite experiment. Highly variable path weights make a
normal approximation unreliable at this sample size. Even the exact-value
case does not guarantee nominal finite-sample coverage. The wrong-both
case happens to have modest average bias under this particular law; that
is not evidence of a new robustness property. A misspecified estimator can
perform tolerably under one mechanism and fail under another.

The product recommendation is therefore to report weight concentration,
extreme score contributions, horizon sensitivity and repeated-sample
calibration on relevant synthetic mechanisms before selecting an interval
method. Weight truncation can improve stability while introducing bias;
any chosen threshold is part of the estimator. A shorter feasible policy
comparison or a prospective experiment may be more informative than a
nominally sophisticated long-horizon estimate.

\section{What partial observation adds to this study}
Observed stress states in the present experiment are sufficient for the
transition and action laws. Real collections staff may observe private
conversation details, perceived ability to pay or impending income events.
If those signals also affect action assignment and outcomes, the
recorded-history exchangeability used here fails.

Chapter~\ref{ch:latentestimation} constructs that alternative explicitly.
Signals imperfectly measure a hidden state, and historical assistance
depends on that state. A joint action--observation likelihood improves
recovery under the stated parametric environment but retains substantial
policy uncertainty. Proximal and instrumental-variable OPE offer other
routes under different measurements and restrictions
\citep{tennenholtz2020,bennettprox2023,xu2023}. A bank cannot choose
among those routes solely by whichever model gives the largest estimated
recovery gain.

The first bank pilot should therefore investigate the action log alongside
the outcome data. It must document which staff information was available,
which variables were recorded before decisions, how discretion operated,
and whether an existing randomized or quasi-experimental variation can
anchor transition effects. That work determines whether the module returns
an identified policy contrast, a conditional scenario or a bounded effect.

\section{A deployable study application, before an automated policy}
The standalone collections workbench would display complete eligible
trajectories, action and observation clocks, fitted continuation values,
support diagnostics, fixed-policy comparisons and customer-outcome
constraints. It would retain the original cohort after default, sale or
closure and explain how each event enters the economic account. A model
developer can investigate latent dynamics without promoting a policy.

A downstream application can consume the evidence package to compare
approved strategies. Promotion to operational decision making additionally
requires current permissions, execution capability, customer protection,
prospective or otherwise adequate value evidence and ongoing monitoring.
Those are implementation obligations of the proposed bank product, not
claims established by the present synthetic studies.

\chapter{Credit supply, refinancing and macro transmission}\label{ch:credit}
\section{A group-level causal product needs several levels of analysis}
Retail treatment effects do not exhaust the questions of a universal bank.
The institution must understand how funding conditions change credit supply,
how borrowers substitute between lenders, and how mortgage refinancing
transmits policy to household cash flow. Those questions use bank-firm
networks, market-level variation and aggregate time series. They also expose
why the product cannot apply one customer-level identification strategy
to every business problem.

Three published designs illustrate complementary information. The liquidity
shock study of \citet{mian2006} compares differently affected banks lending
to the same firm. \citet{agarwal2020} uses exposure to a refinancing
program to distinguish policy effects from borrowers' selected refinancing
choices. \citet{benetton2021} estimates mortgage demand and lender behavior
to study market-wide regulation. Each makes additional claims that require
additional assumptions as the analysis moves from local variation toward
aggregate counterfactuals.

\section{Separating loan supply from borrower demand}
Let $\Delta L_{fb}$ be loan growth for firm $f$ at bank $b$, and
let $Z_b$ measure an externally induced funding shock. A simple model is
\begin{equation}
 \Delta L_{fb}=\alpha_f+\beta Z_b+\epsilon_{fb}.
 \label{eq:bankfirm}
\end{equation}
The firm effect absorbs its common change in demand across lenders. For
firms with several bank relationships, subtracting the within-firm mean
gives
\[
 \Delta L_{fb}-\overline{\Delta L}_{f\cdot}
 =\beta(Z_b-\bar Z_{f\cdot})+
   (\epsilon_{fb}-\bar\epsilon_{f\cdot}).
\]
Identification comes from different lender shocks facing the same borrower,
under the restriction that their residual relationship-specific demand
components do not co-vary with those shocks. A firm's general investment
opportunity can correlate with the shocks of its chosen banks and still
be removed; a bank-specific export-credit need correlated with the funding
shock may remain a confounder.

The distinction is explicit in the archived 2006 paper
\citep[Section III.B, equation 5]{mian2006}. It also notes the importance
of unanticipated shocks and the restriction to multiple-bank borrowers.
The later question of total firm borrowing cannot simply insert the same
firm effect after aggregating to one firm-level observation. Borrower
substitution and aggregate real effects need additional evidence.

\section{An executed matched-lender study}
The synthetic construction has 600 firms, 24 banks and three relationships
per firm. Bank shocks $Z_b$ and firm demand terms $U_f$ are initially
independent standard normals. Relationship selection favors banks with
similar shock exposure through probabilities proportional to
$\exp(0.65U_f Z_b)$. This produces correlation between a firm's demand
and its observed lenders' shocks. Generated loan growth is
\begin{equation}
 \Delta L_{fb}=2-3Z_b+4U_f+\eta_b+\varepsilon_{fb},
 \label{eq:credit-dgp}
\end{equation}
with bank noise standard deviation $0.45$ and relationship noise standard
deviation $1.5$. Units are percentage points of loan growth per unit
funding shock. The true supply coefficient is $-3$ by construction.

Pooled regression leaves selected borrower demand in the error. Firm
demeaning removes the common $4U_f$ term. The study estimates both on
the same 1,800 relationships. The within estimator's illustrative interval
clusters residual score contributions by bank and uses a $t_{23}$
critical value. There are 24 independent bank shocks, not 1,800 independent
policy shocks. The demonstration is one sample under the stipulated law;
it does not qualify few-cluster inference for arbitrary lending networks.

\begin{table}[htbp]\centering\small
\begin{tabular}{@{}lrr@{}}
\toprule
Method & Coefficient & 95\% interval\\
\midrule
Generating effect & -3.000 & --\\
Pooled & -1.080 & --\\
Borrower effects & -2.923 & [-3.027, -2.819]\\
\bottomrule
\end{tabular}

\caption{Synthetic matched-bank loan-growth comparison. Relationship
selection biases the pooled association, while borrower effects recover
the stipulated supply coefficient much more closely.}
\label{tab:credit-supply}
\end{table}

The estimate changes from approximately $-1.08$ pooled to $-2.92$ with
borrower effects. The difference is economically substantial. It arises
because the design compares lenders serving the same borrower, rather
than because fixed effects are a universal solution to endogeneity.
If relationship-specific demand were correlated with $Z_b$, the within
estimator could remain biased. If all a firm's lenders had identical
exposure, that firm would contribute no within-firm identifying variation.

\section{What the bank would need to observe}
A single bank sees its own lending, not necessarily the same firm's other
relationships. It may need an appropriately permitted credit registry,
syndicated-loan dataset or borrower-supplied exposure record to implement
this design. A model for outside borrowing does not substitute for those
measurements without new assumptions. Entity resolution must preserve
legal borrower, consolidated group, facility and lender roles; several
facilities at one bank do not create several independent lender shocks.

The source snapshot must retain pre-shock relationships to avoid conditioning
on post-shock survival or new borrowing. Committed lines, drawn balances,
new originations, prices and collateral changes are different outcomes.
A borrower drawing an existing commitment during a shock can increase
observed balances while the bank tightens new credit. The product therefore
requires an outcome type before fitting the panel.

For a group credit committee, the first report can show the identified
relative supply response and its supported borrower population. A second
scenario may translate loan changes into funding, margin and capital
requirements. That translation needs the contract account, maturities,
loss assumptions and borrower substitution. The platform should preserve
the first empirical result even if the second remains model-dependent.

\section{Refinancing: eligibility, take-up and ownership}
Refinancing illustrates another frequent causal mistake: comparing borrowers
who refinance with those who do not. Refinancing choices respond to
attention, creditworthiness, equity, expected mobility and competing offers.
Those factors can also predict later spending and default. An observed
post-refinancing improvement does not by itself estimate the effect of
an offer or policy that induces refinancing.

\citet{agarwal2020} uses HARP eligibility to define treatment exposure,
with conforming high-LTV loans and a comparison group of observationally
similar non-agency loans. The paper examines pre-program comparisons,
refinancing, borrower savings and broader spending outcomes, and studies
competitive frictions in pass-through. The appropriate product lesson is
the separation of program exposure, take-up and downstream outcomes. Its
design and estimates depend on that institutional setting and should not
be transplanted unchanged to a new bank's retention program.

The bank also needs an ownership account. A lower mortgage rate can improve
the household's cash flow while reducing the value of a mortgage asset
held by the group. Origination or servicing fees may accrue to a different
entity. Paying off an old claim and originating a replacement cannot be
valued by adding the new loan's gross interest while ignoring the old
asset, funding and principal movements. The causal module supplies a
supported behavioral contrast; the mortgage account supplies the signed
cash and continuation value for each owner.

A practical study can define assigned refinancing assistance as treatment,
measure application, offer, acceptance and closing as separate stages,
and value the original eligible cohort. Failed delivery and non-take-up
remain in an intention-to-treat analysis. A complier effect or mechanism
analysis needs additional assumptions. Outcome maturity can be long even
when take-up is observed quickly, which affects both study power and the
product investment calendar.

\section{Why equilibrium counterfactuals need more structure}
A broad change in capital requirements or a group-wide mortgage price
policy can alter competitor prices, market shares and loan amounts.
\citet{benetton2021} models borrowers' discrete product choice and
continuous loan demand, together with multiproduct lender pricing. It
uses variation in predetermined risk-weighted capital requirements as
cost shifters and separates lender and product variation. This construction
allows equilibrium questions beyond a local reduced-form effect, under
the demand, supply, exclusion and behavioral restrictions of that model.

For the proposed platform, this is a reason to expose a structural extension
interface. A fitted local response may be an input or a validation target
for a structural model, but it does not determine the entire market
counterfactual. Supply response, outside options and competitive conduct
remain separate parameters or assumptions. Good market-share prediction
under the incumbent is informative about fit and incomplete evidence about
an unprecedented policy intervention.

\citet{beraja2019} makes the aggregation issue particularly relevant to
monetary transmission: regional variation can reveal heterogeneous responses
while common aggregate effects cancel from the comparison. Adding local
coefficients does not recover what was differenced out. A group-level
report should distinguish directly identified local effects, reweighted
population effects and equilibrium model scenarios.

\section{Treasury transmission and the shock catalogue}
The time-series workspace should preserve a catalogue of candidate shocks:
their economic definition, observation timestamp, normalization, source,
instrument construction and maintained exclusion restrictions. A market
rate change, a policy-announcement surprise and a bank-specific liquidity
shock are different objects. A discovered lead-lag relation can suggest
a candidate mechanism, but it does not turn an ordinary residual into an
identified structural shock.

For a selected shock, local projections and structural VARs can be
challengers with a common outcome, horizon and normalization. The external-
instrument argument in \citet{stock2018} requires the relevant dynamic
exclusions; \citet{jorda2024} explains estimation and inference alternatives.
The report should show response covariance, weak-instrument diagnostics,
regime sensitivity and the difference between a shock response and a
sustained policy path. Few independent macro episodes remain a limitation
even in a bank with billions of transaction records.

\section{What the shared product contributes}
These applications share temporal reconstruction, economic entity roles,
source lineage, score covariance and evidence export. They retain distinct
identification routes: matched borrowers for supply, policy eligibility
for refinancing, structural demand and supply for equilibrium, and external
shock restrictions for macro transmission. This is a stronger reusable
product than a collection of regressions with common plotting defaults.

MacroFinance can consume a supported shock or empirical response while
retaining its structural equilibrium assumptions. SMEwallet can consume
an observed or identified relationship response while preserving uncertainty
about outside borrowing. FinMarketing can use an eligible-cohort response
in a compatible mortgage or relationship account. The proposal specifies
those exchanges; it does not certify current adapters or endorse a copied
historical repository milestone as a current qualification.

\part{Data, software and a reusable product}
\chapter{Automatic organization of banking data}\label{ch:data}
\section{Organize around decisions}
The product needs more than an automated feature table. It needs an account
of which entity could receive which action, what was known when, what
actually happened, and which outcomes belong to that decision.
The basic hierarchy is legal entity, customer or firm, household or connected
group, account, product, decision episode and event. Relationships can change
over time and must be versioned.

Separate four graphs. The enterprise metadata graph describes datasets,
fields, ownership and business terms. The entity graph links people,
households, firms and accounts with uncertainty. The lineage graph records
transformations and versions. The causal graph represents intervention
assumptions. Cross-references between them are valuable; merging their edge
semantics is a scientific error.

\section{The event and availability clocks}
Every event needs occurrence time, economic effective time, system recording
time and the time it became available to the relevant decision process.
These need not coincide. A backdated arrears correction or revised income
field may describe the past while being unavailable at the past decision.
Retrospective warehouse completeness can therefore create prospective leakage.

For a feature used at decision $d$, require
\begin{equation}
 t_{\mathrm{available}}(f)\leq t_d,\qquad
 t_{\mathrm{event}}(f)\leq t_d,
 \label{eq:availability}
\end{equation}
for historical measurements, together with a valid entity mapping and
permission context. An explicitly known future contractual schedule is a
different feature class and needs its own rule. For a derived feature,
availability is no earlier than the latest input availability and its actual
publication to the decision system. An event-time join alone is insufficient.

Retain UTC instants, original time zones and applicable business calendars.
A statement date, settlement date and decision timestamp have different
semantics. Late-arriving data should produce a new analysis snapshot while
preserving the original replay. Restating an old feature in place destroys
the evidence needed to reconstruct the historical decision.

\section{A proposed organization pipeline}
Connectors first inventory schemas and permitted metadata, profile fields,
and propose types, units and entity keys. Statistical matching and language
models may suggest mappings to a banking glossary. Confidence scores
prioritize human review; they do not certify that a column means assigned
treatment, accepted offer or available information.

Next, identity resolution produces deterministic links where authoritative
keys exist and probabilistic links otherwise. A false household merge can
create apparent interference; a missed merge can leak information across
folds. Preserve match probabilities, link provenance and alternative plausible
matches. \texttt{Splink} is a candidate linkage component
\citep{sw_splink}; downstream causal sensitivity remains the platform's
responsibility.

The event builder normalizes money, calendars, action codes and state
transitions. It retains raw units and conversion policies, and validates
stock-flow identities. An account-month table should distinguish a true
zero, a structurally inapplicable value, a missing measurement and a value
not yet available. The trajectory builder then assembles eligibility,
actions, receipts, outcomes and censoring on the declared decision clock.

Finally, the study compiler checks temporal order, independent-unit keys,
population support and the compatibility of the requested estimand with the
available design. It outputs a frozen cohort manifest and a list of
unresolved assumptions. It cannot infer conditional exchangeability from
schema quality.

\section{Treatment and outcome provenance}
An action record should contain assignment probability where known, policy
version, eligibility rule, decision time, assigned treatment, delivery,
receipt and override reason. Overrides and relationship-manager information
can be part of the assignment mechanism. Their absence is a model issue,
not simply a missing-data nuisance.

For financial outcomes, use a component ledger. Deposit funding benefit,
interest expense, servicing cost, credit loss and terminal value each have
one owner, a valuation date and a scenario definition. Internal transfers
must reconcile across products. A group effect aggregates compatible
components; it should not add product-level gains that count the same
customer movement twice.

Terminal values need explicit continuation assumptions. A twelve-month
experiment cannot nonparametrically identify lifetime value without further
information. Distinguish observed-period contribution from modeled
continuation. Discount curves and funds-transfer prices are inputs with
provenance and sensitivity ranges.

\section{Corruption and imputation uncertainty}
Missingness and measurement errors often depend on customer behavior:
unobserved competitor balances, selectively requested income documents,
unlogged calls and sold loans are examples. Generic predictive imputation
does not establish causal recovery. Corruption-aware theory can use
approximate repeated measurements and low-rank structure
\citep{corrupt2022}; those restrictions must be justified for each dataset.

The integration with a data-recovery system should transmit masks,
measurement models, uncertainty draws and recovery assumptions. Imputation
must occur inside the appropriate training partitions when fitted
preprocessing could leak evaluation information. Outcomes may appropriately
enter a congenial statistical imputation model for some estimands, but such
use must never make unavailable outcomes appear as operational policy inputs.
The two roles require separate datasets and explicit contracts.

Combining estimates across recovery draws needs within- and between-draw
uncertainty. A simple multiple-imputation formula is appropriate only under
its assumptions and analysis compatibility. Joint models or sensitivity sets
may be necessary. A single completed table with no missingness record
conceals a material source of uncertainty.

\section{Reuse the bank's data foundations}
DataHub and OpenMetadata provide catalog, ownership and lineage capabilities
\citep{sw_datahub,sw_openmetadata}. OpenLineage records runtime job and dataset
relationships \citep{sw_openlineage}. Feast provides historical/online feature
access and time-travel retrieval \citep{sw_feast}; the inspected API includes
a created-timestamp filtering option. Availability semantics still require
verification of the chosen backend and source timestamps.
Pandera can enforce dataframe contracts \citep{sw_pandera}.

The bank should integrate its existing catalog and feature infrastructure.
The causal product adds study-specific semantics: intervention versions,
decision information, estimands, dependence, graph assumptions and evidence.
It should not require the group to replace a working enterprise catalog.

\section{Data architecture as accountable evidence}
BCBS239 emphasizes integrated taxonomies, identifiers, ownership and robust
reconciliation, without requiring one physical data model
\citep{bcbs239}. These principles support the proposed semantic layer,
but a causal product's conformity to them is not established by this proposal.
Production acceptance needs actual lineage completeness, reconciliation and
recovery tests on the bank's systems.

Useful automation therefore makes a reviewable proposal, detects
inconsistencies and produces reproducible transformations. Consequential
semantic assertions remain attributable to a responsible owner. The user
interface should show the unresolved business question in ordinary language,
such as whether the recorded offer denotes assignment or acceptance, rather
than exposing irrelevant implementation internals.

\section{A canonical model that preserves native accounts}
The common data model should standardize identities and relationships while
leaving native financial objects intact. A household can own several accounts;
a company can be borrower on a facility and guarantor on another; a booking
entity can own a claim serviced elsewhere. Collapsing these roles into one
``customer ID'' creates ambiguity in assignment, interference and financial
ownership. The entity graph needs typed roles with validity intervals and
the source of each link.

At the center is a decision episode: actor, subject, eligible actions,
decision time, information snapshot and policy version. Assignment,
dispatch, receipt, acceptance and settlement are subsequent linked events.
A failed dispatch can incur cost without receipt; acceptance can occur after
several contacts; settlement can be corrected later. A single binary
``treated'' field cannot faithfully represent all these stages.

Outcome components refer back to the original episode and economic owner.
They carry currency, sign convention, economic date, available date and
claim or account identifier. The same customer transfer can generate two
product entries but only one group movement. Reconciliation keys allow the
group account to cancel internal entries without erasing their product-
level operational meaning.

This model should support views for customer trajectories, branch panels,
bank-firm relationships and aggregate series. A view declares its
aggregation weights and denominator. An average per active account differs
from a total per original eligible household, particularly when the policy
changes closure or new account opening. The platform must retain the
denominator as a first-class object.

\section{Revision selection precedes feature-window selection}
For one decision, first identify versions that were available by its
timestamp. Among those, select the authoritative revision of each economic
event. Only then apply the feature's economic-time window and void status.
Reversing those steps can resurrect an obsolete event.

The supplied temporal component exercises four records. A credit of
USD\,100 becomes available on 3 January and is revised to USD\,60 on
10 January. A USD\,30 event dated 4 January arrives on 12 January.
A USD\,20 event initially dated 5 January is revised on 10 January to
an economic date of 1 February. A USD\,40 event is voided on 10 January.
For a complete-feed thirty-day lookback, the correct totals are
USD\,160 at the 8 January decision and USD\,90 at the 13 January decision.

Using the final revised records while ignoring availability gives USD\,90
for the early decision and leaks future knowledge. Filtering economic
time before selecting revisions can retain the obsolete USD\,20 version
in the later decision, giving USD\,110. These errors survive a superficially
correct event-time cutoff. The code preserves selected source versions,
rejects duplicate revisions and checks the calculation against the small
account stated here.

A warehouse implementation can express the same order:
\begin{lstlisting}[language=SQL]
WITH known_versions AS (
  SELECT d.decision_id, d.decision_at, e.*,
         ROW_NUMBER() OVER (
           PARTITION BY d.decision_id, e.event_id
           ORDER BY e.revision DESC) AS revision_rank
  FROM decisions d JOIN events e
    ON d.entity_id = e.entity_id
   AND e.available_at <= d.decision_at
), selected AS (
  SELECT * FROM known_versions WHERE revision_rank = 1
)
SELECT decision_id, SUM(amount) AS amount
FROM selected
WHERE NOT void
  AND occurred_at <= decision_at
  AND occurred_at > decision_at - INTERVAL '30 days'
GROUP BY decision_id;
\end{lstlisting}
This query illustrates the semantics. A production query additionally
preserves decisions with no selected events, enforces source completeness,
resolves currencies and permissions, and uses a declared authoritative
revision order. The executable reference uses Python for inspectability;
the SQL is a proposed warehouse translation, not a claimed benchmark.

\section{Absence is not automatically a zero}
No matching event can mean no event occurred, the feed is incomplete,
the customer was outside coverage, or the relevant identifier was unresolved.
Only the first interpretation supports an observed zero. The reference
temporal component requires a caller declaration that the feed covers the
requested window before returning an additive total. It cannot independently
verify that business assertion.

For the full platform, observation status should distinguish at least
observed value, observed zero, not applicable, not yet available,
missing within coverage and outside coverage. Trajectories retain these
states together with action eligibility and outcome maturity. An account
that closes should not disappear from the original cohort simply because
its subsequent product balance is no longer stored in the active-account
table.

Measurement models can then use the appropriate pattern. A missing income
document requested only from risky applicants is not MCAR. A competitor
balance observed only for customers consenting to aggregation can require
selection assumptions. A masked field supplied by dataml should retain
its recovery model and uncertainty draws. The causal consumer decides
whether the recovered quantity is an adjustment variable, outcome or
scenario input and propagates the corresponding uncertainty.

\section{What can be automated reliably}
Schema and metadata profiling can detect candidate date columns, units,
key uniqueness, foreign-key coverage, value ranges and stock-flow violations.
String and representation matching can suggest glossary terms. A record
linker can propose probabilistic entity links. A transformation generator
can draft a join and its validation queries. These are useful assistance
tasks because their output is concrete and reviewable.

The product should measure mapping quality by field role and consequence.
Confusing gross and net balance can alter valuation; confusing assignment
and acceptance can alter identification. A single global mapping-accuracy
percentage hides that difference. A suitable evaluation uses withheld
schemas from another source system, measures precision and recall for
critical roles, and records the correction effort of a qualified data owner.
Abstention on ambiguous mappings can be preferable to a confidently wrong
automatic link.

An automatic organizer can also propose temporal precedence constraints:
an offer assignment precedes its receipt, and a settlement follows its
underlying transaction. It should attach the source rule and exceptions.
These constraints can reduce discovery search space without asserting that
business chronology uniquely determines causality. Semantic similarity
between ``deposit rate'' and ``funding cost'' is not evidence of a causal
edge or an exclusion restriction.

\section{Entity uncertainty propagates into the study}
Suppose linkage produces plausible entity assignments
$\mathcal L_1,\ldots,\mathcal L_M$ with weights $w_m$. A scientific
analysis can repeat the study under each linkage and report how effects,
fold membership and financial totals change. Interpreting
$\sum_m w_m\widehat\theta_m$ as a posterior mean requires a
defensible joint linkage and outcome model; treating arbitrary match scores
as posterior probabilities is not justified.

False merges can combine one person's treatment with another's outcome.
Missed links can split a household between training and evaluation folds.
Both can produce bias beyond ordinary missing-data variance. A useful
qualification includes controlled false merges, missed links and changes
in connected-component size. The report should reveal whether a result
depends on a few uncertain high-value links.

Deterministic authoritative identifiers are preferable where available,
but they also change through mergers, account migrations and legal
restructuring. The entity service must preserve historical identity and
effective relationships rather than rewriting all past records to the
current organization chart. For a bank-firm credit-supply study, current
group consolidation can otherwise leak information or change the unit
of analysis retrospectively.

\section{Physical execution and access boundaries}
Canonical semantics need not imply a central copy of every transaction.
The platform can publish approved views over existing warehouses, use
columnar snapshots for bounded studies and execute locally within legal
entities. A catalogue identifies dataset owner, schema, permitted purpose,
retention and lineage. A study records the authorized snapshot reference
and columns used, with access checks enforced by the deployment environment.

Large transformations should push filters and projections toward storage,
partition by useful entity/date keys and avoid materializing every lag
for every possible study. Temporal feature caches need keys that include
source revision, availability rule, entity mapping, window, currency and
transformation version. A cache keyed only by feature name and event date
can return a value unavailable at the decision.

Derived features also have publication latency. Even if all raw inputs
existed by noon, a feature generated in an overnight job was not available
to a noon decision unless that computation could actually have run then.
Retrospective scientific adjustment and operational policy information are
different uses. Their admissible data views should be explicit.

\section{Acceptance is a reconstructed account}
The strongest data test is a small account reconstructed from independent
source entries, including late arrivals, reversals, ownership changes and
internal transfers. The resulting feature and outcome totals should be
explained in business terms. Passing a dataframe schema alone cannot
establish that the chosen records describe the intended economic object.

The reference component returns selected event revisions and a content hash,
demonstrating this principle for a narrow additive feature. It is not a
production ingestion engine or automatic semantic classifier. The proposed
next implementation adds warehouse adapters, completeness attestations,
entity-version references and outcome-component reconciliation while
retaining that reconstructible small-account test.

\chapter{Open-source software and build decisions}\label{ch:software}
\section{What was assessed}
The census covers \SoftwareCount{} projects as observed on 7 October 2026. Saved records
include official repository metadata, primary documentation and licenses,
available package metadata, selected API source and content hashes.
The matrix is a source assessment, not an executed comparative benchmark.
An observed last push is not a maintenance guarantee, and monorepo activity
need not describe a particular submodule.

The accompanying JSON records retain package versions and release dates
where retrieved. Some GitHub requests were rate-limited; CausalNex's
repository endpoint was unavailable. A missing response is not evidence
that a repository is archived. Where a commit identifier was retrieved,
it is recorded separately from mutable-branch document hashes rather than
asserted to pin files that were fetched independently.

\section{Estimation and econometrics}
DoWhy's useful contribution is separation of modeling, identification,
estimation and refutation \citep{dowhy2020,sw_dowhy}.
The inspected \texttt{identify\_effect} source accepts an identification
method and returns an estimand object. The proposed wrapper must preserve
that object and unresolved identification status through fitting.
Refutations stress assumptions or estimators; passing them is not proof.

EconML offers heterogeneous-effect estimators, DML, forests, IV and
panel DynamicDML \citep{sw_econml}. The inspected dynamic implementation
uses group/period structure and recursively adjusts future effects.
The inspected implementation requires equal trajectory lengths and treats
within-group occurrence order as chronological; its splitting keeps groups
intact. The adapter must validate these restrictions before calling it,
including the treatment of irregular visits and incomplete histories.
DoubleML provides transparent orthogonal-score estimation and common
inference operations \citep{doubleml2022,sw_doubleml}. Its documentation
supports model-specific nuisance estimation and cluster-aware data structures;
the attempted direct cluster source path returned 404, so no claim here
rests on inspecting that missing file.

CausalML is a useful marketing challenger, while \texttt{grf} provides
honest forest-based estimation, including causal and instrumental settings
\citep{sw_causalml,sw_grf}. They should be compared on the same population,
target and evaluation design. A package's uplift terminology is not
evidence that historical assignment is unconfounded.

For econometric baselines, linearmodels provides panel, IV and system
regression; PyFixest and R fixest provide efficient fixed-effects estimation
\citep{sw_linearmodels,sw_pyfixest,sw_fixest}. Current PyFixest documentation
also lists event-study and robust-covariance capabilities. Faster fixed
effects do not repair a misspecified heterogeneous-treatment target.
Do not infer dynamic-panel GMM support simply from the word panel.
The dedicated \texttt{pydynpd} package documents difference and system GMM,
corrected two-step errors, serial-correlation and overidentification tests,
and collapsed instruments \citep{sw_pydynpd}. It is a more appropriate
comparator for that task. Automated model search must not be interpreted
as validation of moment restrictions.

Use R \texttt{did} as a reference for group-time ATT, HonestDiD for
trend sensitivity, synthdid for its supported synthetic-DiD design,
and gsynth for interactive-factor comparisons
\citep{sw_did,sw_honestdid,sw_synthdid,sw_gsynth}.
CausalImpact and CausalPy provide aggregate time-series and Bayesian
quasi-experimental alternatives \citep{sw_causalimpact,sw_causalpy}.
Retain R adapters where scientific comparators are strongest rather than
requiring premature Python reimplementation.

Longitudinal estimation also has established implementations.
\texttt{lmtp} supports modified treatment policies, targeted estimation,
sequential DR and competing-risk outcomes; \texttt{ltmle} implements
treatment/censoring-specific means and marginal structural models
\citep{sw_lmtp,sw_ltmle}. \texttt{gfoRmula} supplies parametric
g-computation for time-varying treatments and several outcome types
\citep{sw_gformula}. Python \texttt{zEpid} provides weighting,
g-formula, AIPW and TMLE tools \citep{sw_zepid}.
These should be scientific comparators before internal reimplementation.
Their supported interventions, node ordering, censoring behavior and
uncertainty need adapter-specific qualification. zEpid's observed last
repository push was May 2023; compatibility is untested here.

\section{Discovery and latent decisions}
Tigramite is the primary temporal-discovery candidate: its inspected
\texttt{run\_pcmciplus} API exposes lag limits, link assumptions,
conditioning controls and multiplicity options \citep{sw_tigramite}.
Causal-learn offers constraint, score and functional methods
\citep{causallearn2024,sw_causallearn}. Tetrad supplies an independent Java
ecosystem for comparison \citep{sw_tetrad}. Causica's joint graph and
structural-model learning, gCastle, DAGMA and NOTEARS are research
comparators \citep{sw_causica,sw_gcastle,sw_dagma,sw_notears}.
Their fitted graph distributions or exact acyclicity constraints do not
remove causal assumptions.

d3rlpy offers offline policy-learning implementations; SCOPE-RL offers
evaluation and policy selection; OBP addresses contextual-bandit evaluation
\citep{sw_d3rlpy,sw_scope_rl,sw_obp}. None should be advertised as a
generic solution to hidden-confounded bank histories. SCOPE-RL and OBP
had observed last repository pushes in March and June 2024 respectively;
dependency qualification and support planning matter.

\texttt{pomdp\_py} and POMDPs.jl provide modeling/planning ecosystems
\citep{sw_pomdp_py,sw_pomdps}. Their role is distinct from parameter
identification. Dynamax provides state-space inference and learning,
while PyMC and NumPyro support probabilistic model estimation
\citep{sw_dynamax,sw_pymc,sw_numpyro}. These are appropriate reference
engines for a controlled latent-model benchmark. Existing BayesFilter
components may also be reusable, subject to the integration qualification.

\section{Data and execution}
DuckDB is a strong local reference engine for SQL transformations and
auditable small-to-medium studies; Polars provides columnar dataframe
execution; Ray can distribute independent workloads
\citep{sw_duckdb,sw_polars,sw_ray}. Their selection should follow profiling,
not an assumption that every study needs a cluster.
Catalog, lineage, feature and validation tools supply the foundations
described in Chapter~\ref{ch:data}.

\section{Licensing and dependency economics}
The inspected EconML license is MIT and CausalML is Apache-2.0 despite
ambiguous repository API license fields. DoubleML is BSD-3-Clause and
linearmodels uses NCSA. PyFixest is MIT. R fixest and did declare GPL-3;
Tigramite, grf and Tetrad have GPL-3 licensing. HonestDiD and gsynth
declare MIT; synthdid is BSD-3-Clause.
\texttt{lmtp}'s DESCRIPTION and actual license text specify AGPL-3,
despite a GPL-3 badge in its README. Use the license text when qualifying
the dependency. \texttt{ltmle} and \texttt{gfoRmula} declare GPL-2 and
GPL-3 respectively; zEpid and pydynpd use MIT.

These observations affect distribution, integration and procurement.
A separate process or service does not automatically settle copyleft
obligations. The product's distribution model, modifications, dependency
tree and license texts require legal review at productization. No package
should be excluded solely because a metadata service failed to classify
its license, nor assumed permissive without the actual terms.

\section{Choose by a scientific task, not a package ranking}
The relevant unit of comparison is a complete estimand and design. Consider
four requests that an interface might label as a treatment effect: the
population ATE of an assigned offer, a variance-weighted partially linear
coefficient, a group-time ATT following a branch rollout, and a value
contrast between eight-decision policies. These objects have different
populations and weighting. A wrapper that accepts the first available
\texttt{effect} method and returns its array can silently substitute one
for another.

For a static confounded offer, DoWhy can carry an explicit graph and
identification record, DoubleML can expose a declared orthogonal score,
EconML can estimate heterogeneity under the supported model, and grf can
provide a forest-based challenger. They play complementary roles
\citep{sw_dowhy,sw_doubleml,sw_econml,sw_grf}. The bank-specific layer
still owns treatment versions, eligible population, cost conversion and
the covariance of the final policy contrast. Its evaluation should compare
the same nuisance information and held-out units. Giving a flexible method
more features or future observations would confound the software comparison.

For the deposit rollout in Chapter~\ref{ch:deposits}, R did is the principal
group-time reference, while fixed-effects packages reproduce an informative
but potentially inappropriate TWFE baseline. HonestDiD investigates
departures from trend restrictions; an imputation estimator offers another
aggregation route \citep{sw_did,sw_honestdid,borusyak2024}.
These are not votes to average indiscriminately. Differences should be
explained by comparison groups, target weights, support, functional
restrictions and inference. A faster absorbed regression does not resolve
which post-treatment cells belong in the causal target.

For repeated assistance, gfoRmula, ltmle and lmtp provide distinct
constructions \citep{sw_gformula,sw_ltmle,sw_lmtp}. A forward simulator
requires the conditional evolution of time-varying measurements; a targeted
or sequential-regression method can require continuation regressions
without a complete generative environment. Modified-treatment-policy
support in lmtp is particularly relevant for continuous dose changes,
but does not make an arbitrary contractual restructuring a supported shift.
The wrapper must inspect the intervention definition and loss/censoring
nodes before translating a banking request.

The controlled latent-state study needs a different composition.
Dynamax or a probabilistic programming engine can supply estimation
infrastructure; a POMDP planning ecosystem can consume the fitted
environment \citep{sw_dynamax,sw_pymc,sw_numpyro,sw_pomdps,sw_pomdp_py}.
The bridge between them must preserve action selection, reward timing,
policy information and parameter uncertainty. A planning demo with known
transitions does not establish that its upstream banking likelihood is
correct or causally identified.

\section{An adapter qualification is a scientific translation}
Each adapter should publish a capability declaration: supported estimand,
intervention type, outcome/censoring model, panel shape, fold interface,
sample weights, cluster design, uncertainty method and software version.
Unsupported combinations should fail before fitting. For example, a
trajectory adapter cannot quietly sort by the current account identifier
if the method interprets input order as calendar order.

Qualification begins with exact or independently reconstructed cases.
For a group-time estimator, verify the same cohort/time cells and
aggregation weights. For sequential evaluation, verify terminal rewards,
discounting and cumulative ratios against short-path enumeration. For an
HMM, compare marginalized likelihood and expected counts. For discovery,
verify endpoint marks and lag conventions through serialization. A round
trip that turns a circle endpoint into an arrow is a substantive defect
even if the adjacency matrix is unchanged.

Next compare supported synthetic mechanisms over repeated samples, with
appropriate bias, coverage and failure metrics. A fixed numerical tolerance
is suitable for an exact linear calculation but generally inappropriate
for two stochastic nonconvex fits. Their objective, initialization,
termination, label alignment and decision-value differences must be
examined. Package reference agreement is necessary evidence of translation;
it does not prove the maintained model describes the bank.

Finally, record operational behavior: memory, serialization, random-state
control, missing-value conventions, interrupts and upgrade differences.
The software bill of materials and environment lock must match the
executed adapter. A license snapshot and repository HEAD retrieved on the
same day do not establish that independently fetched mutable files came
from that exact commit. A qualified build must retrieve its source at a
specific immutable revision.

\section{Where a shared library should build its own code}
The bank has reasons to own temporal reconstruction and financial
reconciliation: these operations depend on its action logs, legal ownership,
revision systems and funds-transfer pricing. It also has reasons to own
score contracts and inference orchestration, because those connect
estimators to the original population and common shocks. The supplied
components in Chapter~\ref{ch:construction} make two such boundaries
concrete without claiming a complete platform.

By contrast, reimplementing every forest, optimizer, panel estimator and
graph search would create a large validation and maintenance obligation
with limited differentiation. Thin adapters allow established algorithms
to remain identifiable. The scientific layer should export intermediate
quantities needed for independent checks, rather than hiding all methods
behind a proprietary success score.

There are cases for internal reference implementations. The finite-state
controlled HMM in this proposal is small enough that its likelihood and EM
updates can be audited directly. It is a useful comparator for a future
accelerated engine. A compact orthogonal score implementation can verify
normalization and financial contrasts. Such references should remain
deliberately simple, with an explicit supported domain, rather than grow
unnoticed into a second general-purpose library.

\section{Language boundaries and deployment choices}
Python is a practical orchestration language for this workspace, but it
does not determine the best scientific implementation. R has mature
panel and longitudinal comparators; Java supplies the Tetrad ecosystem;
Julia supplies a substantial POMDP interface \citep{sw_did,sw_ltmle,
sw_tetrad,sw_pomdps}. A versioned process boundary can exchange immutable
columnar data and a result manifest without changing the estimand.
The method's license, runtime support and error conventions remain part
of qualification.

An in-process adapter has lower transfer overhead and easier interactive
use, but couples dependency resolution and failure modes. A worker process
can isolate numerical and language environments, at the cost of data
serialization, deployment complexity and duplicated memory. A service
boundary adds scheduling and access controls but does not automatically
solve licensing or scientific compatibility. Choose after measuring the
actual workload and reviewing the distribution arrangement.

For large histories, DuckDB or the bank's warehouse should construct
partitioned snapshots; Polars can support local transformations; Ray can
schedule sufficiently independent fits \citep{sw_duckdb,sw_polars,sw_ray}.
The independent units, common date shocks and folds are declared before
that scheduling decision. Moving row partitions among workers must not
redefine a cross-fitting scheme.

\section{The complete decision matrix}
\small
\begin{longtable}{P{29mm}P{55mm}P{60mm}}
\toprule Project; license & Documented role and observation & Recommendation and boundary\\\midrule\endhead
dowhy \citep{sw_dowhy}\par MIT & Model, identify, estimate and refute; graphical causal mechanisms. Last push: 2026-10-06. Observed package declaration: 0.14. & Wrap identification and refutation; require explicit estimand and graph version. Refutation failure is informative; passing refutations does not verify unmeasured-confounding assumptions\\
EconML \citep{sw_econml}\par MIT & DML, forests, heterogeneous effects and IV; panel DynamicDML. Last push: 2026-10-06. Observed package declaration: 0.17.0. & Primary Python heterogeneous-effect adapter. Estimator-specific exchangeability, IV and panel restrictions; preserve groups in fitting\\
doubleml-for-py \citep{sw_doubleml}\par BSD-3-Clause & Orthogonal scores, cross-fitting, inference and cluster data. Last push: 2026-09-23. Observed package declaration: 0.11.4. & Primary transparent score/inference adapter. A supported cluster layout is not arbitrary temporal dependence\\
causalml \citep{sw_causalml}\par Apache-2.0 & Meta-learners, trees and uplift evaluation. Last push: 2026-10-05. Observed package declaration: 0.17.0. & Marketing challenger. Uplift fit and predictive ranking alone do not establish identification\\
grf \citep{sw_grf}\par GPL-3.0 & Honest causal, instrumental, survival and multi-arm forests. Last push: 2026-04-30. & R benchmark and optional governed adapter. Copyleft review and dependence assumptions precede distribution\\
tigramite \citep{sw_tigramite}\par GPL-3.0 & PCMCI, PCMCI+, LPCMCI, regime and multi-dataset variants. Last push: 2026-01-14. Observed package declaration: 5.2.10.1. & Primary temporal discovery research adapter. Stationarity, graph semantics, CI-test assumptions and copyleft require explicit review\\
causal-learn \citep{sw_causallearn}\par MIT & Constraint, score and functional discovery plus CI tests. Last push: 2026-09-04. Observed package declaration: 0.1.4.8. & Primary static discovery adapter. CPDAG/PAG endpoints cannot be silently converted to a DAG\\
causalnex \citep{sw_causalnex}\par Unverified in this audit & PyPI package metadata available; official GitHub endpoint unavailable. Last push: unavailabl. Observed package declaration: 0.12.2. & Defer new dependency. 404/rate-limited access does not establish archive status or current support\\
causica \citep{sw_causica}\par MIT & DECI joint graph and additive-noise structural-model inference. Last push: 2026-10-06. & Research comparator. Joint discovery/estimation does not remove structural assumptions\\
trustworthyAI \citep{sw_gcastle}\par Apache-2.0 & Causal discovery toolbox within trustworthyAI monorepo. Last push: 2026-09-21. & Research comparator. Monorepo activity is not evidence of submodule maintenance\\
dagma \citep{sw_dagma}\par Apache-2.0 & Linear/nonlinear DAG search with log-determinant acyclicity. Last push: 2024-01-14. & Small benchmark comparator. Nonconvex optimization, causal sufficiency and functional assumptions remain\\
notears \citep{sw_notears}\par Apache-2.0 & Continuous optimization under exact acyclicity constraint. Last push: 2026-05-26. & Reference algorithm. Acyclicity is not causal identification; inspect numerical thresholding\\
tetrad \citep{sw_tetrad}\par GPL-3.0 & Java causal discovery library and GUI. Last push: 2026-10-06. & Cross-language benchmark. Deployment and licensing cost; graph semantics must survive serialization\\
CausalImpact \citep{sw_causalimpact}\par Apache-2.0 & Bayesian structural time-series intervention analysis. Last push: 2026-03-31. & Aggregate campaign comparator. Untreated covariates and stable counterfactual relationship are essential\\
CausalPy \citep{sw_causalpy}\par Apache-2.0 & DiD, synthetic controls, RD, interrupted series and IV. Last push: 2026-10-06. Observed package declaration: 0.9.0. & Bayesian challenger. Posterior uncertainty is conditional on design and model assumptions\\
linearmodels \citep{sw_linearmodels}\par NCSA & Panel, IV, system regression and covariance estimators. Last push: 2026-10-06. Observed package declaration: 7.0. & Econometric baseline adapter. Do not infer dynamic-panel GMM support from the panel-regression label\\
pyfixest \citep{sw_pyfixest}\par MIT & High-dimensional fixed effects, IV, event studies and robust covariance. Last push: 2026-10-07. Observed package declaration: 0.60.0. & Primary Python panel adapter. Fast TWFE remains inappropriate for some heterogeneous staggered effects\\
fixest \citep{sw_fixest}\par GPL-3 & Multiway fixed effects, IV, GLM and inference. Last push: 2026-06-26. Observed package declaration: 0.14.2. & R reference adapter. GPL review; specify the estimand before the regression formula\\
did \citep{sw_did}\par GPL-3 & Group-time ATT with att\_gt and aggregation. Last push: 2026-08-28. Observed package declaration: 2.5.1.902. & R reference for modern DiD. No anticipation, group-specific parallel trends and comparison support\\
HonestDiD \citep{sw_honestdid}\par MIT & Robust confidence sets under restricted trend violations. Last push: 2026-04-12. Observed package declaration: 0.2.8. & Sensitivity adapter. Sensitivity restriction is substantive, not learned truth\\
synthdid \citep{sw_synthdid}\par BSD-3-Clause & Synthetic difference-in-differences estimator. Last push: 2024-01-15. & Optional comparator. Documented block treatment/simultaneous adoption and beta API\\
gsynth \citep{sw_gsynth}\par MIT & Counterfactual imputation using interactive fixed effects. Last push: 2026-09-30. Observed package declaration: 1.4.0. & Optional factor-model comparator. Latent-factor extrapolation and support require validation\\
d3rlpy \citep{sw_d3rlpy}\par MIT & Offline/online RL and distributed training. Last push: 2025-09-10. Observed package declaration: 2.8.1. & Policy challenger after evaluation design. Logged support and causal information set are separate from learner capability\\
scope-rl \citep{sw_scope_rl}\par Apache-2.0 & OPE, policy selection and distributional evaluation. Last push: 2024-03-18. Observed package declaration: 0.2.1. & OPE comparator. Snapshot last push in 2024; no generic hidden-confounding guarantee\\
zr-obp \citep{sw_obp}\par Apache-2.0 & Bandit OPE/learning and logged bandit data tools. Last push: 2024-06-03. & One-step policy comparator. Contextual bandits do not represent multi-period banking carryover\\
pomdp-py \citep{sw_pomdp_py}\par MIT & POMDP modeling, belief updating, planning and simulation. Last push: 2025-10-28. Observed package declaration: 1.3.5.1. & Small reference planner. Planning interface is not a causal parameter estimator\\
POMDPs.jl \citep{sw_pomdps}\par MIT & Julia model interface, tools, solvers and simulators. Last push: 2026-10-04. & Julia reference planner. Solver ecosystem versions/licenses must be pinned separately\\
pymc \citep{sw_pymc}\par Apache-2.0 & MCMC and variational Bayesian modeling. Last push: 2026-10-06. & Hierarchical-model comparator. Convergence and posterior fit do not identify interventions\\
numpyro \citep{sw_numpyro}\par Apache-2.0 & JAX probabilistic programming and HMC. Last push: 2026-10-05. Observed package declaration: 0.22.0. & Accelerated inference adapter candidate. Discrete states often need marginalization; hardware speed unmeasured here\\
dynamax \citep{sw_dynamax}\par MIT & JAX state-space filtering and parameter learning. Last push: 2026-10-01. Observed package declaration: 1.0.2. & First latent-dynamics comparator. Statistical state models need separate controlled causal identification\\
duckdb \citep{sw_duckdb}\par MIT & Embedded SQL analytics with window/nested queries. Last push: 2026-10-07. Observed package declaration: 1.5.6. & Default local reference execution. Measure spill, memory and point-in-time correctness on bank workloads\\
polars \citep{sw_polars}\par MIT & Lazy/streaming analytical DataFrames; optional GPU. Last push: 2026-10-07. & Optional columnar adapter. Capability does not establish distributed causal estimation\\
ray \citep{sw_ray}\par Apache-2.0 & Distributed task execution and ML tools. Last push: 2026-10-07. & Add only when single-node profiling warrants. Fault recovery, deterministic folds and aggregation are product responsibilities\\
datahub \citep{sw_datahub}\par Apache-2.0 & Metadata graph, ownership, column lineage and discovery. Last push: 2026-10-07. & Integrate the bank's existing catalog. Semantic context and lineage are not causal knowledge\\
OpenMetadata \citep{sw_openmetadata}\par Apache-2.0 & Metadata, data quality, lineage, ownership and contracts. Last push: 2026-10-07. & Alternative catalog integration. Avoid operating a duplicate enterprise catalog\\
OpenLineage \citep{sw_openlineage}\par Apache-2.0 & Open standard for runtime job/dataset lineage. Last push: 2026-10-06. & Adopt run and dataset lineage events. Add causal estimands and assumptions as explicit custom facets\\
feast \citep{sw_feast}\par Apache-2.0 & Offline/online feature access and point-in-time retrieval. Last push: 2026-10-07. & Integrate existing feature store. Event-time joins must additionally enforce actual availability time\\
splink \citep{sw_splink}\par MIT & Probabilistic record linkage and match diagnostics. Last push: 2026-10-06. Observed package declaration: 5.0.0. & Optional entity-resolution adapter. Propagate uncertain links; false merges can induce causal bias\\
pandera \citep{sw_pandera}\par MIT & Typed dataframe validation across engines. Last push: 2026-10-07. Observed package declaration: 0.34.1. & Contract-validation adapter. Schema validity does not establish causal semantics\\
lmtp \citep{sw_lmtp}\par AGPL-3.0 & Longitudinal modified treatment policies, TMLE and sequential DR; competing risks. Last push: 2026-07-01. Observed package declaration: 1.5.4. & Qualified R longitudinal comparator. DESCRIPTION and license text specify AGPL-3, despite GPL-3 README badge; support and causal restrictions remain\\
ltmle \citep{sw_ltmle}\par GPL-2 & Longitudinal TMLE of treatment/censoring-specific means and marginal structural models. Last push: 2026-06-03. Observed package declaration: 1.3-0. & R reference implementation. Censoring, intervention support and chronological node semantics need explicit checks\\
gfoRmula \citep{sw_gformula}\par GPL-3 (DESCRIPTION) & Parametric g-formula with time-varying treatments, failure times and competing events. Last push: 2025-11-03. Observed package declaration: 1.1.2. & R g-computation comparator. Model specification and intervention semantics are substantive; full license text not separately retrieved\\
zEpid \citep{sw_zepid}\par MIT & IPTW/IPCW, AIPW, g-formula, TMLE and transportability tools. Last push: 2023-05-07. Observed package declaration: 0.9.1. & Python g-method reference comparator. Observed last push May 2023; qualify dependency compatibility and exact longitudinal method scope\\
pydynpd \citep{sw_pydynpd}\par MIT & Difference/system GMM, corrected two-step errors, AR and Hansen diagnostics, collapsed instruments. Last push: 2025-12-15. Observed package declaration: 0.2.2. & Dynamic-panel GMM comparator. Moment restrictions and instrument proliferation are scientific choices; package model search cannot validate them\\
\bottomrule
\end{longtable}

\normalsize

\section{Recommended construction stack}
Begin all three scientific workspaces with a common temporal and evidence
layer. Use DuckDB and schema validation as local references and integrate
the bank's existing catalog. Qualify a small effect/panel path using
DoWhy where graph-based identification is appropriate,
DoubleML/EconML for supported effect tasks, PyFixest/linearmodels for
econometric baselines and selected R methods for direct scientific
comparison.

In the same first program, build the trajectory/score interface for
longitudinal methods; the finite-state latent-estimation interface with
an independently checked likelihood; and temporal discovery through
Tigramite with a causal-learn or Tetrad comparison. Their qualification
levels differ. Research estimation and ambiguous-graph output can be useful
before historical policy value or a unique causal structure is justified.
Neural challengers and distributed execution enter when measured application
requirements justify their additional cost.

The product must build the missing connections: decision-time semantics,
study compilation, evidence exports, financial outcome ownership,
dependence-aware orchestration, compatibility checks and bank-specific
workspaces. It should not attempt to build a new database, general enterprise
catalog, universal causal discovery theorem or complete RL ecosystem.

\chapter{Architecture, interfaces and repository reuse}\label{ch:architecture}
\section{One scientific core, several applications}
The architecture separates data access, study semantics, estimation and use.
The same immutable study definition should drive a notebook, a scheduled
service and the standalone workbench. Estimator adapters return typed
results; they do not directly approve customer actions. A result can be
computationally complete while scientifically unsupported.

\begin{figure}[htbp]\centering
\begin{tikzpicture}[node distance=8mm and 7mm,
 box/.style={draw=ink,rounded corners,fill=light,text width=42mm,align=center,minimum height=13mm,font=\small},
 flow/.style={-{Latex},thick,color=teal}]
\node[box] (data) {Bank systems\\catalog and event history};
\node[box,right=of data] (sem) {Temporal and entity layer\\snapshots, masks, lineage};
\node[box,right=of sem] (study) {Study compiler\\estimand, assumptions, folds};
\node[box,below=of study] (method) {Method adapters\\effects, panels, policies, graphs};
\node[box,left=of method] (evidence) {Evidence store\\scores, diagnostics, versions};
\node[box,left=of evidence] (apps) {Workbench and sibling apps\\decision reports and APIs};
\draw[flow](data)--(sem);\draw[flow](sem)--(study);
\draw[flow](study)--(method);\draw[flow](method)--(evidence);
\draw[flow](evidence)--(apps);
\end{tikzpicture}
\caption{Proposed architecture. Estimands, assumptions and lineage accompany results.}
\label{fig:architecture}
\end{figure}

The compiler checks required fields, intervention/population compatibility,
temporal availability, method support, dependence and action overlap.
Mechanical checks coexist with documented scientific judgments.
The execution layer schedules fold fits, scores, covariance aggregation
and diagnostics. The evidence store preserves data hashes, code/dependency
versions, seeds, configuration, assumption records and reports.
Large data and models remain in approved storage, referenced by authorized
metadata.

\section{Versioned public objects}
The supplied version 0.1.0 contracts define \emph{Study},
\emph{TemporalEvent}, \emph{EffectResult} and \emph{Graph}.
These are usable JSON Schemas with synthetic examples and selected semantic
checks. A separate version 0.2.0 research package implements the two narrow
components in Chapter~\ref{ch:construction}; the full SDK remains proposed.

Study records decision and assignment units, population/snapshot,
horizon, treatment/comparator, outcome, direction, currency, identification
route, assumptions and inference design. EffectResult preserves the
study hash, scientific status, evidence class, estimate or null,
interval semantics, independent-unit count, limitations and artifact
references. Confidence intervals, credible intervals and identified sets
must not share an unlabeled two-number field.

TemporalEvent distinguishes occurrence, recording, availability and
decision times. Graph preserves DAG/CPDAG/PAG type, lag and endpoint marks.
Checks reject late information, future historical measurements,
missing-as-zero records, unsupported point estimates, reversed intervals,
unlabeled bounds, absent assumptions, PAG edges called DAG edges,
unknown nodes and lags outside the window.

Further objects should include \emph{TrajectoryBatch},
\emph{RecoveryEnsemble}, \emph{PolicySpec}, \emph{BeliefModel},
\emph{OutcomeComponents} and \emph{EvidenceBundle}. Their required
semantics include action masks/propensities, observation/censoring masks,
state/reward timing, uncertainty draws, policy information, cash-flow
ownership and provenance. They are not already implemented schemas.

\section{Proposed API behavior}
A proposed Python workflow is:
\begin{lstlisting}
study = Study.from_json("study.json")
cohort = data.materialize(study.population, as_of=study.snapshot)
design = identify(study, cohort, graph=study.graph)
diagnostics = assess_support(design, cohort)
result = estimate(design, cohort, estimator="doubleml_adapter")
bundle = evidence.export(study, result, diagnostics)
\end{lstlisting}
These workflow names illustrate the proposed interface; they are not callable
functions in this repository. The implemented component names and actual
usage are given in Chapter~\ref{ch:construction}. Proposed service calls
return versioned artifacts and status rather than silently replacing
previous results.

\begin{center}\small
\begin{tabular}{P{76mm}P{68mm}}
\toprule Proposed endpoint & Behavior\\\midrule
\texttt{POST /v1/studies} & Validate a specification\\
\texttt{POST /v1/studies/\{id\}/runs} & Accept an idempotency key\\
\texttt{GET /v1/runs/\{id\}} & Return stage and scientific status\\
\texttt{GET /v1/runs/\{id\}/evidence} & Return authorized evidence references\\
\bottomrule
\end{tabular}
\end{center}
Data access uses approved dataset references rather than arbitrary SQL
executed with service credentials.

Errors should identify unknown treatment versions, unavailable features,
incompatible inference, unsupported actions, absent controls, unresolved
identification or numerical failure. Retries resume deterministic
checkpoints; changed inputs create a new run. Cancellation must not
publish a partial estimate as final.

\section{Cross-repository integration}
Documentation from nine local repositories was inspected and copied with
hashes \citep{localrepos}. The following are proposed interfaces,
not completed integrations.

\begin{longtable}{P{24mm}P{50mm}P{68mm}}
\toprule Repository & Proposed exchange & Constraint to preserve\\\midrule
\endhead
FinMarketing & Campaign, pricing, uplift and policy studies & Assignment versus acceptance, group contribution, horizon\\
dataml & Masks, recovery assumptions and uncertainty draws & Conditional recovery evidence; no completed table treated as error-free truth\\
BayesFilter & Filtering, smoothing and latent estimation & Event clock and numerical diagnostics; fit separate from causality\\
cardnpv & Outcome components and policy-value contrasts & Historical minus perfect-execution value; sign ID \texttt{V\_G0\_MINUS\_V\_GH}\\
SMEwallet & Offer/receipt, wallet/flow scenarios and evidence & Conservation, hidden-wallet identification limits and synthetic status\\
MacroFinance & Regime paths, structural shocks and sensitivities & Parameter identification, normalizations and shock assumptions\\
hponas & Nuisance tuning inside training partitions & No frozen causal evaluation outcomes in tuning\\
MathDev\-MCP & Equation/proof and numerical artifacts & Mathematics does not certify substantive assumptions\\
research-assistant & Literature identity and provenance & Paper organization differs from bank-data semantics\\
\bottomrule
\caption{Proposed integration map based on documentation inspection.}
\label{tab:repos}
\end{longtable}

SMEwallet's model card limits completed evidence to a synthetic prototype
and reports no promoted learned policy. Its scenarios are useful adverse
cases, not competitor-wallet ground truth. The inspected MacroFinance
milestone note discusses covariance provenance and conditional identification;
it is an earlier document, not a current overall status assessment.

For cardnpv, an adapter must map explicitly from any internal
treatment-minus-comparator convention to its canonical historical-minus-
perfect-execution direction. Test the sign on a known two-value fixture.
A numerically accurate sign reversal is a material financial error.

Integration begins with schema validation, then a synthetic adapter,
independent reference agreement and bank qualification. Repositories can
adopt contracts independently. Version mismatch must fail visibly rather
than discard fields.

\section{Deployment and operations}
Offer deployment inside the bank's controlled cloud or premises, plus local
research mode. A control service manages metadata, authorization and jobs;
workers run near data. Sensitive rows do not belong in application logs,
source-control repositories or paper-research artifacts.

Use enterprise identity, least-privilege access, encryption, audited
artifacts, retention policies and secrets management. These are proposed
requirements, to be validated in integration and security tests.
Legal-entity and regional separation applies to caches, features, models
and exports as well as source tables.

Pin adapter environments and software bills of materials. Retain exact
versions, licenses, vulnerability-review results and seeds for qualified
releases. Upgrade through reference, adverse-case and contract tests:
changed treatment coding, covariance defaults or graph endpoints can
alter a scientific conclusion.

Monitor data freshness, missingness, support, outcome delay, policy version,
graph stability and effect drift. Alarms should trigger recomputation,
restricted use or a redesigned study. They must not silently relax
identification assumptions.

\section{Model governance and current guidance}
Federal Reserve SR26-2 distributes April 2026 interagency guidance that
supersedes SR11-7 and SR21-8. It emphasizes risk-based governance,
conceptual soundness, outcomes analysis, effective challenge and vendor
responsibility \citep{fed2026}. Its supervisory applicability is specific;
it is not a global bank rule.
The scope footnote excludes generative and agentic AI models while
covering traditional quantitative and non-generative, non-agentic AI.
The semantic-assistance layer therefore needs the bank's applicable
governance assessment rather than an assertion that the guidance directly
covers every component.

PRA SS1/23 was revised in April 2026. It covers specified UK-incorporated
firms with internal-model regulatory-capital permission, with exclusions
stated in the document \citep{pra2026}. Its principles address model
lifecycle and self-assessment. Map each legal entity and use to applicable
requirements. The product supplies evidence; it cannot confer approval.

\chapter{Construction specification and measured components}\label{ch:construction}
\section{An operation must preserve the study's meaning}
The implementation should make the scientific argument reproducible across
execution environments. A scientist who requests an average effect on the
original eligible cohort should receive that target whether the job runs in
a notebook, a service worker or a distributed warehouse. Removing closed
accounts, changing fold boundaries or averaging branch coefficients with
different weights changes the analysis even if every local function succeeds.
The unit of composition is therefore a versioned study and its artifacts.

The proposal supplies two small working components: historical event
reconstruction and crossed-cluster score reduction. It also supplies fitted
research studies and four JSON contracts. These establish selected semantics
and calculations. The general study compiler, service endpoints, user
workspaces, distributed scheduler and sibling adapters described here remain
proposed construction. Keeping that boundary explicit lets an implementation
team identify its actual starting point.

\section{Package boundaries and owned invariants}
The proposed packages divide responsibility by the invariant they preserve.
The data package owns historical identity, observation status, availability
and reconciled economic components. The design package owns the estimand,
intervention versions, identification route, assumptions and target weights.
The execution package owns folds, model dependencies, resource limits and
retry semantics. Estimation adapters own the translation to an explicitly
supported external method. The inference package owns score meaning,
dependence, resampling and simultaneous contrasts. The evidence package owns
immutable reconstruction, source identity and permission-aware exports.

These packages should exchange documented arrays and records rather than
opaque framework objects. A trajectory artifact records trajectory key,
calendar key, decision index, action and eligibility mask, pre-action
information, outcomes, censoring and model-training role. An irregular
history cannot be padded into a balanced panel without recording which
positions are structural absence, unobserved visits or actual zeros.
An adapter requiring equal-length histories must reject or explicitly
transform an incompatible artifact.

\begin{longtable}{P{31mm}P{54mm}P{57mm}}
\toprule Operation & Required input & Returned artifact or refusal\\\midrule\endhead
Reconstruct snapshot & Decision episodes, event revisions, entity and coverage rules & Selected revisions, features, observation masks, reconciliation\\
Compile design & Study, snapshot, graph assumptions, target population & Estimand and identification status, fold plan, supported methods\\
Fit nuisance/model & Training partition, method/version, seed, resource limit & Model, training lineage, convergence and support diagnostics\\
Evaluate scores & Frozen models, evaluation partition, policy/target weights & Defined score vector, cluster keys, extreme-contribution diagnostics\\
Infer contrasts & Complete score partitions, covariance specification & Point estimate, labeled uncertainty, covariance and limitations\\
Export evidence & Completed run and authorized artifact references & Reconstruction manifest, scientific report and financial account\\
\bottomrule
\caption{Proposed operation contracts. Each refusal identifies a specific
missing condition; it is not a request to silently choose a new target.}
\label{tab:operations}
\end{longtable}

A model fit can finish while its transition parameters remain weakly
determined. A design can be valid while a particular fitted model reaches
its iteration cap. A job can be canceled before a complete covariance is
available. The result should therefore carry separate computational,
scientific and evidence states. Suggested computational states are queued,
running, complete, failed and canceled. Scientific states distinguish
identified under declared assumptions, partially identified, conditional
scenario and unresolved. Evidence states distinguish exact construction,
synthetic study, retrospective bank study and prospective evaluation.
One green success flag cannot communicate all three.

\section{The implemented temporal reference}
The small Python package is named \texttt{causalbank-research}, version 0.2.0.
Its \texttt{Event} and \texttt{Decision} records and
\texttt{compile\_snapshot} function implement the revision ordering in
Chapter~\ref{ch:data}. For example:
\begin{lstlisting}[language=Python]
from causalbank import Event, Decision, compile_snapshot

events = [
    Event("credit", "H1", "2026-01-02T00:00:00Z",
          "2026-01-03T00:00:00Z", 1, 100.0, "USD"),
    Event("credit", "H1", "2026-01-02T00:00:00Z",
          "2026-01-10T00:00:00Z", 2, 60.0, "USD"),
]
decisions = [Decision("d1", "H1", "2026-01-08T00:00:00Z")]
snapshot = compile_snapshot(
    events, decisions, lookback_days=30,
    currency="USD", feed_complete=True)
assert snapshot["rows"][0]["amount"] == 100.0
\end{lstlisting}
The result lists the selected event revision and a hash of the declared
input account. Reversing input event order preserves the result. A duplicate
event revision, changed event owner, nonfinite amount, incompatible currency
or timestamp without a timezone produces an error. The feed-completeness
declaration is required before absence becomes an additive zero.

This reference deliberately uses an inspectable per-decision algorithm.
It is not the implementation to scan a banking group's full transaction
history. A warehouse adapter should implement the same semantics with
indexed intervals, partition pruning and revision-aware joins, then pass the
same independent account fixture. Hashing a complete local input collection
also differs from production content addressing: a large job should use
immutable dataset manifests and partition digests instead of rereading every
byte at each operation.

Availability at exactly the recorded decision instant is treated as known
in the reference. A real decision system may require a strict predecessor
rule or a publication-latency margin. The source sequence number and the
decision system's clock precision determine that choice. It must be a
versioned convention; changing it creates a new snapshot.

\section{The implemented score reducer}
Let $x_j\in\mathbb R^d$ be an already computed, equally weighted score
vector. Write $n$ for its row count, $\bar x=n^{-1}\sum_jx_j$, and
$S_g=\sum_{j\in g}x_j$, $n_g=|g|$ for a declared cluster. The centered
cluster contribution is $U_g=S_g-n_g\bar x$. For $G$ nonempty clusters,
the component covariance used by the reference is
\[
 \widehat V_G=\frac{G}{G-1}\frac1{n^2}\sum_gU_gU_g'.
\]
Entity, date and entity--date intersections give
$\widehat V=\widehat V_E+\widehat V_T-\widehat V_{ET}$.
Retaining intersection groups matters when multiple scores occupy the same
entity/date cell. Replacing the last term by an ordinary row covariance
would then generally compute a different quantity.

Raw counts and sums are additive across partitions. The reducer merges
those sufficient statistics and centers once with the global mean.
Centering each partition and averaging its covariance generally loses
between-partition terms. The supplied check compares the merge result with
an independently centered batch calculation on a crossed 63-row fixture;
maximum discrepancies are below $2\times10^{-16}$ in this environment.

\begin{lstlisting}[language=Python]
from causalbank import ScoreReducer

left = ScoreReducer(n_entities=10000, n_dates=120, dimension=2)
right = ScoreReducer(n_entities=10000, n_dates=120, dimension=2)
# Each worker adds disjoint rows of precomputed scores and integer keys.
left.add(score_part_a, entity_part_a, date_part_a)
right.add(score_part_b, entity_part_b, date_part_b)
result = left.merge(right).result()
\end{lstlisting}
The caller owns score construction, target weighting and the sampling
assumptions. A raw regression coefficient is not automatically a score.
For a nonlinear estimating equation, its Jacobian and normalization must
be applied consistently. Unequal target weights need a defined weighted
score and denominator; they are not supplied by this prototype.

The covariance can be indefinite in a finite sample. The reference reports
its minimum eigenvalue and a flag instead of silently forcing a positive
matrix. A consumer must then use a scientifically justified alternative or
withhold the affected interval. The component also rejects nonfinite
scores, incompatible shapes and too few nonempty clusters. It does not
choose a bandwidth, bootstrap method or independent-unit definition for
the analyst.

\section{What the scale measurement actually times}
The scale fixture creates two-dimensional synthetic scores in chunks of
100,000 rows, with 10,000 declared entity keys and 120 date keys. Three
fresh child processes run at each size. The total timer includes score
generation, accumulation and final covariance; the reduction timer isolates
accumulation. Linux peak resident memory includes the interpreter and
loaded libraries.

\begin{table}[htbp]\centering\small
\begin{tabular}{@{}rrrr@{}}
\toprule Rows & Total seconds & Reduce seconds & Peak RSS (MiB)\\\midrule
100,000 & 0.097 & 0.078 & 77.1\\
1,000,000 & 0.233 & 0.133 & 115.2\\
10,000,000 & 1.186 & 0.630 & 123.1\\
\bottomrule\end{tabular}

\caption{Median of three local runs per size. Precomputed-score reduction
only: no storage reads, temporal reconstruction, nuisance fitting, causal
graph search or POMDP parameter fitting is included.}
\label{tab:scale-results}
\end{table}

The measured ten-million-row operation takes about 1.2 seconds and
123 MiB peak resident memory in this environment. These numbers show
that a bounded score reduction can avoid retaining all rows. They cannot
be advertised as the latency of a ten-million-row causal analysis.
The dense intersection arrays require
$O((E+T+ET)d)$ score storage plus counts. With ten million entities,
even 120 dates makes that representation inappropriate. A sparse grouped
or warehouse aggregation is then required and must be measured separately.
The current measurement does not establish a universal streaming-memory
bound independent of cluster cardinality.

The benchmark records Python and NumPy versions, platform and component
hashes. Shared-host load can change timing. There was no dedicated machine
reservation, storage benchmark or comparison with Ray. A production
qualification should specify cores, memory limits, execution concurrency,
data format, warm/cold caches and complete stage costs before committing
to a service objective.

\section{A reproducible job is a dependency graph}
A compiled job should reference a frozen snapshot and a deterministic fold
assignment. A nuisance fit depends only on its training role. Score tasks
depend on the exact model versions and their evaluation partitions.
Inference depends on the complete set of score partitions, with declared
target weights and covariance method. Report generation depends on the
financial account and both scientific and computational statuses.

The artifact identity should hash canonical configuration, input manifest
digests, code/dependency lock identity and random seeds. A model upgrade,
changed treatment mapping or later outcome snapshot changes that identity.
An authorization decision remains dynamic even when content is immutable:
a previously readable artifact may cease to be exportable after a permission
change. Content identity must not be confused with access authority.

Retries require an exactly-once contribution rule. A failed worker may
write the same score partition twice; the manifest must select one committed
artifact for that partition, rather than summing both. Committing the
partition manifest atomically is preferable to inferring completion from
the presence of a partially written file. A reducer starts only when all
expected partitions are present and valid.

Checkpointing is method-specific. An EM checkpoint includes model
parameters, iteration, objective and sufficient-statistic configuration.
A bootstrap checkpoint includes resample identity and fit diagnostics.
A graph-search checkpoint includes candidate constraints and conditioning
tests already performed. Resuming with changed data or constraints is a
new run. Saving an arbitrary Python process state is not the public
reconstruction contract.

\section{Three concrete repository adoption journeys}
FinMarketing can export assignment episodes, decision-time covariates and
a group contribution ledger. The proposed integration reconstructs the
snapshot, computes cohort effects and returns joint balance/contribution
uncertainty with an explicit target and sign. Acceptance compares the same
study with FinMarketing's existing calculation and reconciles transfers
account by account.

dataml can export a recovery ensemble with observed/missing masks, the
measurement model, draw identifiers and training partition provenance.
The causal consumer repeats the declared analysis under compatible draws
and reports recovery sensitivity. The initial integration should use an
exact small corrupted-data fixture whose observed entries remain unchanged.
It fails if a recovered outcome is exposed as a historical policy feature
or uncertainty is lost when the ensemble becomes a dataframe.

BayesFilter can supply a filtering or likelihood component through a narrow
numerical interface. A controlled model specifies whether the action factor
enters its historical posterior and when the reward is observed. Its
probabilities and gradients are compared with exact short-path sums before
larger fits. The causal library retains ownership of the intervention
interpretation. MacroFinance can similarly consume response-path covariance
and shock restrictions; cardnpv additionally requires an explicit contrast
direction conversion. These are proposed acceptance journeys, not existing
inter-repository calls.

\section{Building the independent application}
The standalone workbench should guide the bank user through subject and
action selection, eligibility, the decision-time account, design declaration,
method comparison and economic outcomes. Discovery should display ambiguous
endpoints and measurements or interventions that could distinguish them.
Latent-state analysis should show filtering uncertainty, fit diagnostics
and conditional policy value, preserving its identification restrictions.

The proposed SDK and evidence export should expose the operation records.
The interface should explain unsupported requests in the user's terms:
no untreated branches remain after the last rollout date, or severe-state
assistance was never observed. Worker identifiers and dependency hashes
belong in expandable technical evidence.

A bank deployment should reuse identity, catalog, warehouse and storage
services. Workers execute near authorized data; the control service holds
study metadata and permitted references. Regional boundaries cover derived
scores and fitted models as well as rows. R, Java and Python adapters each
require environment qualification and applicable license review.

Release acceptance reconstructs an exported study in a clean qualified
environment; detects late-data, support and sign errors; resumes interrupted
jobs without duplicate contributions; and reconciles financial totals
independently. Review-effort savings must then be measured on equivalent tasks.

\part{Evidence, investment and delivery}
\chapter{Evidence and validation}\label{ch:validation}
\section{Different claims require different evidence}
This proposal includes executable illustrations, preserved sources and
checked specifications. They demonstrate selected identities and interfaces,
not estimator superiority, production scale, institutional acceptance or
a profitable bank intervention.

Qualification distinguishes mathematical identity, implementation
correctness, simulation operating characteristics, real-data design
validity, prospective economic value and operational acceptance.
Success at one level informs the next; it does not replace it.

\begin{longtable}{P{34mm}P{60mm}P{48mm}}
\toprule Construction & Observed evidence & Limit\\\midrule
\endhead
Static adjustment & 40,000 independent synthetic units; AIPW $8.472$ versus truth $8$ and raw difference $55.586$ & One seed; no coverage study\\
Fixed policies & Paired AIPW, cost $7$, three prespecified comparisons with multiplicity adjustment & Fixed toy policies, not learned bank policies\\
Two-decision methods & All 16 cells enumerated; g-formula/IPW means $26.5,35.5$ & Exact population calculation\\
Staggered adoption & TWFE $1.2$; treated-cell ATT and clean DiD $3.5$ & Deterministic illustration\\
POMDP filter & Forward likelihood equals path enumeration: $0.0234985$ & Supplied two-state model\\
Transition fit & 700 randomized trajectories, horizon 14; maximum error $0.050$ & Known emissions/initial law; no coverage\\
Nonidentification & Equal observed laws, ATE $0$ versus $0.36$ & Constructive counterexample\\
Contracts & Four valid examples, ten invalid mutations, cross-artifact hashes & Selected semantics only\\
Economics & Recomputed NPVs and delay sensitivity & Planning assumptions\\
\bottomrule
\caption{Executed evidence and scope. Exact outputs and environment are saved.}
\label{tab:executed}
\end{longtable}

Examples use seed 20261007 and preserve numerical-library versions.
Figures derive from saved outputs. Three transition-optimizer starts give
closely agreeing objectives; this is a numerical diagnostic, not causal
identification. Contract checks exercise consequential boundary failures.

\section{A benchmark program that can disconfirm the product}
\subsection{Additional executed studies}
The larger synthetic studies supply fitted estimation, financial accounting
and adverse comparisons beyond the exact introductory examples. They retain
seed 20261007 with documented study-specific offsets. Truth is available
to the evaluator; it is not supplied as an unobserved input to the fitted
banking estimator.

\begin{longtable}{P{34mm}P{60mm}P{48mm}}
\toprule Construction & Observed evidence & Limit\\\midrule\endhead
Deposit rollout & 96 branches, 30 months, 499 branch bootstrap resamples;
balance gain but negative funding contribution & One synthetic rollout;
no validation of bank trend assumptions\\
Sequential collections & 3,000 training and 8,000 evaluation trajectories;
160 repeated datasets; paired policy contrasts & Known logging, complete
independent histories; finite-sample coverage failures retained\\
Controlled latent states & Unknown initial law, transitions and emissions;
hidden logging comparison; 79 bootstraps and 30 repetitions & Model
restrictions stipulated; capped fits and substantial value error\\
Credit supply & 600 firms, 24 banks, 1,800 relationships;
pooled and borrower-effect comparison & Synthetic exogenous bank shock;
one clustered interval\\
Common shocks & 1,200 panels; row, entity and crossed covariance
comparison & Gaussian additive score law; 36 date clusters\\
Temporal reconstruction & Late revisions, moved events, voids and source
completeness; independently stated account & Narrow additive reference,
not enterprise ingestion\\
Score reduction & Independent batch/partition agreement and fresh-process
measurements through ten million rows & Precomputed scores only; no
estimator training or storage workload\\
Reuse economics & Equal-output costs, threshold, adoption ramp and
unfavorable scenarios & Assumed costs and effort savings\\
\bottomrule
\caption{Additional executed evidence developed in the substantive
revision. Chapters containing each study specify its generating law and
what the result can support.}
\label{tab:revision-evidence}
\end{longtable}

The controlled-HMM verification enumerates all sixteen latent paths of
short observed sequences, with and without hidden-state action selection.
It agrees with the scaled forward/backward calculation for likelihood and
initial, transition, emission and action expected counts to below
$4\times10^{-15}$. This is a check on the implementation's probability
factorization, not a test that a bank satisfies it.

A separate three-decision collections enumeration sums 375 supported
state/action paths. The expanded DR score and a backward recursive
calculation agree within $10^{-12}$ on those paths. Expected scores equal
the exact policy value when the action law or action value is correct;
the chosen both-wrong case has bias about USD\,1.16. The current-ratio-only
calculation has expected value about negative USD\,22.91 against true
negative USD\,13.66. This isolates a target mismatch independently of
the larger study's Monte Carlo noise.

The latent-model follow-up also extends the main capped optimization on
the same data and starting seeds. It preserves the sign reversal from
omitting the action mechanism. The thirty-repetition comparison remains
a finite-budget experiment, with caps explicitly reported. An optimizer
stopping test is not a global-optimum certificate.

\subsection{What still requires bank qualification}
Vary outcome noise, treatment imbalance, support, hidden confounding,
serial/cross-sectional dependence, informative missingness, regimes,
interference and horizon. Generate truth separately from observed data;
prevent estimators and tuners from accessing oracle states or counterfactuals.
Freeze evaluation scenarios before comparison.

For effects, compare regression, IPW, g-computation, orthogonal/targeted
methods and justified neural challengers. Report bias, RMSE, coverage,
interval width, unsupported cases, abstention and subgroup behavior.
A method that reports a confident point effect in an unidentified
scenario should fail the scientific test even if optimization succeeds.

Panel cases include heterogeneous effects, staggered/reversible treatment,
anticipation, trend violations, few clusters and spillovers. POMDP cases
vary emission overlap, state dimension, private logging information,
support, horizon, misspecification and state aliasing. Assess predictive
fit, parameter recovery up to labels, value error and conservative decisions.

Discovery benchmarks require latent causes, contemporaneous relations,
nonstationarity, missingness and aggregation. Score adjacency and endpoints
against the correct equivalence-class target. Structural Hamming distance
alone can reward unjustified orientation or penalize justified uncertainty.
Include downstream identification and intervention-prediction tests.

An initial 500-replication coverage study has Monte Carlo standard error
about $\sqrt{0.95(0.05)/500}=0.00975$ at true 95\% coverage.
A 92--98\% observed coverage band is a proposed smoke criterion, not a
guarantee. Report Monte Carlo intervals, diagnose systematic failures,
and increase replication when decision precision warrants it.

\section{Reference agreement and adverse cases}
Adapters should reproduce exact small cases and independent reference
implementations within stated tolerances. Exact finite calculations permit
tight tolerances; stochastic optimization requires a different comparison.

Adverse tests include backdated features, duplicate entities, policy
changes, all-treated strata, missing propensities, outcome leakage,
receipt substituted for assignment, lost graph marks, sign reversal and
mixed currency. Test interruption/retry equivalence and full artifact
reconstruction. Return typed limitations rather than invent missing
information.

\section{Real-bank validation and experiments}
Begin with an existing randomized study and a rollout already understood
by bank economists. Reconstruct the original clocks, population and
analysis; compare results and review effort. Then conduct a prospective
study with prespecified financial/customer outcomes and institutionally
appropriate interventions.

For a balanced two-arm experiment with total independent sample $n$,
common standard deviation $\sigma$ and effect $\Delta$, rough planning gives
\[
 n\approx\frac{4\sigma^2(z_{1-\alpha/2}+z_{1-\beta})^2}{\Delta^2}.
\]
At $\sigma=100$, $\Delta=5$, $\alpha=0.05$ and 80\% power this is
about 12,550 units before attrition, clustering or multiple outcomes.
The equal-cluster design-effect approximation is $1+(m-1)\rho$
for mean cluster size $m$ and intracluster correlation $\rho$;
unequal clusters require more specific calculation.
These are illustrations, not bank sample-size commitments.

Compare matched workflow tasks on time to valid cohort, successful replay,
material findings, cost and decision relevance. A proposed target is 30\%
lower median analyst/reviewer effort without more unresolved material
findings. Establish baseline variability and task complexity first.
This target is a sponsor hypothesis to test.

\section{Value, uncertainty and customer outcomes}
Policy comparisons include contribution, loss, liquidity, capacity and
relevant customer outcomes. Requirements and thresholds belong to the
institution's decision process. Examine subgroup uncertainty, access,
support and error consequences. Predictive fairness metrics do not
establish causal fairness or legal permissibility.

Separate sampling, identification, model and economic-input uncertainty.
A credible interval conditional on a structural prior is not an identified
set. A downstream optimizer must receive the distinction explicitly.
Use sensitivity ranges alongside statistical intervals rather than
compressing every uncertainty source into an unjustified single number.

\section{Release acceptance}
Acceptance requires working workflows, replayable studies, qualified
dependencies, bank-data authorization, model validation and operational
ownership. Independent model validation is a future activity, distinct
from this author's review and synthetic calculations.

Stop or narrow a module when variation is absent, support fails, intervals
are miscalibrated, a simpler reference performs better at lower cost, or
value cannot justify complexity. A \emph{not identified} result can usefully
motivate an experiment; it is not a failed software job.

\chapter{The economic case for a shared product}\label{ch:economics}
\section{Compare the same scientific work}
The credible alternative is a capable internal team using its warehouse,
catalog, versioned notebooks and established open-source packages. That
team can implement the methods surveyed here. The proposed product must
justify the extra expense of reusable semantics, interfaces and operations
through repeated use, lower review effort or additional decision value.
Its comparison cannot assume that the alternative omits temporal correctness,
credible controls or model validation.

Define a completed study as a decision with a frozen population and
information clock, a stated target, defensible identifying restrictions,
at least one suitable reference comparison, an economic account and a
reconstructible evidence package. Both alternatives must deliver the same
number and mix of such studies, with equivalent scientific review.
The platform automates recurring work; expert composition performs more
of it in study-specific code and procedures. Neither is entitled to count
unreviewed model fits as completed studies.

The cost account has initial common construction $K_j$, annual shared
operation $F_{jt}$ and incremental study cost $c_{jk}$ for alternative
$j$. Study cost includes analyst, engineering and review effort, allocated
compute and recurring data reconciliation. Shared cost includes dependency
qualification, orchestration, support, security and reusable validation.
An existing bank catalog is a common baseline input; a new connection
required by one alternative is an incremental cost of that alternative.

\section{The reuse threshold}
For $n_t$ equivalent completed studies in year $t$, present cost is
\begin{equation}
 C_j=K_j+\sum_{t=1}^{H}
 \frac{F_{jt}+\sum_{k=1}^{n_t}c_{jk}}{(1+r)^t}.
 \label{eq:equivalent-cost}
\end{equation}
Suppose for illustration that annual shared cost and per-study cost are
constant, with equal study volume $n$. Let
$A_H(r)=\sum_{t=1}^{H}(1+r)^{-t}$,
$\Delta K=K_P-K_E$, $\Delta F=F_P-F_E$, and
$\Delta c=c_E-c_P>0$, where $P$ denotes the product and $E$ expert
composition. Then the product has lower construction/operation cost when
\begin{equation}
 n>\frac{\Delta K/A_H(r)+\Delta F}{\Delta c}.
 \label{eq:reuse-threshold}
\end{equation}
The numerator is its annualized common-cost premium; the denominator is
the saving per completed equivalent study. If $\Delta c\leq0$ and
common-cost premiums are positive, reuse alone cannot justify the platform.
No amount of repeating an equally expensive workflow removes that premium.

Heterogeneous studies need the original sum in
Equation~\eqref{eq:equivalent-cost}. A new country's undocumented action
logs may require nearly as much work with either stack. A twentieth
deposit rollout over a qualified data account may reuse most of its
transformation and review evidence. Counting both as identical use cases
would inflate the expected saving.

\section{Illustrative costs, including losing scenarios}
Use a three-year horizon and 10\% discount rate. All amounts in this section
are illustrative USD millions. Expert composition requires initial common
work of 1.0, annual shared cost 0.6 and 0.15 per completed study.
The product requires 4.0 initially, 1.8 annually and 0.06 per study.
These are planning assumptions to be replaced by a bank estimate. They are
not observed salary rates, procurement prices or measured productivity.

The study-cost difference is 0.09 million. At the separate staffing
assumption of USD\,25,000 per person-month, that is 3.6 person-months
saved per completed study. That is a large claim: it must be demonstrated
on equivalent tasks, not inferred from the existence of a common API.
The product's annual shared cost includes support and dependency work;
study-specific engineering and review are included in its per-study cost.
The categories must not overlap.

\begin{table}[htbp]\centering\small
\begin{tabular}{@{}lrrr@{}}
\toprule Scenario & Expert cost & Platform cost & Platform saving\\\midrule
Low reuse & 5.476 & 9.670 & -4.194\\
Moderate reuse & 11.445 & 12.057 & -0.613\\
High reuse & 15.921 & 13.848 & 2.073\\
Adoption ramp & 8.402 & 10.840 & -2.438\\
\bottomrule\end{tabular}

\caption{Three-year present cost in illustrative USD millions. Constant
annual volumes are 8, 24 and 36 studies; the ramp is 4, 16 and 30.
Positive saving favors the shared product.}
\label{tab:reuse-economics}
\end{table}

The cost-only threshold is about 26.74 completed studies per year. At
eight annually the product costs about USD\,4.19 million more in present
value. At 24 it still costs about USD\,0.61 million more. At 36 it saves
about USD\,2.07 million. An adoption ramp of 4, 16 and 30 studies leaves a
USD\,2.44 million premium even though final-year throughput exceeds the
constant-volume threshold. Early costs occur before much reuse.

\begin{figure}[htbp]\centering
\includegraphics[width=.9\textwidth]{derivations/reuse_economics.pdf}
\caption{Equal-output cost comparison under declared assumptions. The
crossover is a hypothesis about repeatable work and costs, not a forecast
of bank adoption.}
\label{fig:reuse-economics}
\end{figure}

If the per-study saving is half as large, the threshold doubles to about
53.47 studies annually. If the annual shared-cost premium rises from
1.2 to 2.0 million, the original per-study saving needs about 35.63
studies annually. These sensitivities argue for measuring reusable effort
early and keeping common construction narrow until multiple teams use it.
A platform built for imagined reuse can be economically worse than a
well-maintained expert toolkit.

\section{Separate decision value from construction savings}
Let $D_t$ be additional realized business value attributable to the product
relative to the equivalent expert workflow. The incremental investment case is
\begin{equation}
 \Delta\mathrm{NPV}=C_E-C_P+
       \sum_{t=1}^{H}\frac{D_t}{(1+r)^t}.
 \label{eq:incremental-product-npv}
\end{equation}
The first term already counts workflow cost savings. Do not add the same
released analyst time again as a separate business benefit. If that time
supports additional studies, model their costs and incremental decisions
explicitly, with a capacity constraint.

At 24 studies per year the illustrative cost premium requires about
USD\,0.246 million of additional constant annual decision benefit.
For the adoption ramp it requires about USD\,0.980 million. Those may be
modest amounts for a large banking group, but group size is not evidence
that the product creates them. The expert comparator may already find the
same beneficial intervention. An analyst's higher point estimate is not
realized incremental value.

The deposit study gives a useful distinction. Discovering that a promotion
loses contribution can avoid an expansion loss, but the prevented loss is
attributable only if the comparator would otherwise expand it and the
finding changes the decision. The collections simulation's positive policy
gain is a property of its generating model; it supplies no bank revenue
for Equation~\eqref{eq:incremental-product-npv}. A hidden-state module may
have positive research value while its operational policy benefit remains
unresolved.

Risk reduction is also a counterfactual claim. If the platform catches a
late-data feature, record its frequency, review disposition and consequence
under the baseline process. Do not multiply an arbitrary low-probability
catastrophe by a large bank balance sheet and call the result a measured
benefit. Early validation can instead report concrete errors detected and
the cost of resolving them.

\section{Adoption, maturity and cash timing}
Not every initiated study reaches a usable conclusion. A simple planning
decomposition is
\[
 n_t^{\rm usable}=n_t^{\rm started}\,
 p_t^{\rm data}\,p_t^{\rm design}\,p_t^{\rm completed}
 \,p_t^{\rm review},
\]
where the probabilities are conditional progression rates, not assumed
independent events. The multiplicative form is the chain rule when each
rate conditions on prior progression. Costs incurred by failed studies
remain in the account. A well-designed product may appropriately terminate
more unidentified studies early; usable throughput should then be assessed
alongside avoided futile effort.

Decision outcomes mature later still. A six-month deposit outcome can
become observable within a pilot year; losses from a credit-limit change
may take longer. A continuing value model introduces assumptions rather
than making those outcomes observed. Discounting should follow the actual
date of expenditure and realized benefit, with distinct data-access,
development, study and outcome lags.

The earlier aggregate illustration can be retained as a stress calculation:
a USD\,4 million outlay and USD\,1.8 million annual operation cost need
constant annual incremental benefit of approximately
USD\,\NumBreakEven{} million over three years at 10\%.
At USD\,5 million annual benefit its NPV is USD\,\NumBaseNPV{} million;
losing the first year's benefit makes it about negative USD\,0.588 million.
Those values use a coarser account than the equal-output comparison and
must not be added to it. The explicit comparator model is the preferred
investment account.

\section{Measuring the product hypothesis in the first tranche}
Choose a small set of matched tasks across the three workspaces: reconstruct
a rollout, evaluate a repeated-decision policy under a defensible design,
and organize a temporal discovery dataset with known semantic ambiguities.
Agree their required outputs and reviewer rubric before timing them.
Give the expert-composition arm the same source access, compute budget and
scientific expertise. It may reuse its own sensible templates.

Measure active analyst and reviewer hours separately from calendar delay,
data-owner response time and machine time. Record task complexity,
previous familiarity, failed starts, rework and unresolved material findings.
A counterbalanced assignment of comparable tasks can reduce learning-order
effects; repeated use by several teams is necessary to estimate reuse
rather than first-author familiarity. A proposed target of 30\% lower
median analyst/reviewer effort is an investment hypothesis, not a result
already established by this proposal.

The platform should also export a study to a second analyst who did not
author it. Measure time to reproduce the target, source account and
financial conclusion, and whether the analyst notices deliberately
introduced temporal, support or sign errors. Equivalent review quality is
part of the output definition. Faster execution with more unrecognized
scientific defects is not a saving.

The first tranche is six people for six months: one product/economics lead,
two causal/econometric scientists, two data/platform engineers and one
application engineer. At the illustrative loaded rate, staffing is
USD\,900,000; compute/tooling USD\,50,000 and 20\% contingency bring the
total to USD\,1.14 million. That tranche is included within the illustrative
USD\,4 million platform construction, not added to it. Bank-owner,
independent-validation and security participation remain unpriced and
must be estimated consistently for both alternatives before funding.

Six people cannot production-qualify every method in the census in six
months. They can construct the common account and evidence layer, qualify
a narrow panel/longitudinal path, and deliver finite-state estimation and
temporal discovery as explicit research capabilities. The next investment
depends on observed scientific usefulness and reusable effort rather than
on completing a broad feature checklist.

\section{Independent product economics and procurement}
An independent offering adds customer-specific deployment, support,
upgrades, contracts and onboarding costs. Code reuse does not eliminate
each institution's semantic mapping and validation. A plausible initial
commercial form is a design-partner deployment with a paid construction
scope and recurring support, preserving SDK and evidence export.
Willingness to pay, sales cycle, support load and customer concentration
have not been measured here.

Compare a commercial vendor on the same reconstructed studies, permission
boundaries, supported identification routes, exported uncertainty and
dependency responsibilities. Vendor demos and list prices are insufficient;
no current private enterprise capability or price is asserted in this survey.
For internal adoption, chargeback should distinguish common capability
from a business unit's bespoke model so that demand is not inflated by
unpriced customization.

The sponsor's decision is therefore conditional but concrete: fund a narrow
shared construction if named teams have repeatable decisions, reconstructible
data and a credible path to reusable work. Continue expert composition where
volume is low or assumptions differ too much. Expand the common product
when measured study cost, quality and adoption support the incremental
account. The financial case becomes stronger through evidence of actual
reuse, not through assigning all of a bank's causal-analysis benefit to a
new software layer.

\chapter{Delivery, ownership and stopping decisions}\label{ch:delivery}
\section{Develop the central capabilities together}
The initial program develops three workspaces: policy and transmission,
sequential decisions with controlled latent-state estimation, and temporal
discovery with automatic data organization. They share a temporal account,
study semantics and evidence system. Their scientific claims mature at
different rates. A useful first POMDP module can estimate a declared model
and display conditional policy scenarios; a discovery module can return
ambiguous endpoints and candidate experiments. Neither needs to claim that
observational data uniquely identify every intervention.

The first six-month tranche funds a small common implementation and
bounded application paths. It should qualify a narrow panel/longitudinal
analysis, an inspectable finite-state estimation path and a temporal
discovery adapter. It should not attempt production support for all
\SoftwareCount{} surveyed projects. The illustrative USD\,1.14 million
team and the equal-output economic comparator are developed in
Chapter~\ref{ch:economics}; future funding depends on observed reuse and
scientific usefulness.

\section{A dependency-led construction program}
\begin{longtable}{P{24mm}P{57mm}P{61mm}}
\toprule Stage & Construction & Evidence to continue\\\midrule\endhead
Weeks 1--4 & Select a funding rollout, repeated-decision history and temporal
mechanism question; reconstruct their information and financial accounts &
Named owners, usable variation, action-log provenance and equivalent
expert-workflow baseline\\
Months 2--3 & Common snapshot/compiler/evidence path; cohort and longitudinal
adapters; controlled finite-state likelihood; temporal graph adapter &
Exact/reference agreement, valid endpoint serialization, leakage/support
refusals and explicit identification states\\
Months 4--6 & Complete limited workspaces; bank retrospective studies;
reusable SDK exports; independent user reconstruction & Qualified methods
for stated tasks, measured workflow effort, reconciled outcomes and
identified research limitations\\
Months 7--12 & Broaden supported histories, latent measurements and
discovery designs; production deployment qualification; workload profiling &
Independent model review, calibrated intervals, actual bank-scale costs,
permission/retry tests and a defensible prospective pilot\\
Months 12--18 & Cross-entity reuse and additional policy applications;
qualified distributed execution; design-partner productization &
Repeatable adoption, outcome maturity, support economics and measured
incremental value over expert composition\\
\bottomrule
\caption{Proposed construction after this research commission. Dates are
planning ranges; readiness and evidence determine progression.}
\label{tab:roadmap}
\end{longtable}

The first month is substantive design work. For deposits, identify which
branches changed policy, whether customers anticipated it, whether external
funds can be distinguished from internal transfers and whether untreated
comparisons remain. For collections, inspect how staff discretion and
unrecorded information entered actions. For discovery, define the frequency,
measurement coverage and admissible contemporaneous relationships before
choosing a graph algorithm. A warehouse extract without these answers does
not make a model ready.

The next stage implements one end-to-end path before generalizing. A
cohort study is reconstructed from original events, fitted with a reference
comparison and exported with target weights and contribution uncertainty.
The latent-state path starts with the exact count/likelihood checks in
Chapter~\ref{ch:latentestimation}; the discovery path preserves equivalence
classes and assumptions. These paths exercise common components under
different scientific demands and reveal which abstractions are truly reusable.

By month six the platform should support another analyst repeating a
bounded study without help from its author. That is a stronger acceptance
condition than a polished first-author demo. A research-only state remains
appropriate for a module without a bank identifying design. The bank may
continue useful data and scenario work while funding a measurement or
randomization plan to answer the unresolved causal question.

\section{Who owns each decision}
The business owner specifies the action, feasible comparator, economic
account and decision threshold. The data owner attests to historical
coverage, identity and timestamp semantics. The science owner signs the
identification argument, estimator and inference design. The engineering
owner signs computational reproducibility, access and operation. Independent
model validation challenges the complete analysis under the institution's
applicable framework. None of these roles is replaced by a successful
schema check.

Changes have different approval and requalification consequences within
the future bank process. A documentation correction may leave a result
unchanged. A changed funds-transfer price changes the economic scenario.
A new treatment definition or population changes the estimand. A package
upgrade can alter covariance defaults or missing-value handling. The
study record should show which objects changed and which previous evidence
still applies, rather than carrying forward one undifferentiated approval.

The shared product owner maintains supported combinations and deprecation
policy. An adapter version can remain available for historical replay while
being barred from new decision use after a material defect. Historical
reports should identify the affected dependency and receive an attached
notice; they should not be silently rewritten as if the corrected result
had always been known.

\section{Qualification gates with concrete failure outcomes}
A data gate fails if the requested decision-time account cannot be
reconstructed, material transfers do not reconcile, or missing coverage is
mistaken for zero. The response is a corrected source mapping, a narrower
population or a data-collection plan. It is not permission to use a final
warehouse view as if it were historically available.

An identification gate fails if no adequate treatment variation or proxy,
instrument, randomization or structural route is available. The response
can be an identified set, a conditional scenario or a proposed experiment.
The product can still report useful information, but it cannot label the
result an identified point effect. A discovery graph or predictive latent
fit alone cannot waive this gate.

An estimation gate fails if the declared method does not fit the data
shape, overlap is insufficient for the intended intervention, numerical
fits are unstable or relevant simulated intervals are badly miscalibrated.
The response may be a shorter policy horizon, a supported intervention,
a simpler estimator or an explicit unresolved result. Weight clipping,
state aggregation and policy restriction change the analysis and must be
shown as such.

An economic gate fails if contribution definitions are incompatible,
uncertainty crosses the decision threshold, or the product's added cost is
not justified against expert composition. A scientifically valid null
finding can still save the bank from a poor action. It does not establish
that an expensive shared platform is the cheapest way to obtain that
finding. Module and platform decisions therefore have separate accounts.

An operational gate fails if another authorized user cannot reconstruct
the evidence, if retries duplicate contributions, if regional access is
violated or if dependency ownership is absent. The resolution is an
engineering or operating change. Such defects should not be hidden inside
a scientific sensitivity interval.

\section{Prospective learning and model maintenance}
A prospective pilot should compare a small prespecified set of feasible
strategies, with economic and customer outcomes, a horizon and stopping
rules agreed in advance. Policy learning must respect a held-out evaluation
or a suitable selection-adjusted design. An attractive fitted value from
the same data used to optimize the policy is insufficient evidence of
improvement.

Monitoring separates population change, assignment change, observation
change and response change. A rise in missed measurements can change a
latent-state likelihood; a new staff override process can change action
selection; an interest-rate regime can change deposit response. A common
alert on prediction error does not diagnose all three. The model record
should link each alert to the affected assumption and evidence needed for
continued use.

Outcome maturity affects monitoring as well as economics. Early repayment
can be observed before ultimate losses; closing an account can end product
balance records without ending its group-level economic outcome. A status
page should distinguish complete computation from mature outcomes and
future-value assumptions. This makes a study useful before lifetime results
arrive without misrepresenting what has been learned.

\section{Research priorities that change delivery choices}
The highest-value research reduces a specific uncertainty in a banking
application: whether private logging can be measured, whether weak proxies
support stable policy value, whether common-shock inference is calibrated,
or whether a proposed discovery restriction survives aggregation. Each
research ticket should state the competing constructions, an observable
failure criterion and the product decision it changes.

Chapter~\ref{ch:frontier} discusses current research through the survey
cutoff. Foundation models and automated semantic systems are candidates
for reducing fitting or organization effort, not assumed replacements
for identification. Hidden-state, irregular-time and nonstationary
methods remain central research from the beginning, with release claims
limited to the evidence actually obtained.

The practical stopping rule is modular. Keep a useful temporal account,
panel workflow or conditional latent-model workbench when another module
fails its causal or economic qualification. Stop expanding a shared layer
when each additional application requires bespoke work that erases the
assumed reuse saving. Expand when multiple teams can reconstruct useful
studies at acceptable cost and quality with a clear owner for every
scientific and operational claim.

\chapter{Current research and unresolved product choices}\label{ch:frontier}
\section{A new method must change a decision}
Recent research matters when it changes a feasible estimand, required
measurement, uncertainty calculation or computational cost. Publication
recency alone is not a reason to replace a qualified estimator. The survey
cutoff is 7 October 2026; several cited 2026 works are recent preprints,
including one posted the preceding day. Their archived versions and
provisional status are retained in the source register. None has been
runtime-qualified on bank records in this commission.

The central unresolved questions cut across method families. Can the bank
identify the information used by its historical decision maker? Can it
distinguish a customer response from a common economic shock? Does a
measurement arrive because risk increased? Does a fitted latent model
preserve the response when policy changes? Does a faster estimator merely
move untested assumptions into its training prior? These questions determine
the useful role of current research.

\section{Amortized inference moves work into a prior}
Prior-data fitted networks train on many synthetic datasets and learn a
mapping from a new dataset to an approximate posterior or estimator.
If successful for a supported task, this can reduce repeated tuning and
fitting cost. Let $\Pi$ be a prior over observational/interventional
model pairs and $Q$ a causal query. The learned approximation targets a
prior-conditional object such as
\[
 \Pi(Q\mid D)\ \propto\ \int
       \ind\{Q(M)\in dQ\}\,p(D\mid M)\,\Pi(dM).
\]
The model family and its identifying restrictions enter before observing
the bank dataset. Approximation to this posterior is a computational
achievement; it does not verify that the prior's exclusion or invariance
restrictions hold for that bank.

The inspected ICLR 2026 CausalFM paper makes identification an explicit
input, including declared backdoor, frontdoor and instrumental-variable
settings \citep{causalpfn2026}. Its Section 3 defines an identified query
as a functional of the observational law on the maintained model class.
This distinction is valuable for the product: an amortized estimator can
sit behind the same design contract as a conventional estimator.
It does not infer a valid instrument from a column's predictive usefulness.
The archived work is CausalFM, distinct from the concurrent CausalPFN
work it discusses.

Consider two models with the same observed distribution but different
intervention effects, as constructed in Chapter~\ref{ch:pomdp}.
Their likelihood ratio remains one for every observational sample.
Posterior odds between them therefore remain prior odds. A network cannot
learn a uniformly correct effect from that same observed law without
additional restrictions or information. If its prior excludes one model,
a narrow posterior is conditional on that exclusion. A reliability head
trained to recognize simulated regimes inherits the same boundary.

The temporal PFN preprint describes temporal tokens, a causal-judgment
head, context selection and posterior correction \citep{panelpfn2026}.
These are interesting candidates for reducing repeated computation and
prioritizing uncertain analyses. The stated source does not establish that
an inferred identifiability score is a mathematical certificate for an
arbitrary banking panel. Its industrial examples also have different
measurement and action processes from endogenous bank pricing or
relationship-manager discretion.

Qualification should hold out entire generating mechanisms, not merely
new draws from the training generator. Include weak instruments, proxy
exclusion violations, hidden policy changes, long shared shocks and
delayed revisions absent from the synthetic prior. Compare effects and
interval calibration with conventional estimators on the same target;
measure context selection and preprocessing as well as inference time.
For an amortized product, pretraining and upgrade cost belongs in the
economic account even when per-study fitting is fast.

\section{Proximal learning: the inverse problem survives automation}
The ProximalFM preprint, posted on 6 October 2026, explicitly adopts
latent exchangeability, proxy exclusions, latent positivity and
completeness; it additionally assumes conditions for an outcome bridge
to exist \citep[Sections 2.1--2.2]{proximalfm2026}. It predicts a CATE
posterior from synthetic proximal models instead of solving a new bridge
optimization for each dataset. The target averages over proxies
conditional on the designated covariates $X$.

Those roles have concrete banking meaning. A candidate treatment-inducing
proxy must carry information about hidden risk without violating the
outcome exclusion once action and the latent causes are held fixed.
An outcome-inducing proxy must satisfy its own conditional independence
from action and the first proxy. A staff assessment that directly changes
both assistance and later follow-up may violate these restrictions.
More predictive proxy variables are not automatically better identifying
measurements.

In the finite bridge calculation, small singular values of the relevant
conditional-probability matrix amplify estimation noise. In a functional
bridge, the corresponding operator can be ill posed. Regularization trades
bias and instability under a chosen function class
\citep{ghassami2022,proxreg2024}. Amortization may encode a useful
regularizer through its prior, but does not remove the need to examine
weak-proxy and prior-sensitivity cases. The bank should receive value
sensitivity and unsupported conditions alongside any fast point estimate.

The product research question is therefore whether a specified proxy
measurement design supports useful policy precision. A promising next
experiment can collect an additional pre-action measurement under a
known process, compare candidate bridges and test value sensitivity.
Calling every transaction-derived embedding a proxy would evade the
scientific problem rather than solve it.

\section{Coverage at one date can hide long-horizon difficulty}
Positive action probabilities at each date are necessary for many
off-policy estimators but can leave target trajectories extremely rare.
The collections study demonstrates weight concentration even with
conditional randomization. Markov occupancy-ratio estimators can improve
efficiency when the state representation and sampling restrictions support
their factorization \citep{kallus2020}. A compressed learned state must
retain the relevant information before that argument can be used.

A recent hardness preprint constructs exponentially difficult evaluation
with an unconfounded history-dependent logger, a deterministic memoryless
target and favorable behavior-marginal coverage quantities
\citep[Definition 3.1 and Theorem 4.1]{opehard2026}. Its result concerns
a specific model class and definitions of revealing/coverage. It does
not imply that every bank POMDP is exponentially difficult. It does
refute an overly broad product claim that good one-step coverage metrics
alone establish easy long-horizon evaluation.

This suggests a hierarchy of proposed outputs: short prespecified policy
contrasts where support is adequate, model-conditional longer scenarios,
and learning designs that deliberately improve information about the
target policy. Report where the policy's history distribution departs
from observed experience. A large total transaction count can coexist
with little evidence on the sequence needed for a decision.

\section{Partial identification needs a temporal confounding model}
When point identification fails, bounds can still help choose among
actions. But a dynamic bound cannot be assembled from arbitrary
independently chosen local worst cases if those choices cannot coexist
within one underlying process. The shape of the uncertainty set matters.

The inspected sharp-dynamic-bounds preprint restricts unmeasured
confounding to act within a stage and not persist as a hidden state into
later stages \citep{sharpdynamic2026}. Its local feasible transition sets
then combine under the stated temporal restrictions, supporting a
backward robust evaluation. Persistent wealth, financial fragility or
relationship quality in banking can violate precisely that restriction.
The two-state process in Chapter~\ref{ch:latentestimation} is an example
where persistence is central rather than an incidental nuisance.

The product should therefore distinguish bounds from stage-local
confounding, sensitivity sets for a persistent latent model and
history-based longitudinal restrictions. Each has its own target and
attainability argument. A visually wider interval does not automatically
contain all plausible bank mechanisms. An optimization routine finding
extreme values within a declared set establishes only that conditional
calculation.

\section{Irregular observations are often informative}
Bank records have irregular calls, document requests and account activity.
Regularizing them to a monthly grid changes the measurement process.
The ITSY preprint couples completion and temporal graph estimation,
with its displayed setup assuming causal sufficiency, stationary
mechanisms and MCAR missingness; its implemented structural model uses
first-order lags \citep[Section 3]{itsy2026}. Those restrictions give a
clear comparator domain. They also exclude many risk-triggered measurement
processes that a banking product must represent explicitly.

For example, a customer contacted after deteriorating arrears is more
likely to provide an updated income measure. The absence of that measure
is then informative about the underlying process or policy. A completion
algorithm that treats the gaps as independent can alter both contemporaneous
and lagged associations. Discovery must model or restrict observation
selection, or return a limitation, before interpreting the completed graph.

Continuous-time causal SDE work provides another language for interventions,
with existence, solvability and structural requirements
\citep{sde2026}. A mathematical intervention construction does not by
itself supply an estimator for irregular bank logs. A bank pilot must
specify the sampling process, what action operates between observations
and which parameters are recoverable at that frequency.

Other recent proposals use spectral signatures to route corrections for
different confounding patterns, accelerate temporal conditional-independence
tests, or constrain discovery through group-level background knowledge
\citep{lucid2026,causalts2026,vargas2026}. They motivate explicit
challengers. Their reported speed or graph improvements were not reproduced
here, and background constraints remain substantive assumptions.

\section{The research portfolio has measurable exits}
Each research line should end in a product decision. An amortized estimator
earns an adapter when it improves equivalent-task cost with acceptable
calibration under relevant mechanism shifts. A proxy method earns a
policy-value claim when a defensible measurement design and stability
analysis support it. A discovery extension earns use when its graph
semantics and observation assumptions match the bank's records. A latent
model earns a planning role when its causal route, numerical fit and
decision uncertainty are adequate for the proposed action.

The negative outcomes are useful too: retain a conventional estimator,
collect better measurements, shorten a policy horizon, report a set or
use discovery to plan an intervention. This portfolio treats new methods
as competing solutions to bank questions. It avoids committing the
independent product to a promise that the next model architecture will
eliminate identification, dependence or data-semantic work.

\appendix
\chapter{Research protocol and source concordance}\label{app:evidence}
\section{Coverage and provenance}
The local inventory contains 7,891 PDFs across the user's library and the
FinMarketing paper folder. Discovery used Crossref, OpenAlex, arXiv,
proceedings, author copies and official repositories. It retained 64
public search requests: 59 saved and five failed. These yielded \CandidateCount{} unique
normalized-title candidates. The final census includes \PublicRequestCount{} total registered
public requests. Counts describe collected records, not all literature.

This is a purposive review. Sources were selected when they changed an
assumption, estimator, failure mode, banking design, software choice or
requirement. Recent searches covered 2024 through 7 October 2026;
foundational work had no date restriction. Preprints retain explicit
version labels. The banking-specific survey \citep{banksurvey2023}
supplies discovery leads; primary papers support substantive claims.

Every cited scholarly work has a real local PDF named by title,
first-author surname and citation year. The register records the actual
archived version, which may be a working manuscript corresponding to a
later publication. Original titles remain in metadata when prohibited
filename punctuation is removed. Three official guidance PDFs are also
preserved. Software citations refer to manifests whose component
documents have content hashes.

Records include sources, reading notes, claims, search/request ledgers,
candidates, metadata audit, software assessment, local-repository
snapshots and omissions. The concordance maps citation keys to identities;
an HTML index opens the saved PDFs. Hashes detect file changes but do
not establish scientific correctness.

\section{Integrity repairs and access limits}
Local filenames were treated as hints. A file labeled as a classic
dynamic-panel-bias paper contained a different work and was excluded.
The liquidity-shock working paper is from 2006 and lists Mian first;
the later publication's reversed author order and year were not applied
to that copy. A mortgage-refinancing filename dated 2023 contained an
April 2020 revision. Guessed arXiv identifiers that resolved to unrelated
papers were retained as raw excluded downloads and replaced only after
identity verification.

The expanded 2025 ACM banking survey was located in metadata but its PDF
was inaccessible. The 2023 precursor is labeled accordingly. A 2023
Nature Reviews time-series review was not acquired and is not cited.
An unavailable recurrent marginal structural-network candidate is recorded
as an omission; accessible primary longitudinal sources support the chapter.

Exact normalized-title metadata matches were available for 27 selected
papers in the saved corpus. Available update/retraction fields and citation
counts are preserved. This is not complete publisher, Crossmark or
retraction-database clearance. Citation counts and venues were not used
as evidence of correctness.

Coverage is strongest for effects, panels, longitudinal methods, partial
observation, temporal discovery and infrastructure. It is less comprehensive
for continuous-time econometric identification, network interference,
transportability, online experiments and all jurisdictions' requirements.
Those limits constrain the claims.

\section{Reconstruction and evidence classes}
The source catalog rebuilds bibliography and records from copies and notes.
The example runner regenerates numerical outputs and figures.
Contract examples contain hashes of the frozen numerical artifact;
regeneration requires deliberate refresh of those links.
Contract checks, document compilation and source/consistency audits complete
the reconstruction.

Document builds do not download sources. Public acquisition uses the
separately approved bounded collector. Saved receipts include failures and
dates. Original PDFs and metadata remain available offline.

The author review examined mathematics, evidence boundaries, interfaces,
citations and rendered layout. It is not independent validation.
Bank validation, third-party package runtime qualification and full platform
implementation remain future obligations. The two narrow research components
and executed synthetic studies are described separately in the main text.

\section{Notation and units}
\begin{longtable}{P{28mm}P{112mm}}
\toprule Symbol & Meaning\\\midrule\endhead
$i,t$ & Entity/trajectory and decision/calendar time\\
$A_t,L_t,H_t$ & Action, pre-action covariates and recorded history\\
$Y(a),Y^g$ & Outcome under fixed action or policy\\
$e_t,g_t,\pi_t$ & Behavior assignment, intervention law and target policy\\
$Q_t,V_t,W_t$ & Action value, policy-averaged value and cumulative importance ratio\\
$S_t,O_t,b_t$ & Latent state, available observation and filtered belief\\
$P,B,F,\mu$ & Transition, emission, reward and initial-state distributions\\
$G_i,\ATT(g,t)$ & Adoption date and group-time effect on treated\\
$\theta,\tau(x)$ & Population and conditional average effects\\
$C_0,O_t,B_t$ & Build cost, operating cost and benefit in the economics chapter; $B_t$ there is not an emission\\
\bottomrule
\end{longtable}
The introductory examples use synthetic currency units. The larger banking
studies and investment accounts explicitly use illustrative USD, with balance
tables labeled in millions where applicable. Real studies require currency, valuation
date, horizon and sign.

\section{Cited source concordance}
\small
\begin{longtable}{P{32mm}P{112mm}}
\toprule Audit key & Source identity and archived version\\\midrule\endhead
{\scriptsize\texttt{spectral2016}} & Reinforcement Learning of POMDP's using Spectral Methods \citep{spectral2016}.\par arXiv v1, 25 February 2016; 62-page author preprint, title retains its displayed possessive spelling\\
{\scriptsize\texttt{newey1986}} & A Simple, Positive Semi-Definite, Heteroskedasticity and Autocorrelation Consistent Covariance Matrix \citep{newey1986}.\par NBER Technical Working Paper 55, April 1986; scanned original visually inspected; later 1987 journal identity located but these are 1986 bytes\\
{\scriptsize\texttt{lpaug2020}} & Local Projection Inference is Simpler and More Robust Than You Think \citep{lpaug2020}.\par arXiv v1, 27 July 2020; internal version 29 July 2020; later 2021 Econometrica publication DOI 10.3982/ecta18756 identified in metadata, not substituted\\
{\scriptsize\texttt{inoue2023}} & Significance Bands for Local Projections \citep{inoue2023}.\par arXiv v1, 5 June 2023; cover 6 June 2023; later metadata title Inference for Local Projections is not the inspected version\\
{\scriptsize\texttt{ciganovic2026}} & Double Machine Learning for Time Series \citep{ciganovic2026}.\par arXiv v1, 11 March 2026; subsequent Econometrics Journal identity 10.1093/ectj/utag019 located in metadata; technical assessment uses original v1\\
{\scriptsize\texttt{babii2020}} & Machine Learning Time Series Regressions With an Application to Nowcasting \citep{babii2020}.\par User-preserved arXiv v4, 12 December 2020; internal date 15 December 2020; filename hint 2021 is not the version year\\
{\scriptsize\texttt{johansen2017}} & Cointegration between Trends and Their Estimators in State Space Models and Cointegrated Vector Autoregressive Models \citep{johansen2017}.\par Published Econometrics 5, 36, 22 August 2017; journal PDF retained from user collection\\
{\scriptsize\texttt{roysland2024}} & Graphical criteria for the identification of marginal causal effects in continuous-time survival and event-history analyses \citep{roysland2024}.\par Author preprint arXiv v2, 20 April 2024; journal online identity located by DOI, publisher bytes not substituted\\
{\scriptsize\texttt{rytgaard2021}} & Continuous-time targeted minimum loss-based estimation of intervention-specific mean outcomes \citep{rytgaard2021}.\par arXiv v1, 5 May 2021; cover 6 May 2021; later 2022 Annals article identified in metadata, not the inspected version\\
{\scriptsize\texttt{ryalen2026}} & On causal inference with marked point process data \citep{ryalen2026}.\par Preprint arXiv v1, 14 April 2026; theoretical comparator, not independently verified\\
{\scriptsize\texttt{imbens2002}} & Identification and Estimation of Triangular Simultaneous Equations Models Without Additivity \citep{imbens2002}.\par NBER Technical Working Paper 285, November 2002; later Econometrica 2009 version located in metadata but not substituted for these bytes\\
{\scriptsize\texttt{pearltransport2014}} & External Validity: From Do-Calculus to Transportability Across Populations \citep{pearltransport2014}.\par Statistical Science 29(4), 579--595 (2014); electronic journal reprint on arXiv v1, 5 March 2015, with reprint pagination\\
{\scriptsize\texttt{petersicp2016}} & Causal Inference by using Invariant Prediction: Identification and Confidence Intervals \citep{petersicp2016}.\par JRSS B (2016) identity verified by DOI; preserved author manuscript title omits 'by', archive header v3 dated 24 November 2015, internal cover dated 24 May 2024; technical claims use this exact archived copy\\
{\scriptsize\texttt{howard2021}} & Time-uniform, Nonparametric, Nonasymptotic Confidence Sequences \citep{howard2021}.\par Annals of Statistics (2021) identity verified by DOI; preserved later author revision arXiv v9 dated 6 August 2022, internal date 9 August 2022\\
{\scriptsize\texttt{sarkka2013}} & Bayesian Filtering and Smoothing \citep{sarkka2013}.\par Author-distributed 2013 book PDF from user collection; IMS Textbooks volume 3; personal-use notice preserved\\
{\scriptsize\texttt{kalouptsidi2021}} & Identification of Counterfactuals in Dynamic Discrete Choice Models \citep{kalouptsidi2021}.\par Quantitative Economics 12 (2021), 351--403; published copy in user collection\\
{\scriptsize\texttt{imai2009}} & Bayesian Estimation of Dynamic Discrete Choice Models \citep{imai2009}.\par Econometrica 77(6), November 2009, 1865--1899; published copy in user collection\\
{\scriptsize\texttt{zhang2008}} & Causal Reasoning with Ancestral Graphs \citep{zhang2008}.\par Journal of Machine Learning Research 9 (2008), 1437--1474; published July 2008, publisher PDF\\
{\scriptsize\texttt{perkovic2018}} & Complete Graphical Characterization and Construction of Adjustment Sets in Markov Equivalence Classes of Ancestral Graphs \citep{perkovic2018}.\par Journal of Machine Learning Research 18 (2018), 1--62; revised April and published May 2018, publisher PDF\\
{\scriptsize\texttt{targeted2018}} & Targeted Learning in Data Science: Causal Inference for Complex Longitudinal Studies \citep{targeted2018}.\par Springer 2018 electronic book from the user collection; Chapter 4 is by Sherri Rose and Mark J. van der Laan\\
{\scriptsize\texttt{staiger1994}} & Instrumental Variables Regression with Weak Instruments \citep{staiger1994}.\par NBER Technical Working Paper 151, January 1994; inspected scan, not the later journal version\\
{\scriptsize\texttt{hall2005}} & Nonparametric Methods for Inference in the Presence of Instrumental Variables \citep{hall2005}.\par Annals of Statistics 33(6), 2005, pp.2904--2929; electronic reprint arXiv math/0603130v1 uploaded 6 March 2006, with different pagination\\
{\scriptsize\texttt{hernan2020}} & Causal Inference: What If \citep{hernan2020}.\par Book; local draft dated 21 February 2020\\
{\scriptsize\texttt{peters2017}} & Elements of Causal Inference: Foundations and Learning Algorithms \citep{peters2017}.\par MIT Press book, 2017\\
{\scriptsize\texttt{dml2018}} & Double/Debiased Machine Learning for Treatment and Structural Parameters \citep{dml2018}.\par Econometrics Journal 2018; archived arXiv v7, 3 November 2024\\
{\scriptsize\texttt{kunzel2019}} & Meta-learners for Estimating Heterogeneous Treatment Effects using Machine Learning \citep{kunzel2019}.\par PNAS 2019; arXiv v6, April 2019\\
{\scriptsize\texttt{notears2018}} & DAGs with NO TEARS: Continuous Optimization for Structure Learning \citep{notears2018}.\par NeurIPS 2018 proceedings copy\\
{\scriptsize\texttt{callaway2021}} & Difference-in-Differences with Multiple Time Periods \citep{callaway2021}.\par Journal of Econometrics 2021; local manuscript dated 1 December 2020\\
{\scriptsize\texttt{roth2023}} & What’s Trending in Difference-in-Differences? A Synthesis of the Recent Econometrics Literature \citep{roth2023}.\par Journal of Econometrics 235, 2218–2244; published PDF\\
{\scriptsize\texttt{bacon2021}} & Difference-in-differences with variation in treatment timing \citep{bacon2021}.\par Journal of Econometrics 2021; accepted article copy\\
{\scriptsize\texttt{arellano1991}} & Some Tests of Specification for Panel Data: Monte Carlo Evidence and an Application to Employment Equations \citep{arellano1991}.\par Review of Economic Studies 58, 277–297; scanned published copy\\
{\scriptsize\texttt{bertrand2004}} & How Much Should We Trust Differences-in-Differences Estimates? \citep{bertrand2004}.\par QJE 2004; archived working manuscript\\
{\scriptsize\texttt{abadie2010}} & Synthetic Control Methods for Comparative Case Studies: Estimating the Effect of California’s Tobacco Control Program \citep{abadie2010}.\par JASA 2010; published copy\\
{\scriptsize\texttt{abadie2021}} & Using Synthetic Controls: Feasibility, Data Requirements, and Methodological Aspects \citep{abadie2021}.\par Journal of Economic Literature 2021; June 2020 manuscript\\
{\scriptsize\texttt{drechsler2017}} & The Deposits Channel of Monetary Policy \citep{drechsler2017}.\par QJE 2017; NBER 22152 copy revised January 2017\\
{\scriptsize\texttt{benetton2021}} & Leverage Regulation and Market Structure: A Structural Model of the UK Mortgage Market \citep{benetton2021}.\par Journal of Finance 2021; author accepted manuscript\\
{\scriptsize\texttt{basten2026}} & Cross-Selling in Bank-Household Relationships: Mechanisms and Implications for Pricing \citep{basten2026}.\par Review of Financial Studies 39(6), 2026 issue; earlier online publication; published copy\\
{\scriptsize\texttt{beraja2019}} & Regional Heterogeneity and the Refinancing Channel of Monetary Policy \citep{beraja2019}.\par QJE 134(1), 109–183; published copy\\
{\scriptsize\texttt{kaelbling1998}} & Planning and Acting in Partially Observable Stochastic Domains \citep{kaelbling1998}.\par Artificial Intelligence 101, 99–134; published copy\\
{\scriptsize\texttt{tennenholtz2020}} & Off-Policy Evaluation in Partially Observable Environments \citep{tennenholtz2020}.\par AAAI 2020 proceedings copy\\
{\scriptsize\texttt{nested2022}} & Automatic Debiased Machine Learning for Dynamic Treatment Effects and General Nested Functionals \citep{nested2022}.\par arXiv v4, September 2022; first author corrected from local filename\\
{\scriptsize\texttt{singh2022}} & A Finite Sample Theorem for Longitudinal Causal Inference with Machine Learning: Long Term, Dynamic, and Mediated Effects \citep{singh2022}.\par arXiv v2, 3 October 2022; initial submission 2021\\
{\scriptsize\texttt{bennett2021}} & Off-policy Evaluation in Infinite-horizon Reinforcement Learning with Latent Confounders \citep{bennett2021}.\par AISTATS 2021 proceedings copy\\
{\scriptsize\texttt{ghassami2022}} & Minimax Kernel Machine Learning for a Class of Doubly Robust Functionals with Application to Proximal Causal Inference \citep{ghassami2022}.\par AISTATS 2022 proceedings copy\\
{\scriptsize\texttt{xu2023}} & An Instrumental Variable Approach to Confounded Off-Policy Evaluation \citep{xu2023}.\par ICML 2023 proceedings copy; identity checked from local PDF\\
{\scriptsize\texttt{oberst2019}} & Counterfactual Off-Policy Evaluation with Gumbel-Max Structural Causal Models \citep{oberst2019}.\par ICML 2019; arXiv v3, June 2019\\
{\scriptsize\texttt{boeken2024}} & Dynamic Structural Causal Models \citep{boeken2024}.\par arXiv v2, 22 July 2024\\
{\scriptsize\texttt{bongers2022}} & Causal Modeling of Dynamical Systems \citep{bongers2022}.\par arXiv v4, 27 March 2022\\
{\scriptsize\texttt{tu2019}} & Causal Discovery in the Presence of Missing Data \citep{tu2019}.\par AISTATS 2019 proceedings copy\\
{\scriptsize\texttt{vargas2026}} & Cluster-Dags as Powerful Background Knowledge For Causal Discovery \citep{vargas2026}.\par arXiv v2, 19 January 2026; initial submission 2025\\
{\scriptsize\texttt{wager2025}} & Causal Inference: A Statistical Learning Approach \citep{wager2025}.\par Draft textbook, 26 November 2025\\
{\scriptsize\texttt{applied2024}} & Applied Causal Inference Powered by ML and AI \citep{applied2024}.\par Online book version 0.1.1, 4 March 2024\\
{\scriptsize\texttt{ovb2022}} & Long Story Short: Omitted Variable Bias in Causal Machine Learning \citep{ovb2022}.\par NBER Working Paper 30302, July 2022; not the later published version\\
{\scriptsize\texttt{agarwal2020}} & Mortgage Refinancing, Consumer Spending, and Competition: Evidence from the Home Affordable Refinancing Program \citep{agarwal2020}.\par NBER 21512, revised April 2020; local filename 2023 corrected to copied version\\
{\scriptsize\texttt{corrupt2022}} & Causal Inference with Corrupted Data: Measurement Error, Missing Values, Discretization, and Differential Privacy \citep{corrupt2022}.\par arXiv v5, 9 November 2022\\
{\scriptsize\texttt{swig2013}} & Single World Intervention Graphs: A Primer \citep{swig2013}.\par UAI 2013 workshop primer\\
{\scriptsize\texttt{runge2019}} & Detecting Causal Associations in Large Nonlinear Time Series Datasets \citep{runge2019}.\par Archived arXiv v2, June 2018; precursor of Science Advances 2019 article\\
{\scriptsize\texttt{runge2020}} & Discovering Contemporaneous and Lagged Causal Relations in Autocorrelated Nonlinear Time Series Datasets \citep{runge2020}.\par UAI 2020 method; arXiv v2, 5 January 2022\\
{\scriptsize\texttt{gerhardus2020}} & High-recall Causal Discovery for Autocorrelated Time Series with Latent Confounders \citep{gerhardus2020}.\par NeurIPS 2020; archived arXiv manuscript\\
{\scriptsize\texttt{pamfil2020}} & DYNOTEARS: Structure Learning from Time-Series Data \citep{pamfil2020}.\par AISTATS 2020 published proceedings\\
{\scriptsize\texttt{dagma2022}} & DAGMA: Learning DAGs via M-matrices and a Log-Determinant Acyclicity Characterization \citep{dagma2022}.\par NeurIPS 2022; archived arXiv v3, 16 January 2023\\
{\scriptsize\texttt{hyvarinen2010}} & Estimation of a Structural Vector Autoregression Model Using Non-Gaussianity \citep{hyvarinen2010}.\par JMLR 11, 1709–1731\\
{\scriptsize\texttt{shimizu2006}} & A Linear Non-Gaussian Acyclic Model for Causal Discovery \citep{shimizu2006}.\par JMLR 7, 2003–2030\\
{\scriptsize\texttt{kalisch2007}} & Estimating High-Dimensional Directed Acyclic Graphs with the PC-Algorithm \citep{kalisch2007}.\par JMLR 8, 613–636\\
{\scriptsize\texttt{causallearn2024}} & Causal-learn: Causal Discovery in Python \citep{causallearn2024}.\par JMLR 25, 2024\\
{\scriptsize\texttt{doubleml2022}} & DoubleML: An Object-Oriented Implementation of Double Machine Learning in Python \citep{doubleml2022}.\par JMLR 23, 2022\\
{\scriptsize\texttt{dowhy2020}} & DoWhy: An End-to-End Library for Causal Inference \citep{dowhy2020}.\par arXiv v1, 9 November 2020; current code assessed separately\\
{\scriptsize\texttt{borusyak2024}} & Revisiting Event Study Designs: Robust and Efficient Estimation \citep{borusyak2024}.\par Review of Economic Studies 2024; arXiv v5, January 2024\\
{\scriptsize\texttt{rambachan2023}} & A More Credible Approach to Parallel Trends \citep{rambachan2023}.\par Review of Economic Studies 90, 2555–2591; published PDF\\
{\scriptsize\texttt{synthdid2021}} & Synthetic Difference in Differences \citep{synthdid2021}.\par AER 2021; arXiv v4, July 2021\\
{\scriptsize\texttt{benmichael2021}} & The Augmented Synthetic Control Method \citep{benmichael2021}.\par JASA 2021; archived July 2020 manuscript\\
{\scriptsize\texttt{cluster2021}} & Multiway Cluster Robust Double/Debiased Machine Learning \citep{cluster2021}.\par Journal of Business and Economic Statistics online 2021; arXiv v3, March 2020\\
{\scriptsize\texttt{bica2020}} & Estimating Counterfactual Treatment Outcomes over Time Through Adversarially Balanced Representations \citep{bica2020}.\par ICLR 2020; arXiv v1\\
{\scriptsize\texttt{melnychuk2022}} & Causal Transformer for Estimating Counterfactual Outcomes \citep{melnychuk2022}.\par ICML 2022 proceedings copy\\
{\scriptsize\texttt{kallus2020}} & Double Reinforcement Learning for Efficient Off-Policy Evaluation in Markov Decision Processes \citep{kallus2020}.\par JMLR 21, 2020\\
{\scriptsize\texttt{proximal2020}} & An Introduction to Proximal Causal Learning \citep{proximal2020}.\par arXiv v1, 23 September 2020\\
{\scriptsize\texttt{miao2018}} & Identifying Causal Effects With Proxy Variables of an Unmeasured Confounder \citep{miao2018}.\par Biometrika 2018; arXiv v4, June 2018\\
{\scriptsize\texttt{bennettprox2023}} & Proximal Reinforcement Learning: Efficient Off-Policy Evaluation in Partially Observed Markov Decision Processes \citep{bennettprox2023}.\par arXiv v2, 22 March 2023\\
{\scriptsize\texttt{dynamicdml2021}} & Double/Debiased Machine Learning for Dynamic Treatment Effects via g-Estimation \citep{dynamicdml2021}.\par arXiv v5, 17 June 2021\\
{\scriptsize\texttt{kennedy2019}} & Nonparametric Causal Effects Based on Incremental Propensity Score Interventions \citep{kennedy2019}.\par JASA 2019; arXiv v3, June 2018\\
{\scriptsize\texttt{policy2021}} & Policy Learning with Observational Data \citep{policy2021}.\par Econometrica 2021; arXiv v6, September 2020\\
{\scriptsize\texttt{sequentialdr2018}} & Sequential Double Robustness in Right-Censored Longitudinal Models \citep{sequentialdr2018}.\par arXiv v2, 17 May 2018\\
{\scriptsize\texttt{lucid2026}} & LUCID: Learning Under Confounding for Inference and Discovery in Time Series \citep{lucid2026}.\par Preprint v1, 25 September 2026; preliminary research comparator\\
{\scriptsize\texttt{itsy2026}} & ITSY: Causal Discovery From Irregular Time-Series Data \citep{itsy2026}.\par Preprint v1, 20 September 2026; preliminary research comparator\\
{\scriptsize\texttt{causalts2026}} & Causal-TS: A Python Library for Causal Discovery in High-Dimensional and Nonstationary Time Series \citep{causalts2026}.\par Preprint v1, 27 July 2026; software claims not runtime-validated\\
{\scriptsize\texttt{panelpfn2026}} & Temporal Causal Prior-Data Fitted Networks for Panel Data with Learned Reliability Signals \citep{panelpfn2026}.\par Preprint v1, 18 June 2026; preliminary research comparator\\
{\scriptsize\texttt{paneldml2024}} & Double Machine Learning meets Panel Data: Promises, Pitfalls, and Potential Solutions \citep{paneldml2024}.\par arXiv v1, 2 September 2024\\
{\scriptsize\texttt{causalpfn2026}} & Foundation Models for Causal Inference via Prior-Data Fitted Networks \citep{causalpfn2026}.\par ICLR 2026 paper; arXiv v3, 24 February 2026\\
{\scriptsize\texttt{selective2025}} & Longitudinal Causal Inference with Selective Eligibility \citep{selective2025}.\par arXiv v2, March 2025\\
{\scriptsize\texttt{proxreg2024}} & Regression-Based Proximal Causal Inference \citep{proxreg2024}.\par arXiv v3, 5 June 2024\\
{\scriptsize\texttt{hidden2026}} & Causal Inference with a Hidden Treatment \citep{hidden2026}.\par arXiv v4, 22 June 2026\\
{\scriptsize\texttt{competing2026}} & Nonparametric Efficient Estimation of Dynamic Treatment Regimes with Competing Risks \citep{competing2026}.\par Preprint v1, 28 August 2026; preliminary research comparator\\
{\scriptsize\texttt{hmmident2016}} & Nonparametric Identification and Maximum Likelihood Estimation for Hidden Markov Models \citep{hmmident2016}.\par Biometrika 2016; arXiv v2, September 2015\\
{\scriptsize\texttt{lmtp2023}} & Non-parametric Causal Effects Based on Longitudinal Modified Treatment Policies \citep{lmtp2023}.\par JASA 2023 print publication; arXiv v4, July 2021\\
{\scriptsize\texttt{sun2021}} & Estimating Dynamic Treatment Effects in Event Studies with Heterogeneous Treatment Effects \citep{sun2021}.\par Journal of Econometrics 2021; archived arXiv manuscript\\
{\scriptsize\texttt{dechaisemartin2024}} & Difference-in-Differences Estimators of Intertemporal Treatment Effects \citep{dechaisemartin2024}.\par Review of Economics and Statistics, online 2024; archived arXiv v14, 12 May 2026\\
{\scriptsize\texttt{opehard2026}} & Exponential Hardness of Off-Policy Evaluation under History-Dependent Logging \citep{opehard2026}.\par Preprint v1, 16 September 2026; unreplicated theoretical comparator\\
{\scriptsize\texttt{sharpdynamic2026}} & Sharp Causal Bounds for Dynamic Treatment Regimes \citep{sharpdynamic2026}.\par Preprint v1, 6 September 2026; unreplicated theoretical comparator\\
{\scriptsize\texttt{proximalfm2026}} & ProximalFM: Amortized Proximal Causal Inference under Hidden Confounding \citep{proximalfm2026}.\par Preprint v1, 6 October 2026; very recent research comparator\\
{\scriptsize\texttt{sde2026}} & Causal Graphs, Markov Properties and Do-calculus for Stochastic Differential Equations \citep{sde2026}.\par Preprint v2, 22 July 2026\\
{\scriptsize\texttt{mian2006}} & Tracing the Impact of Bank Liquidity Shocks: Evidence from an Emerging Market \citep{mian2006}.\par NBER Working Paper 12612, October 2006; later AER publication reverses author order\\
{\scriptsize\texttt{gnet2020}} & G-Net: A Deep Learning Approach to G-computation for Counterfactual Outcome Prediction Under Dynamic Treatment Regimes \citep{gnet2020}.\par arXiv v1, March 2020\\
{\scriptsize\texttt{banksurvey2023}} & Causal Inference for Banking Finance and Insurance: A Survey \citep{banksurvey2023}.\par arXiv v1, 2023; 2025 ACM expanded publication discovered separately but publisher copy inaccessible\\
{\scriptsize\texttt{panelreview2024}} & Causal Models for Longitudinal and Panel Data: A Survey \citep{panelreview2024}.\par Archived arXiv v3, 25 June 2024\\
{\scriptsize\texttt{brodersen2015}} & Inferring Causal Impact Using Bayesian Structural Time-Series Models \citep{brodersen2015}.\par Annals of Applied Statistics 2015; arXiv v1\\
{\scriptsize\texttt{jiang2016}} & Doubly Robust Off-policy Value Evaluation for Reinforcement Learning \citep{jiang2016}.\par ICML 2016, PMLR 48; proceedings PDF\\
{\scriptsize\texttt{rabiner1989}} & A Tutorial on Hidden Markov Models and Selected Applications in Speech Recognition \citep{rabiner1989}.\par Proceedings of the IEEE 77(2), February 1989; author-hosted scan\\
{\scriptsize\texttt{jorda2024}} & Local Projections \citep{jorda2024}.\par NBER Working Paper 32822, August 2024; archived working-paper version\\
{\scriptsize\texttt{stock2018}} & Identification and Estimation of Dynamic Causal Effects in Macroeconomics Using External Instruments \citep{stock2018}.\par NBER Working Paper 24216, January 2018; archived working paper\\
{\scriptsize\texttt{pearl2009}} & Causal inference in statistics: An overview \citep{pearl2009}.\par Statistics Surveys 3 (2009), 96--146; published PDF\\
{\scriptsize\texttt{angrist1996}} & Identification of Causal Effects Using Instrumental Variables \citep{angrist1996}.\par Journal of the American Statistical Association 91(434), 444--455; published scan\\
{\scriptsize\texttt{grf2018}} & Generalized Random Forests \citep{grf2018}.\par Archived arXiv v4, 5 April 2018; submitted manuscript. Citation year denotes this inspected version, not the later journal publication.\\
{\scriptsize\texttt{rlearner2020}} & Quasi-Oracle Estimation of Heterogeneous Treatment Effects \citep{rlearner2020}.\par Archived arXiv v4, 6 August 2020; draft August 2020. Citation year denotes this inspected version, not the later journal publication.\\
{\scriptsize\texttt{drlearner2023}} & Towards Optimal Doubly Robust Estimation of Heterogeneous Causal Effects \citep{drlearner2023}.\par Archived arXiv v5, 21 August 2023; inspected manuscript version\\
{\scriptsize\texttt{snm2014}} & Structural Nested Models and G-estimation: The Partially Realized Promise \citep{snm2014}.\par Statistical Science 29(4), 707--731, 2014; arXiv v1 reprint 5 March 2015\\
{\scriptsize\texttt{id2006}} & Identification of Conditional Interventional Distributions \citep{id2006}.\par Proceedings of UAI 2006, 437--444; archived arXiv v1 submitted 27 June 2012, with proceedings identity checked in metadata\\
{\scriptsize\texttt{rohekar2023}} & From Temporal to Contemporaneous Iterative Causal Discovery in the Presence of Latent Confounders \citep{rohekar2023}.\par arXiv v1, 1 June 2023; PMLR 202, 2023 author version\\
{\scriptsize\texttt{slawski2025}} & Causal Secondary Analysis of Linked Data in the Presence of Mismatch Error \citep{slawski2025}.\par arXiv v1, 16 December 2025\\
{\scriptsize\texttt{zhangprox2024}} & On Identification of Dynamic Treatment Regimes with Proxies of Hidden Confounders \citep{zhangprox2024}.\par arXiv v1, 22 February 2024\\
{\scriptsize\texttt{wei2026}} & Off-Policy Evaluation for Missingness-Aware Policies in MDPs with Rewards Missing Not at Random \citep{wei2026}.\par arXiv v1, 18 June 2026\\
{\scriptsize\texttt{blundellbond2023}} & Initial conditions and Blundell–Bond estimators \citep{blundellbond2023}.\par Journal of Econometrics 234 (2023), 101--110; published PDF inspected\\
{\scriptsize\texttt{casini2025}} & Dynamic Local Average Treatment Effects in Time Series \citep{casini2025}.\par arXiv v1, 16 September 2025; internal date 17 September 2025\\
{\scriptsize\texttt{ran2026}} & Doubly Robust Functional Representation Learning for Longitudinal Causal Inference with Irregular Histories \citep{ran2026}.\par arXiv v1, 30 July 2026\\
{\scriptsize\texttt{kleibergen2021}} & Double robust inference for continuous updating GMM \citep{kleibergen2021}.\par arXiv v1, 18 May 2021; internal 19 May 2021. 2025 Quantitative Economics DOI 10.3982/QE2347 verified separately; journal revision not substituted\\
{\scriptsize\texttt{gerhardus2021}} & Characterization of causal ancestral graphs for time series with latent confounders \citep{gerhardus2021}.\par arXiv v1, 15 December 2021; 2024 Annals DOI 10.1214/23-AOS2325 verified separately; publisher endpoint did not return PDF\\
{\scriptsize\texttt{guha2022}} & Bayesian Causal Inference with Bipartite Record Linkage \citep{guha2022}.\par Bayesian Analysis 17(4), 1275–1299, publisher version 2022; compared with archived arXiv v2 September 2020\\
{\scriptsize\texttt{conformal2022}} & An Exact and Robust Conformal Inference Method for Counterfactual and Synthetic Controls \citep{conformal2022}.\par Preserved arXiv v10 labelled 20 May 2021, manuscript title page dated 25 May 2022; JASA DOI published online 1 June 2021; claims tied to preserved text\\
{\scriptsize\texttt{sw\_dowhy}} & py-why/dowhy: software documentation and source snapshot \citep{sw_dowhy}.\par Documentation snapshot 7 October 2026, qualified 8 October 2026; content hashes preserved; selected APIs commit-pinned; runtime untested\\
{\scriptsize\texttt{sw\_econml}} & py-why/EconML: software documentation and source snapshot \citep{sw_econml}.\par Documentation snapshot 7 October 2026, qualified 8 October 2026; content hashes preserved; selected APIs commit-pinned; runtime untested\\
{\scriptsize\texttt{sw\_doubleml}} & DoubleML/doubleml-for-py: software documentation and source snapshot \citep{sw_doubleml}.\par Documentation snapshot 7 October 2026, qualified 8 October 2026; content hashes preserved; selected APIs commit-pinned; runtime untested\\
{\scriptsize\texttt{sw\_causalml}} & uber/causalml: software documentation and source snapshot \citep{sw_causalml}.\par Documentation snapshot 7 October 2026, qualified 8 October 2026; content hashes preserved; selected APIs commit-pinned; runtime untested\\
{\scriptsize\texttt{sw\_grf}} & grf-labs/grf: software documentation and source snapshot \citep{sw_grf}.\par Documentation snapshot 7 October 2026, qualified 8 October 2026; content hashes preserved; selected APIs commit-pinned; runtime untested\\
{\scriptsize\texttt{sw\_tigramite}} & jakobrunge/tigramite: software documentation and source snapshot \citep{sw_tigramite}.\par Documentation snapshot 7 October 2026, qualified 8 October 2026; content hashes preserved; selected APIs commit-pinned; runtime untested\\
{\scriptsize\texttt{sw\_causallearn}} & py-why/causal-learn: software documentation and source snapshot \citep{sw_causallearn}.\par Documentation snapshot 7 October 2026, qualified 8 October 2026; content hashes preserved; selected APIs commit-pinned; runtime untested\\
{\scriptsize\texttt{sw\_causalnex}} & quantumblacklabs/causalnex: software documentation and source snapshot \citep{sw_causalnex}.\par Documentation snapshot 7 October 2026, qualified 8 October 2026; content hashes preserved; selected APIs commit-pinned; runtime untested\\
{\scriptsize\texttt{sw\_causica}} & microsoft/causica: software documentation and source snapshot \citep{sw_causica}.\par Documentation snapshot 7 October 2026, qualified 8 October 2026; content hashes preserved; selected APIs commit-pinned; runtime untested\\
{\scriptsize\texttt{sw\_gcastle}} & huawei-noah/trustworthyAI: software documentation and source snapshot \citep{sw_gcastle}.\par Documentation snapshot 7 October 2026, qualified 8 October 2026; content hashes preserved; selected APIs commit-pinned; runtime untested\\
{\scriptsize\texttt{sw\_dagma}} & kevinsbello/dagma: software documentation and source snapshot \citep{sw_dagma}.\par Documentation snapshot 7 October 2026, qualified 8 October 2026; content hashes preserved; selected APIs commit-pinned; runtime untested\\
{\scriptsize\texttt{sw\_notears}} & xunzheng/notears: software documentation and source snapshot \citep{sw_notears}.\par Documentation snapshot 7 October 2026, qualified 8 October 2026; content hashes preserved; selected APIs commit-pinned; runtime untested\\
{\scriptsize\texttt{sw\_tetrad}} & cmu-phil/tetrad: software documentation and source snapshot \citep{sw_tetrad}.\par Documentation snapshot 7 October 2026, qualified 8 October 2026; content hashes preserved; selected APIs commit-pinned; runtime untested\\
{\scriptsize\texttt{sw\_causalimpact}} & google/CausalImpact: software documentation and source snapshot \citep{sw_causalimpact}.\par Documentation snapshot 7 October 2026, qualified 8 October 2026; content hashes preserved; selected APIs commit-pinned; runtime untested\\
{\scriptsize\texttt{sw\_causalpy}} & pymc-labs/CausalPy: software documentation and source snapshot \citep{sw_causalpy}.\par Documentation snapshot 7 October 2026, qualified 8 October 2026; content hashes preserved; selected APIs commit-pinned; runtime untested\\
{\scriptsize\texttt{sw\_linearmodels}} & bashtage/linearmodels: software documentation and source snapshot \citep{sw_linearmodels}.\par Documentation snapshot 7 October 2026, qualified 8 October 2026; content hashes preserved; selected APIs commit-pinned; runtime untested\\
{\scriptsize\texttt{sw\_pyfixest}} & py-econometrics/pyfixest: software documentation and source snapshot \citep{sw_pyfixest}.\par Documentation snapshot 7 October 2026, qualified 8 October 2026; content hashes preserved; selected APIs commit-pinned; runtime untested\\
{\scriptsize\texttt{sw\_fixest}} & lrberge/fixest: software documentation and source snapshot \citep{sw_fixest}.\par Documentation snapshot 7 October 2026, qualified 8 October 2026; content hashes preserved; selected APIs commit-pinned; runtime untested\\
{\scriptsize\texttt{sw\_did}} & bcallaway11/did: software documentation and source snapshot \citep{sw_did}.\par Documentation snapshot 7 October 2026, qualified 8 October 2026; content hashes preserved; selected APIs commit-pinned; runtime untested\\
{\scriptsize\texttt{sw\_honestdid}} & asheshrambachan/HonestDiD: software documentation and source snapshot \citep{sw_honestdid}.\par Documentation snapshot 7 October 2026, qualified 8 October 2026; content hashes preserved; selected APIs commit-pinned; runtime untested\\
{\scriptsize\texttt{sw\_synthdid}} & synth-inference/synthdid: software documentation and source snapshot \citep{sw_synthdid}.\par Documentation snapshot 7 October 2026, qualified 8 October 2026; content hashes preserved; selected APIs commit-pinned; runtime untested\\
{\scriptsize\texttt{sw\_gsynth}} & xuyiqing/gsynth: software documentation and source snapshot \citep{sw_gsynth}.\par Documentation snapshot 7 October 2026, qualified 8 October 2026; content hashes preserved; selected APIs commit-pinned; runtime untested\\
{\scriptsize\texttt{sw\_d3rlpy}} & takuseno/d3rlpy: software documentation and source snapshot \citep{sw_d3rlpy}.\par Documentation snapshot 7 October 2026, qualified 8 October 2026; content hashes preserved; selected APIs commit-pinned; runtime untested\\
{\scriptsize\texttt{sw\_scope\_rl}} & hakuhodo-technologies/scope-rl: software documentation and source snapshot \citep{sw_scope_rl}.\par Documentation snapshot 7 October 2026, qualified 8 October 2026; content hashes preserved; selected APIs commit-pinned; runtime untested\\
{\scriptsize\texttt{sw\_obp}} & st-tech/zr-obp: software documentation and source snapshot \citep{sw_obp}.\par Documentation snapshot 7 October 2026, qualified 8 October 2026; content hashes preserved; selected APIs commit-pinned; runtime untested\\
{\scriptsize\texttt{sw\_pomdp\_py}} & h2r/pomdp-py: software documentation and source snapshot \citep{sw_pomdp_py}.\par Documentation snapshot 7 October 2026, qualified 8 October 2026; content hashes preserved; selected APIs commit-pinned; runtime untested\\
{\scriptsize\texttt{sw\_pomdps}} & JuliaPOMDP/POMDPs.jl: software documentation and source snapshot \citep{sw_pomdps}.\par Documentation snapshot 7 October 2026, qualified 8 October 2026; content hashes preserved; selected APIs commit-pinned; runtime untested\\
{\scriptsize\texttt{sw\_pymc}} & pymc-devs/pymc: software documentation and source snapshot \citep{sw_pymc}.\par Documentation snapshot 7 October 2026, qualified 8 October 2026; content hashes preserved; selected APIs commit-pinned; runtime untested\\
{\scriptsize\texttt{sw\_numpyro}} & pyro-ppl/numpyro: software documentation and source snapshot \citep{sw_numpyro}.\par Documentation snapshot 7 October 2026, qualified 8 October 2026; content hashes preserved; selected APIs commit-pinned; runtime untested\\
{\scriptsize\texttt{sw\_dynamax}} & probml/dynamax: software documentation and source snapshot \citep{sw_dynamax}.\par Documentation snapshot 7 October 2026, qualified 8 October 2026; content hashes preserved; selected APIs commit-pinned; runtime untested\\
{\scriptsize\texttt{sw\_duckdb}} & duckdb/duckdb: software documentation and source snapshot \citep{sw_duckdb}.\par Documentation snapshot 7 October 2026, qualified 8 October 2026; content hashes preserved; selected APIs commit-pinned; runtime untested\\
{\scriptsize\texttt{sw\_polars}} & pola-rs/polars: software documentation and source snapshot \citep{sw_polars}.\par Documentation snapshot 7 October 2026, qualified 8 October 2026; content hashes preserved; selected APIs commit-pinned; runtime untested\\
{\scriptsize\texttt{sw\_ray}} & ray-project/ray: software documentation and source snapshot \citep{sw_ray}.\par Documentation snapshot 7 October 2026, qualified 8 October 2026; content hashes preserved; selected APIs commit-pinned; runtime untested\\
{\scriptsize\texttt{sw\_datahub}} & datahub-project/datahub: software documentation and source snapshot \citep{sw_datahub}.\par Documentation snapshot 7 October 2026, qualified 8 October 2026; content hashes preserved; selected APIs commit-pinned; runtime untested\\
{\scriptsize\texttt{sw\_openmetadata}} & open-metadata/OpenMetadata: software documentation and source snapshot \citep{sw_openmetadata}.\par Documentation snapshot 7 October 2026, qualified 8 October 2026; content hashes preserved; selected APIs commit-pinned; runtime untested\\
{\scriptsize\texttt{sw\_openlineage}} & OpenLineage/OpenLineage: software documentation and source snapshot \citep{sw_openlineage}.\par Documentation snapshot 7 October 2026, qualified 8 October 2026; content hashes preserved; selected APIs commit-pinned; runtime untested\\
{\scriptsize\texttt{sw\_feast}} & feast-dev/feast: software documentation and source snapshot \citep{sw_feast}.\par Documentation snapshot 7 October 2026, qualified 8 October 2026; content hashes preserved; selected APIs commit-pinned; runtime untested\\
{\scriptsize\texttt{sw\_splink}} & moj-analytical-services/splink: software documentation and source snapshot \citep{sw_splink}.\par Documentation snapshot 7 October 2026, qualified 8 October 2026; content hashes preserved; selected APIs commit-pinned; runtime untested\\
{\scriptsize\texttt{sw\_pandera}} & unionai-oss/pandera: software documentation and source snapshot \citep{sw_pandera}.\par Documentation snapshot 7 October 2026, qualified 8 October 2026; content hashes preserved; selected APIs commit-pinned; runtime untested\\
{\scriptsize\texttt{sw\_lmtp}} & nt-williams/lmtp: software documentation and source snapshot \citep{sw_lmtp}.\par Documentation snapshot 7 October 2026, qualified 8 October 2026; content hashes preserved; selected APIs commit-pinned; runtime untested\\
{\scriptsize\texttt{sw\_ltmle}} & joshuaschwab/ltmle: software documentation and source snapshot \citep{sw_ltmle}.\par Documentation snapshot 7 October 2026, qualified 8 October 2026; content hashes preserved; selected APIs commit-pinned; runtime untested\\
{\scriptsize\texttt{sw\_gformula}} & CausalInference/gfoRmula: software documentation and source snapshot \citep{sw_gformula}.\par Documentation snapshot 7 October 2026, qualified 8 October 2026; content hashes preserved; selected APIs commit-pinned; runtime untested\\
{\scriptsize\texttt{sw\_zepid}} & pzivich/zEpid: software documentation and source snapshot \citep{sw_zepid}.\par Documentation snapshot 7 October 2026, qualified 8 October 2026; content hashes preserved; selected APIs commit-pinned; runtime untested\\
{\scriptsize\texttt{sw\_pydynpd}} & dazhwu/pydynpd: software documentation and source snapshot \citep{sw_pydynpd}.\par Documentation snapshot 7 October 2026, qualified 8 October 2026; content hashes preserved; selected APIs commit-pinned; runtime untested\\
{\scriptsize\texttt{fed2026}} & Revised Guidance on Model Risk Management \citep{fed2026}.\par 17 April 2026; SR26-2 supersedes SR11-7 and SR21-8\\
{\scriptsize\texttt{pra2026}} & Model Risk Management Principles for Banks \citep{pra2026}.\par SS1/23, April 2026 revision; original implementation date retained in text\\
{\scriptsize\texttt{bcbs239}} & Principles for Effective Risk Data Aggregation and Risk Reporting \citep{bcbs239}.\par BCBS239, January 2013\\
{\scriptsize\texttt{localrepos}} & Local banking and scientific repositories: inspected documentation snapshots \citep{localrepos}.\par Documentation snapshots taken 7 October 2026; some documents describe earlier milestones\\
{\scriptsize\texttt{runtime\_duckdb}} & duckdb 1.5.5: installed API documentation and operation receipt \citep{runtime_duckdb}.\par Installed version 1.5.5 observed 8 October 2026; release date at the literature cutoff not inferred\\
{\scriptsize\texttt{runtime\_polars}} & polars 1.44.1: installed API documentation and operation receipt \citep{runtime_polars}.\par Installed version 1.44.1 observed 8 October 2026; release date at the literature cutoff not inferred\\
{\scriptsize\texttt{runtime\_scikit\_learn}} & scikit-learn 1.9.0: installed API documentation and operation receipt \citep{runtime_scikit_learn}.\par Installed version 1.9.0 observed 8 October 2026; release date at the literature cutoff not inferred\\
{\scriptsize\texttt{runtime\_scipy}} & scipy 1.18.1: installed API documentation and operation receipt \citep{runtime_scipy}.\par Installed version 1.18.1 observed 8 October 2026; release date at the literature cutoff not inferred\\
{\scriptsize\texttt{runtime\_numpy}} & numpy 2.5.2: installed API documentation and operation receipt \citep{runtime_numpy}.\par Installed version 2.5.2 observed 8 October 2026; release date at the literature cutoff not inferred\\
{\scriptsize\texttt{runtime\_doubleml\_scope}} & doubleml 0.11.4: inspected installed causal operations \citep{runtime_doubleml_scope}.\par Installed 0.11.4; release 2026-08-10 from preserved PyPI metadata; observed 8 October 2026\\
{\scriptsize\texttt{runtime\_tigramite\_scope}} & tigramite 5.2.10.1: inspected installed causal operations \citep{runtime_tigramite_scope}.\par Installed 5.2.10.1; release 2026-01-14 from preserved PyPI metadata; observed 8 October 2026\\
{\scriptsize\texttt{public401k}} & Archived SIPP 1991 401(k) application extract from DMLonGitHub \citep{public401k}.\par Unversioned public branch retrieved 8 October 2026, identified by exact byte hash; 1991 SIPP application extract\\
{\scriptsize\texttt{runtime\_statsmodels\_dynamic}} & statsmodels 0.15.0 installed OLS and Bartlett HAC operations \citep{runtime_statsmodels_dynamic}.\par Installed 0.15.0; PyPI upload 27 August 2026, before cutoff; installed source inspected 9 October 2026\\
{\scriptsize\texttt{public\_macrodata\_dynamic}} & United States macroeconomic data: 1959Q1–2009Q3 archived extract \citep{public_macrodata_dynamic}.\par Installed macrodata.csv in statsmodels 0.15.0; 203 rows, source access 15 December 2009\\
{\scriptsize\texttt{fed\_h8\_dynamic}} & H.8 Assets and Liabilities of Commercial Banks in the United States: About the release \citep{fed_h8_dynamic}.\par Public page retrieved 9 October 2026; snapshot observation after 7 October cutoff, used for data-description scope not a new scholarly result\\
{\scriptsize\texttt{runtime\_networkx\_dynamic}} & NetworkX 3.6.1 and Tigramite 5.2.10.1 finite-DAG d-separation operations \citep{runtime_networkx_dynamic}.\par Installed NetworkX 3.6.1 and Tigramite 5.2.10.1, observed 9 October 2026; exact operation snapshots\\
\bottomrule
\end{longtable}

\normalsize

\backmatter
\bibliographystyle{plainnat}
\bibliography{causal_banking_refs}
\end{document}
